% ============================================================
%  H. Høffding, Totalitet som Kategori:
%  en erkendelsesteoretisk Undersøgelse
%  D. Kgl. Danske Vidensk. Selsk. Skrifter, 7. Række,
%  historisk og filosofisk Afd. III, 2.
%  København: Andr. Fred. Høst & Søn / Bianco Lunos Bogtrykkeri, 1917.
%
%  DIPLOMATIC TRANSCRIPTION (Danish)
%  Status: ALL PAGES TRANSCRIBED 2026-09-30 (78/78: body 3–78,
%  Kategoritavle 79, Indhold and Trykfejl 80).
%
%  SECOND READING (2026-09-30): every page 3–80 read word for word at the
%  image by a reader who had not transcribed it, text layer withheld (12 pp.
%  as a random sample, the other 66 pp. by seven readers; in the last three
%  readers' pages, 6 planted control errors were all caught). 5 word errors
%  found and corrected: printer's wrong sorts silently normalised on pp. 7
%  („Dømmen“), 13 („Sammenbængen“), 51 („Forboldet“), 75 („pet“), now
%  restored as printed with comments; and „KANTs“ for the printed „KANT's“
%  (p. 75 note 3). Two e's without a crossbar (pp. 45 n., 61) are logged as
%  probable broken sorts. Pp. 14–23, 25–32, 34–43 (27 pp.) were read without
%  planted controls and returned no errors; that reading is recorded as done
%  but uncontrolled (decision 2026-09-30).
%
%  Word-accuracy: every batch diffed against the scan's embedded ABBYY layer
%  before splicing (ocrdiff.py). 10 batches, 370 candidates, 0 transcription
%  errors found by the diff, 0 unresolved (all 370 were the layer's faults,
%  chiefly t read as l). Separately, the transcribers caught 4 slips of their
%  own at the image that the diff had not flagged (pp. 11, 28 ×2, 58), all
%  corrected. See TRANSCRIPTION-PLAYBOOK.md §3 step 4.
%  Structural control: 103 footnotes, matching the treatise's note count
%  chapter by chapter (12 + 9 + 24 + 39 + 13 + 6).
%
%  CONVENTIONS. Diplomatic: 1917 orthography as printed; low-high „…“ as
%  printed; printer's errors kept and logged "% PRINTED AS IS" at the site;
%  small caps (personal names) \textsc; italic titles \textit; letterspacing
%  \emph; footnotes per printed page as in the print; signature marks and
%  stray marks in the scanned copy omitted and logged in % comments.
% ------------------------------------------------------------
%  THE WITNESS
% ------------------------------------------------------------
%  The Royal Danish Academy's own digitisation (ABBYY FineReader Server, 2019),
%  80 pp., 400 ppi colour, antiqua text layer with italics tagged.
%    Scan: ~/bibliotek/Høffding, Harald/totalitet.pdf
%    sha256 40900a16c165d0b7674f315d06942116b07feb5d5b7899851b25b1fa17b562dd
%    Online: http://publ.royalacademy.dk/books/419/2912
%
%  PAGE MAP: printed page = PDF page (offset 0), pp. 3–80. Verified
%  2026-09-30 from the running heads of every page 4–78, which print both the
%  treatise's own numeral and the volume's continuous numeral (own + 256; the
%  treatise is vol. pp. 259–334); p. 3 opens ch. I and carries no head.
%  PDF 1 = title page, PDF 2 = blank verso. The markers \opage{N} use the
%  treatise's own numbering, by which Høffding himself cites it
%  (Erkendelsesteori og Livsopfattelse, 1925).
%
%  THE PRINT'S OWN ERRATA (Trykfejl, p. 80) are NOT applied in this
%  transcription: the text reads as printed, with each erratum noted in a
%  % comment at its site. (The English translation applies them.)
% ------------------------------------------------------------

\documentclass[12pt]{article}
\usepackage[utf8]{inputenc}
\usepackage[T1]{fontenc}
\usepackage[danish]{babel}
\usepackage[protrusion=true]{microtype}
\usepackage[margin=1.4in]{geometry}
\usepackage{setspace}
\usepackage{amsmath}
\usepackage{libertinus}
\usepackage{libertinust1math}
\usepackage{textalpha}
\usepackage{fancyhdr}
\usepackage[colorlinks=true,linkcolor=black,urlcolor=black,
  pdftitle={Totalitet som Kategori},
  pdfauthor={Harald Høffding}]{hyperref}
\setlength{\headheight}{15pt}
\pagestyle{fancy}
\fancyhf{}
\fancyhead[LE,RO]{\thepage}
\fancyhead[C]{\textit{Totalitet som Kategori}}
\setlength{\parindent}{1.5em}
\setlength{\parskip}{0pt}
\setstretch{1.3}

\title{\textbf{Totalitet som Kategori}\\[0.5em]
  \large en erkendelsesteoretisk Undersøgelse}
\author{Harald Høffding}
\date{D.\ Kgl.\ Danske Vidensk.\ Selsk.\ Skrifter, 7.\ Række,
  historisk og filosofisk Afd.\ III.\ 2.\\
  København: Hovedkommissionær Andr.\ Fred.\ Høst \& Søn,
  Kgl.\ Hof-Boghandel.\\ Bianco Lunos Bogtrykkeri, 1917}

\usepackage{marginnote}
\usepackage{graphicx}
% The id-est sign „ɔ:“ (a turned c), p. 39 and later.
\DeclareUnicodeCharacter{0254}{\reflectbox{c}}
\newcommand{\opage}[1]{\marginnote{\footnotesize\textit{[#1]}}}
\newcommand{\apage}[1]{\marginnote{\footnotesize\textit{[#1]}}}
\begin{document}
\maketitle
\tableofcontents
\bigskip

% --- p. 3 ---
% CHAPTER HEAD as printed on p. 3, two centred lines ("I." / title with
% final stop) over a short rule; p. 3 has no running head.
% DIVISION NUMERALS (1--29, continuous through the book) are run-in at the
% head of an indented paragraph, roman, not bold, not spaced: set plain
% "1. ". The print opens a little extra space before each division after
% the first; set \medskip on its own line before it.
\section*{I.\ Kategorilærens Stilling i Filosofien.}
\addcontentsline{toc}{section}{I.\ Kategorilærens Stilling i Filosofien}

\opage{3}1. Kategorier er Grundbegreber, som indeholdes i de Forudsætninger, paa
hvilke al videnskabelig Forsken hviler, og som kan paavises ved Analyse af disse
Forudsætninger. Men til fuld Belysning af deres Natur er det ikke nok at gaa ud
fra den til en vis Tid bestaaende Videnskab. Det maa ogsaa paavises, hvorledes
de har arbejdet sig frem gennem Videnskabens Historie, --- hvorledes de opstaar
som Former for Tankearbejde, --- og hvorledes de kan forsvinde igen, naar de ikke
længere er Udtryk for den arbejdende Tanke. Og fra en saadan historisk
Betragtning maa der igen gaas tilbage til en psykologisk Undersøgelse af
Tænkningens Udvikling i det sædvanlige menneskelige Bevidsthedsliv, af hvilket
den videnskabelige
% DEFECT (p. 3): a stray low dot stands after "en" ("en. ved"); a loose
% sort, not punctuation the sense wants. Not transcribed.
Tænkning er en ved bestemte Opgaver betinget særegen Udformning. De
Grundbegreber, hvormed Videnskaben arbejder, maa være psykologisk mulige, og
deres første Udspring maa kunne findes i det uvilkaarlige Tankeliv. Forskellen
mellem den psykologiske og den erkendelsesteoretiske Betragtning af Kategorierne
bestaar deri, at hin finder Mulighed saavel for uvidenskabelig som for
videnskabelig Tænkning, medens denne paaviser, at der under visse givne Forhold
kun gives en eneste snever Vej til Forstaaelse. Det gælder om vore Tanker, at
mange er kaldede, faa udvalgte, --- ja, vi maa som Psykologer føje til, at mange
er ukaldede.

I erkendelsesteoretiske Arbejder vil man finde forskellig Vægt lagt paa
Grundbegreber og paa Grundsætninger. Nogle Erkendelsesteoretikere holder sig ene
til Grundsætningerne, de Sætninger, videnskabelig Begrundelse hviler paa. Og da
Sætninger, der er Forudsætninger for al Begrundelse, naturligvis ikke selv kan
begrundes, fremstilles Grundsætningerne ofte som beroende paa et Valg, et
Spring, en Akt, der i og for sig ligesaa godt kunde være falden ud i helt anden
Retning. Videnskabens første Principer er efter denne Opfattelse ikke selv
Genstand for Videnskab. Dog indrømmes det fordetmeste af dem, der hylder denne
Opfattelse, at Valget er bundet indenfor en vis Kreds af Muligheder, og denne
Kreds af Muligheder er given med den menneskelige Tænknings Natur, saaledes som
denne nu engang er beskaffen. Denne faktiske Beskaffenhed er det Psykologiens
Opgave at undersøge, og der stiller sig da det Spørgsmaal, hvorledes de
forskellige Synspunkter, fra hvilke Erkendelsesteori og Psykologi vurderer
menneskelig Erkendelse, for%
% SIGNATURE MARK "33*" at the foot of p. 3 omitted as printer's apparatus.
% --- p. 4 ---
\opage{4}holder sig til hinanden. Men allerede før vi tager dette Spørgsmaal op,
vil det være klart, at der er et saadant Vexelforhold mellem Grundbegreber og
Grundsætninger, ligesom i det hele mellem Begreber og Domme, at Undersøgelsen
over de enkelte Begreber, som Grundsætningerne indeholder, vil kaste Lys over
selve Grundsætningerne. Fremskridende Klarhed bliver overhovedet kun mulig ved,
at vi med udtrykkelig Bevidsthed undersøger og bestemmer de Begreber, der
indeholdes i vore Domme, ligesom vi paa den anden Side undersøge og bestemme
Arten af den Forbindelse, der i Dommene hævdes mellem Begreberne. De to
Undersøgelser supplerer hinanden. Fremskridende Klarhed i Begrebsbestemmelse er
Betingelse for fremskridende Klarhed i Forstaaelse af de Domme, i hvilke
Begreberne indgaar, og omvendt. Om man lægger størst Vægt paa den ene eller den
anden Undersøgelse, vil afhænge af den specielle Opgave, man stiller sig. Men
til et fuldstændigt Opgør kan intet af de to Synspunkter undværes. Dertil kommer
endnu, at Tanketypen er den samme i Begreber som i Domme. Dommen er en
Forbindelse af Begreber, men Begrebet er en Forbindelse af Kendemærker. Det er
samme Grundegenskab ved vor Tænkning, der udtrykker sig begge Steder, i
Begrebet mere som Struktur, i Dommen mere som Funktion.

Medens man ved blot at holde sig til Videnskabens første Forudsætninger kommer
til en stejl Postulatteori, der stiller alt videnskabeligt Tankearbejde i et
fjernt Forhold eller endog i skarp Modsætning til alt andet aandeligt Arbejde,
vil man ved tillige at undersøge Grundbegreberne kunne føres til Forstaaelse af
det videnskabelige Tankearbejdes Forhold til Aandslivet i det hele. Og da det nu
netop er en af Filosofiens Opgaver at undersøge dette Forhold, vil det være af
Betydning at forsøge en Bestemmelse af Kategorilærens Stilling indenfor
Filosofien.

\medskip
2. Naar begynder Videnskaben? Og hvorledes forholder den videnskabelige
% DEFECT (p. 4): a stray raised tick stands over the comma after
% "Tankeliv"; not transcribed, the comma kept.
Bevidsthed sig til det Tankeliv, der udfolder sig forud for og jævnsides med det
videnskabelige, i større eller mindre Uafhængighed af dette?

Til Belysning af disse Spørgsmaal maa det først fastholdes, at der gør sig en
fælles Type gældende i alt menneskeligt Tankeliv, baade det førvidenskabelige,
det videnskabelige og det uvidenskabelige.

Selve det praktiske Liv, dets Fornødenheder og dets Opgaver, fører til at
tænke. For at et Menneske kan opnaa de Formaal, det stræber efter, er der visse
Midler, der maa skaffes tilveje. Her lærer Mennesket faktisk og praktisk at
kende, hvad Nødvendighed er, d.~v.~s. det lærer det ubønhørlige Forhold mellem
Grund og Følge at kende. Saalænge det fastholder Formaalet, saalænge maa det
ogsaa indrømme Nødvendigheden af Midler, tilsidst visse bestemte Midler, og det
lærer omvendt at se, at Formaalets Opnaaelse er en Følge af visse Midlers
Anvendelse. Tillige lærer det, at naar de samme Fornødenheder og de samme ydre
Forhold gentager sig, maa der gribes til de samme Midler. En Centralaustralier
samler f.~Ex. Planter, tærsker dem med en Stok, sigter Frøet, rister det, knuser
det paa Stene og spiser det: en Række af Handlinger, der, hvergang han er
sulten, maa foretages
% --- p. 5 ---
\setcounter{footnote}{0}%
\opage{5}paa samme Maade og i samme Rækkefølge. I Grunden kan allerede i denne
Rækkedannelse Videnskabens første Forudsætninger siges at være tilstede.

Man har i den nyeste Tid lagt særlig Vægt paa det sociale Livs Indflydelse paa
Tænkningens første Udvikling. I Stammens Inddeling i forskellige Grupper eller
Klasser forelaa der et Forbillede for al Klassifikation, og dermed for al Logik.
Man taler jo endnu stedse om Slægter og Familier ved Inddeling af
Naturgenstande, til Trods for, at ivrige Arvelighedsteoretikere her finder
falske Analogier. Tidsforestillingen og Tidsinddelingen fik sin Udvikling gennem
religiøse Festers og Ceremoniers rytmiske Forløb og regelmæssige Genkomst.
Endnu angiver jo vor Almanak Kirkeaarets Skema. Og den religiøse Kultus krævede,
at liturgiske Handlinger skulde foretages paa en aldeles bestemt, i alle
enkelte Tilfælde ens Maade, hvis Resultatet skulde naas. Her indskærpedes
saaledes den Sætning, at det Lige maa frembringes af det Lige. Rummets
Inddeling fik sin Klarhed formedelst de forskellige Retningers og de
forskellige Steders religiøse Betydning\footnote{\textsc{Émile Durkheim} har,
især i sit Værk \textit{Les formes élémentaires de la vie religieuse} (1912),
gjort denne Synsmaade gældende og skærpet den til den Paastand, at alle
Kategorier er af religiøs-social Oprindelse. Smlgn. min Afhandling
\textit{Sociologi og Religionsfilosofi} (Vor Tid, 1914---15), (paa Fransk i
Revue de la Métaphysique et de Morale. Novbr. 1914).}.

At de sociale Forhold, under hvilke Mennesker udvikler sig, øver en meget stor
Indflydelse paa deres Tankegang og deres Tankeformer, er utvivlsomt, selv om
Alt ikke kan forklares ad denne Vej. Den sociale Indflydelse brydes stedse med
den praktiske Erfaring, de Enkelte kan gøre paa egen Haand, Ansigt til Ansigt
med de Livsforhold, der ogsaa faar Indflydelse paa
% PRINTED AS IS (p. 5): "Samfundete" for Samfundets.
Samfundete Færd i det hele.

Hvad enten man nu vil lægge størst Vægt paa de enkelte Individers praktiske
Erfaringer eller paa sociale Institutioners Indflydelse, en Kendsgerning er
det, at der til hver given Tid og hos hvert enkelt Folk fremtræder en vis,
uvilkaarlig og tildels ubevidst udviklet Verdensopfattelse, der er bygget efter
samme Hovedformer som de, der fremtræder paa mere stringent og speciel Maade i
videnskabelig Tænkning. Den vidner om Trang og Evne til at ordne og tyde
Iagttagelser og Traditioner og derved bringe dem i en vis Sammenhæng. Det
Tankearbejde, her gøres, foregaar uden klar Bevidsthed. Der henvendes ikke
nogen Opmærksomhed paa de enkelte Tankeovergange eller paa Berettigelsen af at
forbinde de enkelte Elementer, som Iagttagelse og Tradition frembyder, paa en
vis bestemt Maade. For at bruge den Terminologi, jeg andensteds har anvendt:
det er ikke analytisk og syntetisk Intuition, men praktisk og konkret
Intuition, der her foreligger\footnote{Se \textit{Om Begrebet Intuition} (Vid.
Selsk. Oversigt 1914) p. 77.}. Den saakaldte sunde Menneskeforstand opererer
gennemgaaende med disse sidstnævnte Arter af Intuition. Derigennem kan den, for
at bruge \textsc{Meyerson}'s Udtryk\footnote{\textsc{Meyerson} har paa en
fortræffelig Maade drøftet Begrebet sens commun i sin Bog \textit{Identité et
% PRINTED AS IS (p. 5, note 3): "Réalté" for Réalité (verified at 500 dpi).
Réalté} (1908). Chap. XI.}, give „et første, meget plumpt Udkast til et
videnskabeligt og metafysisk System“. --- I Sproget, i de Ord, Sproget danner,
og de Ordforbindelser, det former, indeholdes allerede et saadant System, der
stadig gør sin Indflydelse gældende paa de videnskabelige Tankeformer.

% --- p. 6 ---
\opage{6}Videnskabens Udspring og Udvikling, saaledes som Videnskabens Historie
viser os den, faar egentlig Spørgsmaalet om dens første Begyndelse til helt at
falde bort, da det viser sig, at dens første Spire er tilstede, saasnart
Mennesker begynde at tænke. Der gør sig paa alle Trin og under alle Forhold en
vis Modsætning gældende mellem ét Givet, Foreliggende og dets Tydning, mellem
det enkelte Emne og den Sammenhæng, i hvilken det indføjes. Fremskridtet
bestaar i, at der stedse bliver et nøjere og strengere Forhold mellem
Iagttagelse og Tydning. Mere og mere bliver man sig bevidst, at al Tydning
tilsidst selv maa hentes fra Iagttagelser, og at Indføjelsen af det enkelte
Emne som Led i en hel Sammenhæng maa ske efter nøjagtig Prøvelse baade af den
enkelte Iagttagelse selv og af de andre Iagttagelser, med hvilke den forbindes.
Som det i det Følgende vil blive vist, er det her især Forskellen mellem
Identitet og Analogi, samt Grænsen for Analogiens Berettigelse, det gælder at
have Øje for.

\medskip
3. Den fælles Type for alt Tankeliv, det videnskabelige saavel som det
førvidenskabelige og det uvidenskabelige, er det Psykologiens Opgave at finde.

Denne Type kan i sin Almindelighed betegnes som en Helhedsdannelse. Allerede de
simpleste Sansefornemmelser viser sig at staa i Forbindelse med hverandre, saa
at den enkelte ikke kan forstaas, hverken hvad Styrkegrad eller Kvalitet angaar,
uden ved at ses som Led i en hel Række eller Gruppe af Fornemmelser. Allerede
Tidsforholdet mellem de forskellige Fornemmelser er afgørende for den Grad af
Selvstændighed, med hvilken de optræde. Dertil kommer, at Reproduktioner af
tidligere Fornemmelser mere eller mindre nøje forener sig med de nye opdukkende
Fornemmelser, ofte lige fra deres Fødelsesstund. Derved bliver Genkendelse,
partiel Perception og Sanseillusion mulige. Hvor de reproducerede Elementer
faar afgjort Overhaand, optræder Erindring og Fantasi som en Helhedsdannelse af
højere Orden. Gennem Association, Forskydning og Sammenligning kan Sanse-,
Erindrings- og Fantasianskuelser undergaa Ændringer, Artikulationer, uden at
Karakteren af Anskuelighed forsvinder. Medvirkende ved alle disse
Helhedsdannelser er fra først til sidst praktiske Motiver, Trang og Drift, Lyst
og Ulyst.

Karakteristisk for de her omtalte Helhedsdannelser (der fører til, hvad jeg har
kaldt konkret og praktisk Intuition) er den uvilkaarlige Maade, paa hvilken de
bliver til. Deraf følger den Selvfølgelighed og Oprindelighed, de synes at
besidde. De Elementer, af hvilke de er opbyggede, er ikke selv tilgængelige for
os paa samme Maade som de Elementer, af hvilke vi med fuld Bevidsthed danner
Begreber, Domme og Slutninger, idet Begrebets Kendemærker, Dommens Led og
Slutningens Præmisser da er givne forud for de nye Helhedsdannelser, i hvilke
de indgaar. Derved træder, hvad jeg har kaldt analytisk og syntetisk Intuition,
i Modsætning til konkret og praktisk Intuition. Nu er det faktisk ikke alle
Begreber, Domme og Slutninger, vi danner med fuld Bevidsthed; de kan ogsaa
opstaa ved uvilkaarlig Helhedsdannelse. Men skal Begreber, Domme og Slutninger
have videnskabelig Betydning, maa de kunne dannes paany, saaledes at der er
fuld og klar Bevidsthed tilstede
% --- p. 7 ---
% p. 7: stray marks at the head (a dot beside the folio; the layer reads
% "*", "►", "*") and the scan's margin arrows are not text of the work.
\setcounter{footnote}{0}%
\opage{7}ved hver enkelt Overgang fra Tanke til Tanke (analytisk Intuition)
saavel som ved hver enkelt Forbindelse af Tanke og Tanke (syntetisk Intuition).
Selv ved uvilkaarlig Dannelse af Begreber, Domme og Slutninger fremtræder dog
Eftertanken som noget fra den umiddelbare Anskuen (i Sansning, Erindring og
Fantasi) forskelligt, idet der finder en mere eller mindre tydelig Overgang og
Sammenligning Sted mellem de Elementer, der bringes i Forbindelse, medens de
Elementer, af hvilke umiddelbare Anskuelser dannes, kun indirekte og tildels ad
Analogiens Vej kan efterspores.

% PRINTED AS IS (p. 7): „Dømmen“ for „Dommen“ (wrong sort; second reading
% 2026-09-30, clear at 500 dpi).
Da Begrebsdannelse sker ved Dømmen, og da Slutningen er en Doms Dannelse ud fra
andre Domme, kan vi i denne Sammenhæng holde os til Dommene.

Dommen dannes fra først af paa Grundlag af umiddelbar Anskuen. Den angiver et
bestemt Forhold mellem Elementer i Anskuelsens Indhold, Elementer, der i selve
Anskuelsen er i kontinuerlig Forbindelse eller maaske endog sammensmeltede med
hverandre. Dommen „Bordet er firkantet“ dannes paa Grundlag af en Anskuen, i
hvilken Egenskaben „firkantet“ allerede var indeholdt („laa“), men uden at
blive draget særlig frem og uden at stilles i Forhold til Bordet som Helhed,
d.~v.~s. som Indbegreb af alle Egenskaber.

Men hvorfor dannes overhovedet Domme, --- hvorfor bliver vi ikke staaende ved
Anskuen? Det er sikkert ikke, som nogle Forfattere synes at opfatte Sagen, en
Art Syndefald, der her har fundet Sted. Det er Iagttagelsens Rigdom og
Mangfoldighed og Erindringens Forsøg paa at omspænde denne Rigdom, der gør, at
en anden Art Helhedsdannelse end den umiddelbare Anskuen bliver nødvendig. Den
primære Helhedsform sprænges, og der søges en ny.

Det kan være vanskeligt at drage en skarp Grænse mellem Anskuelse og Dom. Som
allerede antydet, kan Anskuelsen ændres, artikuleres. Det sker ved, at enkelte
Dele skyder sig frem i Forgrunden, medens andre træder tilbage. Ved Bordet kan
f.~Ex. Farven særlig gøre sig gældende fremfor Figuren, ved selve Figuren
Konturen fremfor Fladens Form\footnote{Smlgn. \textsc{Edgar Rubin}:
\textit{Synsoplevede Figurer.}}. Saadanne Forskydninger foregaar uafbrudt ved
vore Sansninger, Erindringer og Fantasier. Men en Dom dannes først, naar vi
bliver os disse Overgange bevidst istedetfor at gaa helt op i de Anskuelser,
der er tilstede i de enkelte Øjeblikke, og som afløser hverandre. Nu kan der vel
ogsaa, hvad denne Bevidsthed angaar, være mange Overgange og Grader; men derfor
staar Endepunkterne alligevel i karakteristisk Modsætning til hinanden, og det
vilde være uberettiget at overføre Alt, hvad der gælder for Dommen, paa
Anskuelsen. Der er kun Analogi imellem dem, en Analogi, der grunder sig paa, at
de er to forskellige Maader, paa hvilke den psykiske Energi, d.~v.~s. Trangen
og Evnen til Helhedsdannelse, kan ytre sig.

Grundlæggerne af Kategorilæren, \textsc{Aristoteles} og \textsc{Kant}, har
rigtig set, at Erkendelsens videnskabelige Grundformer maatte kunne findes ved
Analyse af Dommene, som fældes. Hvad de ikke klart nok har set, er, at den
psykiske Energi, der virker i Domsdannelsen, ogsaa, under andre Betingelser,
virker i andre Hel%
% --- p. 8 ---
\setcounter{footnote}{0}%
\opage{8}hedsdannelser. Kun gennem den Kontinuitet og Analogi, der gør sig
gældende mellem den primære Helhedsdannelse og Dommen, forstaas det, hvorledes
den videnskabelige Bevidsthed staar i nøje Sammenhæng med Bevidsthedslivet i
det hele og er en speciellere Udformning af dette. Men det er ogsaa netop kun
et Analogiforhold, der er tilstede, som mellem to Exempler paa samme Begreb. Da
Domme kan udledes af, bygges paa Anskuelser, kan man jo gerne sige, at de
„indeholdes“ i dem eller „ligger“ i dem. Men det er billedlig Tale. Fordi en
Dom er logisk mulig, er det nemlig ikke sagt, at den er psykologisk mulig, og i
hvert Tilfælde er det en Omsætning til en anden Form, der sker ved Overgangen
fra Anskuelse til Dom, en Overgang, der forudsætter visse bestemte Betingelser.
Man maa ikke forvexle den umiddelbare og uvilkaarlige Tryghed, hvormed Sindet
hviler i Sanse- og Erindringsanskuelser, med den Forvisning, der først er
Eftertankens Værk. En saadan Forvexling begaar \textsc{Heinrich
Maier}\footnote{\textit{Psychologie des emotionalen Denkens} (1908), p. 2. ---
Smlgn. om hele dette Spørgsmaal \textit{Den menneskelige Tanke}, p. 53---59; 74
f; 99---103.}, naar han udtaler: „Elementære Domme ligger i enhver Iagttagelse,
ethvert Erindringsbillede, enhver kognitiv Fantasiforestilling, forsaavidt alle
disse Forestillinger har virkelige Objekter, virkelige Begivenheder, Tilstande,
Ting o.~s.~v. til Genstand“. Og naar han føjer til: „De foregivne
Iagttagelses-, Erindrings- og Fantasibilleder, som skal være indifferente
overfor Modsætningen mellem Sandt og Falsk, Gyldigt og Ugyldigt er Fiktioner,
der aldrig og intetsteds er virkelige“, --- saa kan jeg kun give ham Ret heri,
forsaavidt hin umiddelbare Tryghed ikke kan kaldes Ligegyldighed. Om
Ligegyldighed kan man strengt taget kun tale, naar begge Modsætningsled kendes
og staar som lige gyldige. Men dette Kendskab er ikke fra først af tilstede.
Der er endnu ikke spist af Kundskabens Træ. Ingen Tvivl har endnu rørt sig;
Tilstanden er helt igennem positiv. Naar den umiddelbare Tryghed brydes gennem
Modsigelse eller gennem de Krav til nøjere Bestemmelse, nye Iagttagelser
stiller, saa begynder Eftertanken sit Arbejde for at tilvejebringe nye
Helhedsdannelser, og da kan Dommen afløse Anskuelsen. Da først bliver der
Anvendelse for Begreberne Sandt og Falsk som Modsætninger.

\medskip
4. Forskellen mellem det elementære, uvilkaarlige Tankearbejde og det
videnskabelige kan belyses fra tre forskellige Sider.

% ENUMERATOR (p. 8): italic Latin "a." -- the first of the "tre forskellige
% Sider" (b. p. 11, c. p. 12). At 500 dpi the italic a is nearly the same
% glyph as the Greek α of p. 9; UNSURE: rejected α.
\textit{a.} Eftertanken vil allerførst søge at ordne de forskellige Emner, den
skelner, efter deres indbyrdes Forhold. Relation, Forholdsbestemthed er derfor
(som Kategorilærens Historie ogsaa viser) den mest iøjnefaldende Kategori. Men
de forholdsbestemte Ordninger kan være af forskellig Betydning for Erkendelsens
Udvikling. Gennem en Række af Trin skrider Eftertanken frem fra rent udvortes
Ordninger, der ikke kan føre videre, til saadanne Ordninger, der gør det muligt
at udlede nye Forhold af de allerede foreliggende. Der er en hel Række af Trin
fra Kaos til Rationalitet, og Rationaliteten kan selv have forskellige Grader.
Disse Trin belyses ved de Rækker, der kan dannes ved Ordning af Emnerne. En
Oversigt over
% --- p. 9 ---
\setcounter{footnote}{0}%
\opage{9}disse Rækker har jeg givet i „Den menneskelige Tanke“ (p. 162---207);
men jeg kommer tilbage til dem her i al Korthed og fra et noget andet
Synspunkt.\footnote{Som Tillæg til nærværende Afhandling er der aftrykt en
Liste over Kategorierne, i det væsentlige som jeg har fremstillet dem i den
nævnte Bog. Og der er tillige anført de forskellige Rækkedannelser, der
udtrykker Overgangen fra den førvidenskabelige Bevidsthed til den
videnskabelige.}

α. Forsøger vi at tænke os de for menneskelig Tænkning mest uheldige Forhold,
staar vi overfor det logisk Irrationelle. Vi staar overfor et saadant Forhold
mellem givne Emner, at de alle er lige meget forskellige fra hverandre
indbyrdes, og at de derfor kan anbringes i hvilkensomhelst Orden. Man kan skyde
et hvilketsomhelst Led ud, uden at Forholdet mellem de tilbageblivende Led
ændres, og uden at der overhovedet følger noget deraf. En saadan Række har jeg
kaldt en kaotisk Forskelsrække. Den fremtræder ikke i umiddelbar Iagttagelse.
Thi allerede rent psykologisk er Addendernes Orden ikke ligegyldig. Vore
sjælelige Tilstande bestemmes, baade hvad Grad og Kvalitet angaar, ved deres
Plads i vore Oplevelsers Række. Men der kan gives alle mulige Tilnærmelser til
en saadan logisk Irrationalitet, der er forskellig fra den matematiske
Irrationalitet derved, at den ikke som denne ved selve sin Tilblivelsesmaade
kan føre til Reduktion af Forskellighederne. Vi ved, at $\sqrt{2}$ ligger i
Tallenes Række mellem 1 og 2 og ikke andre Steder, og indenfor det nævnte
Mellemrum kan Pladsen bestemmes stedse nøjere, jo flere Decimaler vi udregner.
Det logisk Irrationelle betyder derimod en Stansning.

Det næste Skridt er, at der dannes Rækker, hvis Led ligger fast, idet
Forskelsforholdene varierer fra Led til Led. En saadan ubestemt varierende
Forskelsrække tillader ikke Udskydning, uden at de Led, mellem hvilke
Udskydningen finder Sted, kommer i andet Forhold til hinanden end før. Men
Fremskridtet bestaar i, at klar Oversigt har afløst den kaotiske Uro.

Et nyt Fremskridt gøres, naar Forskellighederne mellem Leddene varierer i en
bestemt Retning indenfor Rækken. En saadan regelmæssig varierende
Forskelsrække muliggør for Eftertanken Dannelsen af en Lov eller Regel, efter
hvilken Emnernes Plads i Rækken kan bestemmes.
% DEFECT (p. 9): a stray low dot stands before "Og" (". .Og"); a loose
% sort, not transcribed.
Og den Opgave kan stilles at finde nye Led i Rækken, Mellemled mellem de
allerede givne.

Hvis Forskellen mellem Led, der har deres bestemte Plads i en Række, bestandig
er den samme, Led og Led imellem, faar vi en identisk varierende Forskelsrække.
Identitetsbegrebet begynder at dæmre, foreløbig som Identitet af Forskellen
mellem kvalitativt varierende Emner. Saadanne Rækker kan dannes af
Øjeblikkene, Tallene, Graderne og Stederne. Hvert Øjeblik, med dets
ejendommelige Oplevelse, har en bestemt Plads i Øjeblikkenes Række, og hvert
Øjeblik staar i samme Forhold til det foregaaende og det følgende Øjeblik.
Ligesaa med Tallene. Hvert Tal har fra først af, som \textsc{Poincaré} har
udtrykt det, sin egen Individualitet, er kvalitativt forskelligt fra andre
Tal, ikke blot ved sin Plads i Talrækken. At man ad forskellige Veje kan komme
til samme Individualtal, er en Opdagelse, der forudsætter et højere Trin end
det, vi her staa ved. \textsc{Goethe} mente, at Matematiken bevæger sig i
lutter Identiteter, fordi $2 + 2 = 4$, men ogsaa $4 = 2 + 2$, saa at vi
% SIGNATURE LINE at the foot of p. 9 ("D. K. D. Vidensk. Selsk. Skr., 7.
% Række, hist. og filos. Afd. III, 2." and "34") omitted as printer's apparatus.
% --- p. 10 ---
\opage{10}ikke faar andet, end vi havde forud. Men de forskellige Veje, ad
hvilke man kan komme til et Tal (f.~Ex. ved at tælle til 4 eller lægge 2 og 2
sammen), er forskellige Operationer, om hvilke man først bagefter opdager, at
de fører til samme Resultat og forudsætter de samme Principer. Tallet er
foreløbig Ordenstal, Nummer. Ogsaa Gradsforskellene ere umiddelbart
Kvalitetsforskelligheder, og dog ligger de i en bestemt Orden i en Række, hvis
Led kan variere paa identisk Maade. Ligesaa forholder det sig med Stederne i
Rummet. Selv om de maaske endnu betragtes som kvalitativt forskellige, saa er
det en identisk Varieren, ved hvilken der gaas fra Sted til Sted.

β. I de hidtil omtalte Rækker er Leddene kvalitativt forskellige, og det drejer
sig blot om deres indbyrdes Forhold med Hensyn til Plads i en Række. Det er et
stort Fremskridt, naar denne Plads bliver aldeles bestemt. Klassifikation og
Sammenligning bliver mulig derved. Men virkelig udføres kan disse Operationer
dog kun, naar der kan dannes Rækker, i hvilke hvert Led er indbefattet i det
foregaaende. Dette kan igen fremtræde paa to Maader. I en fremskridende
Forskelsrække kan f.~Ex. ethvert Led være Del af det foregaaende. I en
Klassifikation er Arten en Del af Slægten, og Individet en Del af Arten. Deraf
kan sluttes, at naar et Individ A hører til Arten B, og Arten B til Slægten C,
hører A ogsaa til C. Det er væsentlig herpaa, den aristoteliske Slutningslære
bygger. I en partiel Identitetsrække forholder det følgende Led sig til det
foregaaende som Prædikat til Subjekt, eller som Følge til Grund. Ogsaa her kan
Klassifikation bruges som Exempel, thi at A hører til Arten B, vil sige, at de
Prædikater, B betegner, tilkommer A. Ved begge de nævnte Rækker sker
Slutningen ved Udskydning af Mellemled.

Naar en Række, i hvilken Leddene ligger fast, er symmetrisk, d.~v.~s. ligesaa
godt kan læses bagfra som forfra, kan der fra den sluttes til den omvendte
Række.

% ENUMERATOR (p. 10): γ read at 500 dpi; the text layer drops it.
γ. Fremskridende Forskelsrækker, partielle Identitetsrækker og symmetriske
Rækker kan kaldes rationelle Rækker, idet de gør Slutning mulig. Men Leddenes
kvalitative Forskellighed er ikke ophævet, fordi de kunne ses som Dele eller
Helheder, eller som Subjekter og Prædikater i Forhold til hverandre. Den
fuldkomneste Art af Slutning bliver først mulig, naar ikke blot Forholdet
mellem Leddene stedse er identisk, men selve Leddene ogsaa er det, saa at ikke
blot Led kan skydes ud, men ethvert Led kan indsættes istedetfor ethvert andet.
Vi faar da en absolut
% FORMULA (p. 10): inline in the print, followed by five spaced dots.
Identitetsrække: $A = B = C = D$ .~.~.~.~. Her er, som først \textsc{Leibniz}
ret saa, Grundlaget --- og tillige Idealet --- givet for al rationel Tænken. En
saadan Række betegner den fuldstændige Modsætning til den kaotiske
Forskelsrække. De to Rækker betegner de ideale Grænser, mellem hvilke
menneskelig Erkendelse bevæger sig.

Skridt for Skridt kommer vi saaledes ved at gennemløbe de Trin, som de
forskellige Rækker betegner, op til de Forudsætninger, streng Videnskab hviler
paa. Men vi ser tillige, at de rationelle Rækker ikke er de eneste mulige, og at
der indenfor de rationelle Rækker maa skelnes mellem saadanne, der blot gør
Klassifikation og Sammenligning, og saadanne, der gør Slutning i strengeste
Forstand mulig. De forskellige Rækker danner en Skala, og det vil maaske være
muligt at finde endnu flere Trin paa
% --- p. 11 ---
\setcounter{footnote}{0}%
\opage{11}denne Skala, end der er nævnt i det foregaaende. Tankearbejdet, der gaar
ud paa Ophævelse af det kaotiske Forhold mellem Emnerne, frembyder en hel Række
af Grader, indtil det rent rationelle Trin naas, hvor ikke blot Leddenes
Forhold, men selve Leddene er identiske. Dette højeste Trin staar dog stedse kun
som et Ideal, der kun kan naas tilnærmelsesvis og under Bortseen fra en Mængde
faktisk foreliggende Forskelligheder. Gennem Ordning, Klassifikation og
Sammenligning arbejder videnskabelig Bevidsthed sig op imod sin højeste Form,
overalt hvor Emnerne gør det muligt.

\textsc{Kant} tog i „Kritik der reinen Vernunft“ kun Hensyn til det højeste
Trin af Videnskab, altsaa til det Trin, hvor det var muligt at drage Slutninger
% PRINTED AS IS (p. 11): "absotut" for "absolut" (checked at 600 dpi).
paa Grundlag af absotut Identitet. Klassificerende og sammenlignende Videnskab
tog han intet Hensyn til; den regnede han til den blotte Empiri. Og overfor den
Indvending, at Klassifikation og Sammenligning jo dog ikke mindre forudsætter
Gyldigheden af logisk Tænkning end den matematiske Naturvidenskab, svarer han,
at den rene Logiks Principer vel er negative Kriterier for Sandhed, men ikke
% PRINTED AS IS (p. 11): "tilsrækkelige" for "tilstrækkelige" (checked at 600 dpi).
tilsrækkelige til at begrunde den Art Videnskab, han vil undersøge Muligheden
af\footnote{\textit{Kritik der reinen Vernunft}\textsuperscript{1}, p. 151 f.}.
Men vi tillægger jo dog vore Klassifikationer og Sammenligninger objektiv
Gyldighed, ligesaavel som de egentlige Naturlove, og der er derfor ligesaa megen
Grund til at undersøge Berettigelsen af at antage rent logiske Principers
objektive Gyldighed, som der er til at undersøge de matematiske Principers og
Aarsagsprincipets objektive Gyldighed. Allerede \textsc{Salomon Maimon} har
indset dette, idet han overfor \textsc{Kant} hævdede, at de logiske
Grundformer, Bekræftelse og Benægtelse, ingen Betydning har, naar de ikke
forudsættes gyldige for de Emner, de anvendes paa. Medens \textsc{Kant} lader
Erkendelsesteorien („Transcendentalfilosofien“) forudsætte den formale Logik som
en en Gang for alle færdig Disciplin, hævder \textsc{Maimon}, at
Erkendelsesteorien maa gaa forud for den formale Logik, naar denne skal have
real Betydning.\footnote{\textit{Versuch einer neuen Logik oder Theorie des
Denkens}, Berlin 1794, p. 404---408.} --- \textsc{Kant} havde naturligvis Ret
til at begrænse Omfanget for sine Undersøgelser, men han havde ikke Ret til at
mene, at han derved tillige havde begrænset Videnskabens Omfang.

\textit{b.} Fra en ny Side fremtræder Overgangen fra den elementære,
førvidenskabelige Bevidsthed til den videnskabelige, naar vi fremdrager en
Forudsætning, som al den Art Rækkedannelse, der har været Tale om i det
foregaaende, hviler paa. Vi have i det foregaaende kun taget Hensyn til Kampen
mod Forskellighederne og set, hvorledes Lighedsgraderne, baade hvad Forhold og
hvad Led angaar, Skridt for Skridt nærmer sig til ren Identitet. Men al Forskel
og Lighed, og dermed al Rækkedannelse, forudsætter Muligheden af, at Emner kan
samles, sammenfattes og ordnes. Her kommer Begrebet Syntese frem som en
Forudsætning, der gælder baade for formal og for real Videnskab, baade for
Klassifikation og for Loverkendelse. Rækkedannelse beror paa, at der paavises
bestemte Forhold mellem Emnerne, og Forhold, Relation tænkes kun, naar Leddene,
mellem hvilke Relationen gælder, er Genstand for en og samme Tankeakt. Relation
forudsætter Syntese. Men da
% SIGNATURE MARK at the foot of p. 11 ("34*") omitted as printer's apparatus.
% MARGIN MARKS (p. 11): several small arrow-like marks (►) stand in the left
% margin of the scanned copy; not text of the work, not transcribed.
% --- p. 12 ---
\setcounter{footnote}{0}%
\opage{12}Syntese, Sammenfatten netop bestaar i at bringe Emner i et eller andet
Forhold til hverandre, er Relation og Syntese korrelate
% PRINTED AS IS (p. 12, note): "155 f" without a stop after f.
Begreber\footnote{I \textit{Den menneskelige Tanke} p. 155 f opstillede jeg
Syntese som den første Kategori, Relation som den anden. En skarpsindig Læser
har gjort mig opmærksom paa, at de burde være opstillede som Korrelatbegreber,
ligesom Kontinuitet og Diskontinuitet, Lighed og Forskel.}.

Nu er de i det foregaaende omtalte Rækker kun typiske Exempler paa, hvad der kan
finde Anvendelse i mange specielle Tilfælde. Der kan foretages mange forskellige
Klassifikationer, anlægges mange forskellige komparative Synspunkter, findes
mange forskellige Love. Men alle disse maa igen bringes i Sammenhæng med
hverandre, og herved vil nye Rækkedannelser komme frem. Der forudsættes da
Synteser af højere Orden end ved de enkelte Rækker hver for sig. Det gælder da
om, at Syntesens Omfang bliver saa stort som muligt. Der gør sig her en Stræben
gældende henimod et idealt Maximum. Der vil her være større Modstand at
overvinde end ved de enkelte Rækker, da Forskellen mellem de Rækker, der kan
dannes paa de forskellige Omraader, vil være større end Forskellen mellem Led i
samme Række og paa samme Omraade. Spørgsmaalet bliver tilsidst, om vor Tænkning
raader over logiske Former for fuldstændig Sammenfatten, eller om ikke mere
eller mindre dristige og mere eller mindre poetiske Analogier paa et vist Punkt
afløser den logiske Tænkning. I nærværende Afhandlings sjette Kapitel ville vi
vende tilbage til dette Spørgsmaal. Men før den Grænse er naaet, hvor Tænkning
gaar over til Poesi, er der endnu et Omraade, hvor videnskabelige Hypoteser med
fuld Ret bevæger sig. Trangen til Sammenfatten vil stadig gøre sig gældende,
ikke blot overfor Emner, men ogsaa overfor Rækker. Og der kan ved Forsøg paa
stedse mere omfattende Syntese kastes Lys over de enkelte Rækker og de enkelte
Emner, ligesom de enkelte Emner allerede belyses ved den Plads, de indtager i den
Række, de er Led i. Syntese og Analyse (og analogt: Kontinuitet og
Diskontinuitet, Lighed og Forskel) staar i stadig Vexelvirkning med hinanden, og
Videnskabens rytmiske Svingninger finder deres naturlige Forklaring ved denne
Grundlov for den menneskelige Tanke.

\textit{c.} Et tredie Synspunkt til Karakteristik af Forskellen mellem
førvidenskabelig og videnskabelig Tænkning er givet deri, at jo mere Videnskaben
udvikler sig, des større Differentiation af de forskellige Problemer finder der
Sted. I den uvilkaarlige Tænken skelnes der ikke mellem Emne, Form og Værdi. Man
gør sig ikke kritisk klart, hvor Grænsen ligger mellem det umiddelbart Givne
(Emnet) og den Tankebearbejdelse, man anvender ved dets Tydning. Og det bliver
heller ikke erkendt, hvorledes allerede Emnets første Optræden for os betinges
ved vor Modtagelighed, der aldrig er absolut passiv. Heller ikke skilles Værdi
bestemt fra Emne og Form. Værdisynspunktet vil under primitive Forhold have
været det overvejende, idet Emner især bemærkes paa Grund af deres praktiske
Betydning, og idet ogsaa den Maade, hvorpaa Emnet tydes, betinges ved dets
Sammenhæng med praktiske Interesser. \textsc{Bergson} burde derfor have regnet
Værdielementet med til „les données immédiates“; han har forsøgt at eliminere
Formelementet og mener, at hvis det kunde udskilles, blev Emneelementet alene
tilbage. Han er
% --- p. 13 ---
\opage{13}endog tilbøjelig til at mene, at Grunden til Formelementets
Fremhersken i Videnskaben maa søges i praktiske Hensyn. Men allerede i
Emnevalget medvirker Interesse. Naturligvis medvirker den ogsaa i den primitive
Tydning. Det teleologiske Synspunkt raader i Animismen under alle dens Former,
lige fra Fetishismen op i den mest sublime Teologi. Og det kan i de primitive
Former eller Antydninger af, hvad vi nu kalder Kategorier, spores, at praktiske
Motiver har været medvirkende ved deres Opstaaen. Rummet betyder først Afstand
fra et tilstræbt Maal, Tiden et Forhold mellem en ønsket Fremtid og en
utilfredsstillende Nutid, og Aarsagsforholdet gør sig gældende under Forsøget
paa at ophæve disse Afstande i Rum og Tid. Under Videnskabens Udvikling spores
en stadig Stræben efter at eliminere Værdielementet, at se bort fra al Interesse
med Undtagelse af den intellektuelle Interesse, der baade som Motiv og som
Virkning er paa inderligst Maade knyttet til selve Tankearbejdet. Dermed falder
andre Værdier ikke nødvendigvis bort, men Værdibegrebet kan nu blive Genstand
for et særligt Problem, der bl.~A. ogsaa drøfter Forholdet mellem intellektuel
Værdi og andre Værdier.

Sondringen mellem de tre Elementer betegner et vigtigt Fremskridt. Den er endnu
ikke strengt gennemført i Videnskaben, og endnu mindre i den sædvanlige
Bevidsthed. Denne Sondring løser naturligvis ingen Problemer. Tvertimod er dens
Betingelse for, at Problemer kan stilles skarpt og bestemt, og selve Brydningen
mellem de tre Elementer stiller tilsidst det største Problem af alle. ---
Iøvrigt fører dette fra Differentiationen hentede Synspunkt tilbage til de
foregaaende Synspunkter, idet der ved stigende Differentiation stilles stigende
Fordringer til Syntese, dersom Aandslivets Enhed skal bevares, og idet
Differentiationen foranlediger nye Relationer.
% MARK (p. 13): a stray dot stands below "Relationer."; not transcribed.

\medskip
5. Psykologiens Selvstændighed overfor Erkendelsesteorien fremgaar af, at
videnskabelig Tænken ikke kan emancipere fra Bevidsthedslivets almindelige
Love. (Smlgn. min Psykologi\textsuperscript{6} V B, 11). Baade Opdagelse og
Bevisførelse skyldes psykisk Arbejde. De maa være psykologisk mulige. Man kan
ikke deducere Erkendelsesteori ud fra Psykologien, men man kan undersøge, hvilke
psykologiske Betingelser videnskabelig Erkendelse forudsætter, og derefter
undersøge, om disse Betingelser er tilstede. Psykologiens Stilling overfor
Videnskaben er ganske analog med dens Stilling overfor Kunst og Religion. Den
spørger om, hvorledes Videnskab er mulig, ligesom den spørger om, hvorledes
Kunst og Religion er mulig. Disse forskellige aandelige Livsformer vokse ud af
det uvilkaarlige Sjæleliv og taber, saalænge Betingelserne for deres Bestaaen er
tilstede, ikke Sammenhængen med dette, trods den Ejendommelighed, hvormed de
fremtræder under deres højere Udvikling.

% PRINTED AS IS (p. 13): „Sammenbængen“ for „Sammenhængen“ (b for h;
% second reading 2026-09-30, clear at 500 dpi).
Naar Erkendelsesteorien fremstilles syntetisk, viser Sammenbængen med
Psykologien sig tydeligst. Psykologien fremtræder da som en Indledning, som en
% PRINTED AS IS (p. 13): "der Geistes" (for Hegel's "des Geistes").
Art „Phänomenologie der Geistes“. Ved Undersøgelse af det uvilkaarlige
Tankearbejde bliver vi opmærksomme paa det Forhold mellem Emne og Tydning, der i
videnskabelig Tænkning spiller en afgørende Rolle. De Former, under hvilke
Emnerne
% --- p. 14 ---
\opage{14}bearbejdes ved uvilkaarlig Tankevirksomhed, maa tilslibes, udvides
eller begrænses, for at det Maal, Videnskaben stiller sig, nemlig at vinde
Tanker, der kan bestaa trods Emnernes Skiften, kan naas.

\textsc{Kant} gaar i „Kritik der reinen Vernunft“ syntetisk frem. Hans
„metafysiske“ Deduktion i den „transcendentale Æstetik“ og hans „subjektive“
Deduktion i den „transcendentale Analytik“ er i Virkeligheden psykologiske
Analyser, ved hvilke Begrebet Form udskilles fra Begreberne Emne og Værdi. I min
Bog „Den menneskelige Tanke“ gaar jeg ogsaa syntetisk frem, kun med større Vægt
paa det psykologiske og historiske Grundlag. Jeg søger at vise, hvorledes
allerede Overgangen fra Anskuen til Dømmen indeholder Spiren til videnskabelig
Tænken, og hvorledes denne gennem Videnskabens Historie vaagner til stedse
klarere Bevidsthed om sig selv. Deraf fulgte Inddelingen af min Fremstilling i
tre Afsnit: Funktioner, Kategorier og Problemer.

Men ogsaa ved en analytisk Fremstilling af Erkendelsesteorien vil det stedse
vise sig, at der er en nødvendig Sammenhæng mellem Erkendelsesteori og
Psykologi. Man gaar da ud fra den faktisk bestaaende Videnskab og spørger,
hvilke Forudsætninger den hviler paa, og hvilke Tankeformer den virker med. Og
derefter kan man ikke undgaa at spørge, hvorvidt disse Forudsætninger og Former
faktisk er tilstede i Bevidsthedslivet, og hvorledes de forholde sig til dettes
uvilkaarlige Tankearbejde. Denne Vej gik \textsc{Kant} i „Prolegomena“. Vejen
fører her fra Grundsætningerne gennem Kategorierne tilbage til de psykiske
Funktioner, som det er Psykologiens Sag at paavise og at beskrive. Denne Trehed
træder tydelig frem i nyere erkendelsesteoretiske Værker, der gaar analytisk
frem ligesom „Prolegomena“. \textsc{Meyerson} (Identité et Réalité. 1908)
kommer gennem Analyse af moderne exakt Naturvidenskabs Forudsætninger tilbage
til Modsætningen mellem Identitet og Succession. Kausalitet i streng Betydning
indeholder efter hans Opfattelse tvesidig Ækvivalens, Mulighed for Substitution,
altsaa Identitet. Men der kan være Lovmæssighed (légalité) uden tvesidig
Ækvivalens; her er Substitution ikke mulig, og derfor vil \textsc{Meyerson} ikke
bruge Udtrykket Kausalitet her. \textsc{Carnot}'s Princip (Entropien) viser
Betydningen af denne Forskel. \textsc{Meyerson} finder gennem Videnskabens
Historie en stadig, ubetvingelig Tendens i den menneskelige Aand til at danne
Begreber om noget, der bestaar trods al Succession. Men naar nu Tiden ikke kan
elimineres, hvilket \textsc{Carnot}'s Princip tydelig godtgør ved at paavise en
Grænse for Substitution af Tilstande i Naturen, saa staar vi ved Modsætningen
mellem Identitet og Succession, og den dermed sammenhængende Modsætning mellem
Kvantitet og Kvalitet, der tilsidst viser tilbage til Forholdet mellem Tænkning
og Sansning, som det er Psykologiens Opgave at undersøge. Altsaa kommer vi ogsaa
her tilbage fra Grundsætninger gennem Kategorier til psykiske Funktioner. Ogsaa
hos \textsc{Cassirer} (Substanz und Funktion. 1910) fremtræder tydelig
Forskellen mellem Funktioner, Kategorier og Problemer, i omvendt Orden af den,
jeg, formedelst syntetisk Fremstilling, følger i min Bog. \textsc{Cassirer}'s
Analyse af Videnskabens Forudsætninger fører ham til at indskærpe Betydningen af
Begrebet Relation, der for ham, som
% --- p. 15 ---
\opage{15}for saa mange nyere Erkendelsesteoretikere, bliver den vigtigste
Kategori. Og derved stilles der tilsidst Psykologien den Opgave at undersøge
Relationen, Forholdsbestemthedens Rolle indenfor Sjælelivet. Med Urette mener
\textsc{Cassirer}, at denne Undersøgelse hidtil har været aldeles forsømt. Men
Hovedsagen er i denne Sammenhæng, at der er en nøje Forbindelse mellem
Erkendelsesteori og Psykologi. Ja, \textsc{Cassirer} ender endog sin Bog med den
Udtalelse, at tilsidst falder Modsætningen mellem dem bort, idet „Psykologien
selv giver et Tilløb (Ansatz) til de Problemer, som maa søge deres
fremskridende Løsning i Logiken og i dennes Anvendelse paa Videnskaben“. Ogsaa
i \textsc{Henri Poincaré}'s erkendelsesteoretiske Afhandlinger finder vi den
omtalte Trehed. Hans Undersøgelser af den moderne Matematiks og Fysiks Teorier
fører ham tilbage til Modsætningen mellem Kontinuitet og Diskontinuitet og deres
stadige Brydning indenfor Tankens og Videnskabens Udvikling. Og denne Modsætning
er det Psykologiens Opgave at undersøge, idet den hænger sammen med Modsætningen
mellem Anskuen og Tænken, en Modsætning, der stedse paany maa gøre sig
gældende. Derfor nægter \textsc{Poincaré} (se især „Dernières Pensées“, p. 139.
158), at der gives en af al Psykologi uafhængig Erkendelsesteori: en saadan er
ligesaa umulig, som at der skulde kunne gives Videnskab uden Videnskabsmænd!

Dersom det uvilkaarlige Tankeliv ikke forberedte det videnskabelige Tankeliv,
vilde Videnskaben miste sin Naturgrund. Videnskaben vokser frem af det
uvilkaarlige Tankeliv. Der er da ogsaa en hel Gruppe af Kategorier, de
fundamentale Kategorier, som allerede det uvilkaarlige Tankeliv opererer med:
Syntese og Relation, Kontinuitet og Diskontinuitet, Lighed og Forskel er Former,
intet Bevidsthedsled kan undvære. De afgiver en Definition paa, hvad det er „at
tænke“ i Ordets videste Betydning, den Betydning, i hvilken \textsc{Descartes}
efter sin egen Forklaring tog Ordet, naar han sagde „je pense, donc je suis“. De
andre Grupper er specielt udformede Anvendelser af de fundamentale Kategorier. I
Klassifikation og i logisk og matematisk Tænken lægger de formale Kategorier
(Identitet og Rationalitet) sig for Dagen. I Naturvidenskab, Psykologi og
Sociologi arbejdes med de reale Kategorier (Kausalitet, Totalitet, Udvikling). I
Økonomi, Æstetik, Etik og Religionsfilosofi ligger Værdibegreber til Grund. En
rationel Klassifikation af Videnskaberne er et med en systematisk Kategorilære.

De formale Kategorier kom allerede til Bevidsthed i Oldtiden. Her ligger den
græske Tænknings Stordaad. I nyere Tid kostede det en Kamp, inden de reale
Kategorier kunde ses i deres karakteristiske Forskel fra de formale. Her ligger
\textsc{Hume}'s og \textsc{Kant}'s store Indsats. Og det har kostet --- og
koster --- stadig et Arbejde at udsondre Værdibegreberne fra de andre Kategorier
og dog se dem i deres Sammenhæng med dem. Problemernes og Kategoriernes Historie
staar i nøje Forbindelse indbyrdes.

Nogle Exempler vil nærmere oplyse baade Forholdet mellem de forskellige Grupper
af Kategorier og Forholdet mellem Psykologi og Erkendelsesteori.

Det er ved Analyse af de i psykologisk Erfaring givne Tilstande, at Erken%
% --- p. 16 ---
\opage{16}delsesteori forberedes. I de psykiske Tilstande er der givet andet og
mere, end Erkendelsesteorien behøver. Det er Opgaverne, som bestemmer, hvad der
vælges indenfor det foreliggende Indbegreb af Muligheder. Der fremtræder
f.~Ex. i psykologisk Erfaring forskellige Grader eller Arter af Lighedsforhold.
Baade ved Genkendelse, Association og Sammenligning lægger Lighedsforholdet sig
for Dagen, dels som Dækningslighed, dels som sammensat Lighed (Dækningslighed
$+$ Forskel), dels som Kvalitetslighed, dels som Forholdslighed. Disse Ligheder
faar erkendelsesteoretisk Betydning derved, at de gør Dannelsen af Rækker mulig,
de fastliggende Rækker, i hvilke Leddene er ordnede efter aftagende Lighed eller
tiltagende Forskel, og ved hvilke Klassifikation og sammenlignende Videnskab
bliver mulige. Men Logik og Matematik kræver Dækningsligheden i dens fuldkomne,
ideale Form som absolut Identitet, fordi kun derved streng Slutning og Beregning
bliver mulig. Identitet defineres da som et Grænsetilfælde af Lighed. Det er en
analytisk Definition, vi her anvender, ligesom naar vi i Geometrien gaa fra
Legemet („Klodsen“) til Fladen, fra Fladen til Linien („Kanten“) og fra Linien
til Punktet. Begyndelsen er psykologisk (eller ligger paa et for Psykologi og
Erkendelsesteori fælles Omraade), men ud fra dette Begyndelsespunkt udvikler det
erkendelsesteoretiske Synspunkt sig successivt til Selvstændighed. Denne
Udvikling falder for en Del sammen med en Udvikling fra Praxis til Teori. I
Praxis kan vi ofte lade os nøje med et Skøn, en Lighedsopfattelse af en Art, der
ligger mellem Dækningslighed og Forholdslighed. Men den strenge Teori kræver
Identitet som Grundlag for Slutning og Maaling. Det er en saadan Identitet,
\textsc{Platon} mener med, hvad han kalder „Lighedens Ide“. De platoniske Ideer
er baade Almenbegreber og Grundbegreber; \textsc{Platon} følte ingen Trang til
at sondre skarpt mellem Klassifikation og exakt Tænkning, mellem beskrivende og
formal Videnskab. Ligesom for \textsc{Platon} den højeste Grad af Lighed er en
„Ide“, til hvilken vi i alle givne Tilfælde kun har Tilnærmelser, saaledes er
ifølge \textsc{Kant} det rene Rum en „Ide“, en Regel, efter hvilken Fænomener
kan ordnes og bedømmes, men som ikke selv er Fænomen. Striden mellem de
forskellige erkendelsesteoretiske Retninger drejer sig egentlig om, hvorvidt det
% DEFECT (p. 16): the e of "er" is damaged (reads like "ęr"); set as "er".
er nødvendigt at gaa til Grænsen --- til Identitet som Lighedens Ideal. Alle er
vel enige om, at absolut Identitet ikke kan paavises i nogensomhelst
Iagttagelse. Det gælder blot, om man straks vil gaa til den praktiske Anvendelse
af Videnskaben, eller om man først vil fremstille Videnskaben i den strenge Form,
den faar ved, at man gaar til Grænsen, og derefter stige ned til de i Erfaring
og Praxis mulige Tilfælde. Det er et Forhold mellem Helhed og Del, eller mellem
en Række af stigende Lighedsgrader og denne Rækkes sidste Led, som her ligger
til Grund. De forskellige Definitioner, som er givne af Identitet, har det
fælles, at de alle forudsætter Lighed og Forskel som kendte, som umiddelbart
foreliggende, saasnart Eftertanken er begyndt at virke, og at man indenfor de
Muligheder, Ligheds- og Forskelsforholdet indeholder, udtager den, der kan blive
frugtbar for Tanken. Naar \textsc{Leibniz} definerer Identitet ved Mulighed af
Substitution, saa vil det jo sige, at han holder sig til den Grad af Lighed, der
gør det muligt at sætte det Ene istedetfor
% --- p. 17 ---
\opage{17}det Andet. Ligesaa i en berømt Matematikers Definition af Identitet
som den Forskel, der er mindre end enhver given Forskel; her lægges blot en
Række af aftagende Forskelligheder til Grund. Den mindst mulige Forskel er en
„Ide“, ligesom den størst mulige Lighed. En tredie Definition
(\textsc{Bolzano}'s), der forklarer Identitet som en Tings Lighed med sig selv,
fører tilbage til \textsc{Leibniz}'s Definition, da det er saadan Lighed, der
betinger Substitution. Eller ogsaa fører den tilbage til den mindst mulige
Forskels Begreb; thi mellem en Ting og den selv er der i Regelen mindre Forskel
end mellem en Ting og andre Ting. Og da ingen Ting fremtræder ganske paa samme
Maade i forskellig Sammenhæng, bliver Identiteten ogsaa efter denne Definition
en „Ide“.

Det er ikke Psykologien ligegyldigt, om man definerer Identitet ved Lighed eller
Lighed ved Identitet. I sidste Tilfælde bliver Definitionen syntetisk; man
opfatter Identitet som et simpelt Begreb og Lighed som Identitet $+$ Forskel.
Men herimod maa Psykologien, og i Grunden al Erfaringsvidenskab protestere. Thi
der gives Ligheder, der ikke kan opløses paa denne Maade, som f.~Ex. Ligheder
mellem to Farver (Orange og Gult, Blaat og Violet). Ikke blot
Kvalitetslighed, heller ikke Forholdslighed (Analogi), kan uden videre opløses
paa denne Maade. ---

Et andet Exempel paa Overgang fra Psykologi til Erkendelsesteori gennem Analyse
haves i Forholdet mellem Virkelighedskriteriet og Aarsagsbegrebet. Psykologien
maa for sit eget Vedkommende stille sig det Spørgsmaal, hvorledes Mennesker bære
sig ad, naar de kommer til at tvivle om deres Iagttagelsers eller Forestillingers
Gyldighed, d.~v.~s. om Muligheden af at holde dem fast og at handle efter dem i
alle Tilfælde. Nu viser Erfaring, at vi i Tvivlstilfælde uvilkaarlig giver os
til at prøve, om de enkelte Iagttagelser og Forestillinger stemmer overens
indbyrdes og med de nye Iagttagelser, som kan gøres. Kun hvor der viser sig en
uundgaaelig Sammenhæng mellem alle de forskellige Iagttagelser, der staar til
Raadighed, stoler vi paa, at vi staar overfor en „Virkelighed“. Men en saadan
uundgaaelig Sammenhæng er netop Aarsagsbegrebets mest elementære Form.
Videnskaben bliver ikke staaende ved denne elementære Form, men søger, overalt
hvor det er muligt, at gøre Sammenhæng absolut rationel, saaledes at Leddene
blive tvesidig ækvivalente. I det ideale Kausalitetsbegreb har Rationaliteten
Overvægten over det rumlige eller tidslige Forhold mellem Leddene. ---
Psykologisk spørges der, hvorledes vi bærer os ad i Tvivlstilfælde;
erkendelsesteoretisk spørges der, med hvilken Ret vi slutter fra Fortid eller
Nutid til Fremtid, eller fra Nutid til Fortid.

\medskip
6. Det er ingen Indvending mod den her givne Bestemmelse af Forholdet mellem
Psykologi og Erkendelsesteori, at selve Psykologien, ligesom enhver anden
Videnskab, forudsætter et erkendelsesteoretisk Grundlag. Erkendelsesteorien, som
undersøger al Videnskabs Grundlag, undersøger naturligvis ogsaa Psykologiens.
Psykologiens videnskabelige Karakter beror paa, at der paa det Omraade, den
dyrker, kan dannes Rækker, som kan faa videnskabelig Betydning. Det er
foreløbig Rækker med fastliggende Led, der vil kunne dannes paa det psykologiske
% SIGNATURE LINE at the foot of p. 17 ("D. K. D. Vidensk. Selsk. Skr., 7.
% Række, hist. og filos. Afd. III. 2." and "35") omitted as printer's apparatus.
% --- p. 18 ---
\opage{18}Omraade; absolute Identitetsrækker, og derved Mulighed for Paavisning
af exakte Ækvivalensforhold, ligger endnu fjernt.

Der er ingen Modsigelse i, at Erkendelsesteorien, der selv historisk har
udviklet sig af Psykologien, senere vender sig mod sit eget Ophav, bl.~a. for at
forstaa sin egen Tilblivelse, ligesaalidt som der overhovedet er nogen
Modsigelse i, at Videnskabens Historie bliver Genstand for Videnskab. Det er jo
ikke Erkendelsesteoriens Opgave at frembringe de specielle Videnskaber. Den tager
dem som givne, skønt de jo hviler paa Forudsætninger, som det er
Erkendelsesteoriens Opgave at undersøge.

Den Vanskelighed, man her ofte finder, vil maaske forsvinde, naar man betænker,
at selve Erkendelsesteorien jo hviler paa visse Forudsætninger, saasnart den
skal være Videnskab. Der kan altsaa gives en Erkendelsesteori om
Erkendelsesteorien, en Erkendelsesteori i anden Potens. Naar
Erkendelsesteorien f.~Ex. hævder, at al Erkendelse frembyder et Forhold mellem
et Subjekt ($S$) og et Objekt ($O$), kan selve dette Forhold og dets Led ($S$ og
$O$) igen tages op og betragtes som et erkendelsesteoretisk Problem. Dette sker
jo i Grunden, naar Erkendelsesteoretikere kritiserer hinanden. Naar man ikke,
som \textsc{Kant}, mener, at det erkendelsesteoretiske Arbejde snart kan
afsluttes, maa man indrømme, at der her er Mulighed for en uendelig Række. Man
kan nu engang ikke én Gang for alle bryde Skepticismens Brod. Det er en
Grundejendommelighed ved det aandelige Liv, at det forudgaaende aandelige Arbejde
kan blive Genstand for nyt aandeligt Arbejde. Subjektet kan gøre sig selv til
Objekt. Hvis Tegnet $\lbrace$ udtrykker et Objektiveringsforhold, kan vi faa en
uendelig Række:
% FORMULA (p. 18): displayed and centred in the print; the sign is a left
% curly bracket, followed by a row of dots.
\[ S_1 \lbrace S_2 \lbrace S_3 \lbrace S_4 \ldots\ldots \]

Der er i denne Række ganske set bort fra $O$. I $O$ kan der indtræde
Forandringer, der ikke afhænger af $S$ alene, ligesom der jo ogsaa i $S$ kan
indtræde Forandringer, der ikke afhænger af $O$ alene. Forholdene vil altsaa i
Virkeligheden være mere komplicerede end i den anførte Række. Men denne er
tilstrækkelig til at vise, hvad det her kommer an paa. --- Smlgn. om hele det
her berørte Spørgsmaal: „Den menneskelige Tanke“ p. 297---305.

\medskip
7. Som for alle Begreber, der er uddragne umiddelbart af Anskuelse (Sansning,
Erindring eller Fantasi), gælder for Kategorierne den Regel, at der er et
omvendt Forhold mellem Indhold og Omfang. De mere omfattende Kategorier er de
fundamentale (Syntese og Relation, Kontinuitet og Diskontinuitet, Lighed og
Forskel), idet de gælder for alt Bevidsthedsliv, for al Sansning og Tænkning, og
for ugyldig Tænken ligesaavel som for gyldig Tænken. En snævrere Klasse af
Kategorier er de formale (Identitet, Analogi, Negation, Rationalitet), der
lægger sig for Dagen, naar der stilles bestemte Fordringer til Dannelsen af
Begreber, Domme og Slutninger. Naar der spørges om, hvorledes disse formale
Kategorier anvendes paa de kvalitative Forandringer, Iagttagelsen paaviser,
fremtræder de reale Kategorier
% --- p. 19 ---
\opage{19}(Kausalitet, Totalitet og Udvikling). Og naar der særlig spørges om, hvorvidt der
er Betingelser tilstede for en vis Helheds Opstaaen, Bestaaen og Udvikling gennem
Tiden, fremtræder de ideale Kategorier, Værdibegreberne, idet al Værdi forudsætter
et Forhold mellem en Helhed og dens Betingelser.

Kategorierne kan derfor ordnes i en fremskridende Forskelsrække, altsaa ---
naar vi holder os til de vigtigste Kategorier indenfor hver enkelt Gruppe ---
% DISPLAY (p. 19): the series stands centred on a line of its own.
\begin{center}
Syntese $>$ Identitet $>$ Rationalitet $>$ Kausalitet $>$ Totalitet $>$ Værdi.
\end{center}

De følgende Led i Rækken forudsætter stedse de forudgaaende, indeholdes paa
en Maade i dem, er specielle Udformninger eller Anvendelser af dem. Historisk er
Kategorierne fremtraadte sporadisk, alt efter den Trang, der var til at anvende dem,
de Opgaver, det gjaldt at løse. De forskellige Videnskaber, der hver for sig arbejder
med sine Kategorier, opstaar ogsaa sporadisk, selv om \textsc{Comte} har Ret i, at der i
den Rækkefølge, i hvilken de opnaar en streng, positiv Karakter, kan spores en vis
Orden. Gennemgaaende har der været en Tilbøjelighed til at arbejde med Totaliteter
og Værdier førend med de foregaaende Kategorier. De fundamentale Kategorier
bliver først med Klarhed fremdragne af \textsc{Leibniz} og \textsc{Kant}. Netop fordi de er mest
inderlig knyttede til, eller er et med selve vor Aands egen Virkemaade, er det
forstaaeligt, at de opdages sent. Den logiske Orden er i hvert Tilfælde en anden end
den historiske. Ligesom ved den af \textsc{Comte} foreslaaede Klassifikation af
Videnskaberne, kan vi ordne Kategorierne saaledes, at der skrides frem fra de mest
nødvendige og derfor mest omfattende til de speciellere.

Det skal i en senere Sammenhæng undersøges, om det ikke er muligt paa én
Gang at udvide et Begrebs Indhold og dets Omfang. Den omtalte logiske Regel
om det omvendte Forhold mellem Indhold og Omfang gælder i hvert Tilfælde for
Kategorierne, da disse er umiddelbart uddragne af Tankebevægelsens Forløb i
Bevidsthedslivet og specielt i Videnskaben. Det er umuligt af de fundamentale
Kategorier at udlede de andre mere specielle Kategorigrupper. Syntesens Lov kunde
saaledes gælde, selv om Rationalitet og Kausalitet ingen Anvendelse kunde faa; den
gælder f.~Ex. ogsaa for en kaotisk Forskelsrække. Kausaliteten kunde gælde, selv
om vi ikke kunde faa Anvendelse for Begreber som Totalitet og Værdi. \textsc{Hegel}'s
storartede Forsøg paa en rent deduktiv Udvikling af Kategorilæren ud fra de
simpleste og mest omfattende Grundbegreber strandede paa Nødvendigheden af stadig
at optage nye Elementer og Præmisser fra Erfaringen. Omvendt er det umuligt at
gøre de mest konkrete Kategorier, Værdibegreberne, til de mest universelle, saaledes
som det forsøges i alle teologiske eller teologiserende Anskuelser. Begge Dele fører
kun til et Spil med Analogier, uden Mulighed for at paavise baade disse Analogiers
Berettigelse og deres Begrænsning.

Men det følger allerede af Syntesebegrebets Stilling som første Led i
Kategorirækken, at der er en Type, som er fælles for alle Kategorier og Kategorigrupper.
Dersom Syntesen er den fælles Form for al Tankevirksomhed, paa hvilket Omraade
den saa øves, maa alle Tanker ogsaa frembyde en mere eller mindre ud%
% SIGNATURE MARK "35*" at the foot of p. 19 omitted as printer's apparatus.
% --- p. 20 ---
\setcounter{footnote}{0}%
\opage{20}præget Helhedskarakter. Og dersom det er rigtigt, at alle grundlæggende
Definitioner maa være analytiske, saa ligger deri, at vi stedse begynde med Helheder, og
at Tankearbejdet bestaar i at finde Dele eller Elementer indenfor saadanne
Helheder. Saaledes er Identitet et specielt Tilfælde af Lighed, Fladen en Del af
Klodsen o.~s.~v. Afset fra Totalitetsbegrebets specielle Plads i Kategoriernes Række
vil det være at formode, at det foregribes i de forudgaaende Led af Rækken og i
speciellere Form ligger til Grund i de følgende Led. Det er fra denne Side, at
den foreliggende Afhandling genoptager Undersøgelsen af Kategorierne. Den
tidligere, i „Den menneskelige Tanke“ fremstillede Ordning af dem vil herved faa en
nærmere Belysning.

\medskip
8. Historisk gør, som allerede antydet, Totalitetsbegrebet sig meget tidlig
gældende paa Grund af Tendensen til at begynde med de mest konkrete Begreber,
om hvilke man ikke straks kan se, at de tillige er de mest komplicerede, eller
maaske snarere fordi de forskellige Kategorigrupper endnu ikke havde udsondret
sig fra hverandre, men dannede et Hele eller betragtedes som identiske indbyrdes.

Det antike Verdensbillede stod for Tanken som en afsluttet Helhed, indenfor
hvilken hvert Væsen, hver Kvalitet og hver Værdi havde sin bestemte Plads.
Allerede i Førsokratikernes Spørgsmaal: „Hvad er det, som er til?“ kundgør denne
Forudsætning sig. Der søges en Egenskab ved Tilværelsen, som gælder for
Helheden og alle dens Dele, i lige Grad in toto et in parte. Men herved opstod
tillige et Problem, der strakte sig gennem hele Tænkningens Historie. Thi naar hin
Grundegenskab skal være den samme overalt, hvorfra skriver sig da Forskellen
mellem Væsener, mellem Kvaliteter og mellem Værdier? Forskellighederne kan
ikke udledes af den rene Identitet, og der opstaar derfor hos strenge Logikere, som
\textsc{Parmenides}, en Tilbøjelighed til at nægte Forskellighedernes Realitet eller henvise
dem til „de Dødeliges Meninger“. Der fremtræder her en Modsætning mellem Enhed
og Helhed, mellem Rationalitet og Totalitet, en Modsætning, der i nyere Tid
fremtræder som en Modsætning mellem formal og real Videnskab. I \textsc{Aristoteles}' Lære,
at Videnskab kun har med det Almene at gøre, men at alt Virkeligt er individuelt,
altsaa har Totalitetens Præg, fremtræder det Problem, vi her staar ved, paa
karakteristisk Maade --- og for alle Tider.

I Renaissancen trænger Helhedsbegrebet sig frem i første Linie, og der gøres
Forsøg paa at overvinde den antike Modsætning mellem Tanke og Totalitet ved at
fortolke \textsc{Platon}'s „Ideer“ som Love for Fænomenernes reale Sammenhæng. De
sande Ideer, siger \textsc{Bruno}, er ikke Begreber om Ting eller Egenskaber i deres
abstrakte Væsen, men Begreber om Love, som udtrykker forskellige Deles
Sammenhæng i en Helhed, f.~Ex. ikke Begreber om Hoved, Lunge, Hjerte o.~s.~v. hver for
sig, men om Sammenhæng mellem alle disse Ting, hvorved de kommer til at
udgøre en Helhed\footnote{\textit{„Den nyere Filosofis Historie“}, I, p. 126
(\textsc{Bruno}). Jfr. p. 194---196 (\textsc{Bacon}).}. Den saaledes antydede
Tankegang fortsættes hos \textsc{Spinoza} i hans
Lære om individuelle Helheder, der kan bestaa, selv om Delene skifte og de
% --- p. 21 ---
\setcounter{footnote}{0}%
\opage{21}enkelte Processer indenfor Helheden ændrer Retning, naar blot Loven for
Sammenhængen vedbliver at være den samme, --- og om Individualiteter af forskellige Grader,
der danner en Trinrække, indtil tilsidst hele Naturen staar for os som en stor
Individualitet, hvis enkelte Dele varierer uden at det hele store Individ undergaar
nogen Forandring.

Den Loverkendelse, den moderne Videnskab tilstræber, angaar Forholdet mellem
de Fænomener, Erfaringen frembyder. Man vil bort fra de ufrugtbare Almenbegreber,
til hvilke man kan stige op fra de enkelte Iagttagelser, men fra hvilke man ikke
igen kan stige ned til de Enkeltheder, Iagttagelsen godtgør. Dette er den
betydningsfulde Forandring af Opgaven, som den moderne Naturvidenskab skylder sin
Tilblivelse, og som fandt sit Udtryk i \textsc{Galilei}'s Metodelære ved Fordringen om
Samvirken og Supplering mellem Deduktion og Induktion. Men ad denne Vej kan
man ikke komme til Tanken om en saadan afsluttet Helhed, som Oldtiden
fastholdt, --- kun til en stor, uoverskuelig Sammenhæng, der stedse kan udvides.
Allerede \textsc{Bruno} antyder, at hverken Begrebet Del eller Begrebet Helhed strengt taget
kan finde Anvendelse paa Tilværelsen\footnote{L'universo è tutto in tutto (se pur
in modo alchuno se può dir totalità dove non è parte ne fine). \textit{Opere
italiane}, ed Lagarde, p. 315.}. \textsc{Spinoza} forsøger, som vi i en senere
Sammenhæng (Kap. V) vil komme tilbage til, at forene Loverkendelse og
Helhedsopfattelse i ét Begreb. Men \textsc{Hobbes} erklærede bestemt, at om Verden som Helhed kan
man kun stille faa Spørgsmaal --- og ingen af disse Spørgsmaal kan
besvares\footnote{De mundo, ut uno multarum partium aggregato, qvæ qværi
possunt, paucissima sunt, quæ determinari, nulla. \textit{De Corpore}, Cap. 26,
1.}. Den
Tankegang, \textsc{Hobbes} antyder, fuldendes i \textsc{Kant}'s Lære, at Begreber om absolut
Totalitet, „Ideer“, som han kalder dem, ikke kan tjene til Grundlag for
videnskabelig Erkendelse. Og til „Ideerne“ regner han bl.~a. Begrebet Verden; hvis
dette Begreb skal kunne dannes, vil det forudsætte en uendelig og dog fuldendt Syntese.

Allerede de fundamentale Kategorier stiller den Fordring til alt, hvad der skal
kunne erkendes, at det maa kunne sammenfattes med noget andet og stilles i
Forhold til dette som kontinuerligt eller diskontinuerligt med det, og som ligt med eller
forskelligt fra det. Det er derfor, \textsc{Hobbes} har Ret i, at man ikke kan faa Svar
paa, hvad man spørger om angaaende Verden som Helhed. Men des flere
Spørgsmaal bliver der til Gengæld at stille, med Udsigt til Svar, angaaende de begrænsede
Individualiteter eller Totaliteter, som Erfaringen viser os. Da enhver begrænset
Helhed staar i Forhold til andre Helheder eller til Elementer, der er forskellige
fra den selv, bliver det muligt her at anvende vore Kategorier.

En begrænset Helhed kan ikke have Grunden til sin Opstaaen, Bestaaen og
Forgaaen i sig selv. Den er Led i en Række, forudsætter foregaaende Led i denne
Række, som har forberedt den, og fra hvilke den successivt udskiller sig; den maa
i hvert Øjeblik blive til paany, da dens Bestaaen afhænger af bestemte
Betingelser; og forsaavidt den er Aarsag til Forandringer udenfor den, vil den stedse
vise sig at virke sammen med en Mangfoldighed af andre Aarsager, uden hvilke
% --- p. 22 ---
\opage{22}dens ejendommelige Indgriben ikke vilde være mulig eller føre til noget
Resultat.

Men et Problem opstaar ikke blot ved en Helheds Forhold til, hvad der er
forskelligt fra den; ogsaa i Forholdet mellem Helheden som saadan og de
Elementer, der indeholdes i den, ligger der et Problem, som maaske kunde friste til
overhovedet at nægte Helhedsbegrebet erkendelsesteoretisk Berettigelse. Der staar i hvert
Tilfælde her forskellige Retninger overfor hinanden, som svarer til Modsætningen
mellem Deduktion og Induktion og betinges af Vanskeligheden ved her at anvende
den af \textsc{Galilei} krævede Vekselvirkning mellem de to Metoder.

Fra den ene Side lægges hele Vægten paa de Elementer, som kan skelnes
indenfor en given Helhed, som er af mere eller mindre indgribende Betydning for
dens Existens, og som enten i samme eller i analoge Former kan genfindes
udenfor den. --- Fra den anden Side lægger man hele Vægten paa den
Ejendommelighed, hvormed Helheden fremtræder i Erfaringen, paa den indre Sammenhæng
mellem dens Elementer, saalænge de er Led i den, paa den samlede Kraft, den
formaar at anvende i aldeles bestemte Retninger, og uden hvilken den ikke kunde
bestaa. Formedelst denne indre Sammenhæng og denne koncentrerede Optræden
frembyder den Egenskaber, som intet af dens Elementer hvert for sig har, og den
træder derved frem som forskelligartet fra disse. I ethvert Forsøg paa at skrive et
Bjergs, en Flods, en Skovs eller en Klodes Historie træder denne Forskel frem.
Men det er især paa det biologiske, det psykologiske og det sociologiske Omraade,
at der her gør sig en Modsætning og dermed et Problem gældende. Det er
Modsætningen mellem de Retninger, der kan betegnes respektive som Mekanicisme og
Vitalisme, som Associationisme og Spiritualisme, som Individualisme og Socialisme.

Man har villet finde Forskellen mellem Naturvidenskab og Historievidenskab
(eller, som nogle hellere vil sige, Kulturvidenskab) deri, at hin kun gaar ud paa
at finde almene Love for Elementers Samvirken, medens denne har med
individuelle Helheder i deres Ejendommelighed at gøre. Ogsaa overfor en saadan
Opfattelse vil en nærmere Undersøgelse af Totalitetsbegrebet være af Betydning. Dersom
dette Begreb, ligesom alle andre Kategorier, bærer den menneskelige Tankes
ejendommelige Præg, vil det ikke helt kunne fornægtes paa nogetsomhelst
Erkendelsesomraade.

I nogle berømte Linier har \textsc{Goethe} anslaaet en Forskers Stemning ved
Helhedsbegrebet:
% VERSE (p. 22): four lines in smaller type, indented, line breaks as printed.
\begin{verse}
Freuet euch des wahren Scheins,\\
freuet euch des ernsten Spieles!\\
Kein Lebendiges ist eins,\\
immer ist es vieles.
\end{verse}

Helheden har Skinnets og Legens Karakter, naar den tages som fuldt færdig
paa nogetsomhelst Trin, og naar man paa mystisk Maade tillægger den en Kraft
til Selvopholdelse og Udvikling, ganske afset fra Vekselvirkningen mellem indre
Elementer og med ydre Forhold. Men der er Sandhed og Alvor i Skinnet og i
% --- p. 23 ---
\opage{23}Legen, fordi Helhedssammenhængen betinger Elementernes Karakter, ikke mindre
end de bestemmer Helhedens Karakter, og fordi ydre Betingelser kun faar
% DEFECT (p. 23): a stray accent-like mark stands over the first o of
% "Ejendommelighed" (300 dpi); a mark in the scanned copy, not transcribed.
Indflydelse under Medvirken af den Ejendommelighed, den faktisk givne Helhed
fremtræder med. Derhos er selve Tankevirksomheden, paa hvilket Omraade den saa
bevæger sig, Udtryk for en Helhed, der søger at udvide sig under stedse
% DEFECT (p. 23): a stray tick under "forskellige"; not transcribed.
inderligere Sammenslutten af sine Elementer. --- Fra disse forskellige Synspunkter vil
vi i det følgende betragte Totalitetsbegrebet.

% CHAPTER HEAD as printed on p. 23: a short centred rule closes ch. I; then
% two centred lines "II." / title with final stop. First division indented.
\section*{II.\ Totalitetsbegrebet og de fundamentale Kategorier.}
\addcontentsline{toc}{section}{II.\ Totalitetsbegrebet og de fundamentale Kategorier}

9. Uden uvilkaarlig Helhedsdannelse vilde Eftertanken, der begynder med
Analyse, ikke være mulig. En konkret Intuition, Frugt af en Syntese, der er
foregaaet uden at være Genstand for Bevidsthed, danner Grundlaget for al Erkendelse.

Hverken et absolut Enkelt og Enestaaende, eller et absolut Kaos vilde kunne
danne Grundlag for Eftertanke. Det Faktum, at analyserende Eftertanke er mulig,
--- dette Faktum, der er Grundlag for al Videnskab, --- godtgør, at den primære
psykiske Proces bestaar i en Helhedsdannelse.

Eftertanken retter sig dels mod Forholdet mellem de skiftende Helheder (paa
Sansningens, Erindringens og Fantasiens Omraader), dels mod Forholdet mellem
de forskellige Dele eller Egenskaber ved den enkelte Helhed. Ethvert Emne er en
lille Verden, men denne lille Verden viser sig ogsaa at være Del af en større
Verden eller at staa i Forhold til andre smaa Verdener. Ethvert Emne har, som
man har sagt, Rifter eller Frynser i Randen, der røber, at det er taget ud af en
nøje Forbindelse med andre Emner. Skridt for Skridt opdages det, at enhver
Bestemmelse, f.~Ex. enhver Tids- eller Stedsbestemmelse, enhver Konstatering af en
Egenskab, i Grunden forudsætter en langt større Helhed end den, Eftertanken var
begyndt med.

Allerede den enkelte Sansefornemmelse bestemmes, baade hvad Styrkegrad og
Kvalitet angaar, ved den hele Sammenhæng, indenfor hvilken den optræder. Ved
Genkendelse gør sig tillige en Sammenhæng med tidligere Fornemmelsers
Eftervirkninger gældende. Der foregaar her Processer, som Sproget ikke har særegne Navne
for, men som frembyder Analogier til, hvad der paa mere udviklede
Bevidsthedstrin fremtræder i Forholdet mellem givne Emner, deres Tydning og de Slutninger,
der drages fra dem, --- en Analogi, man ikke maa forveksle med Identitet. I Erin%
% --- p. 24 ---
\setcounter{footnote}{0}%
\opage{24}dring og Fantasi dannes Helhedsbilleder i Kraft af den til Grund for al
Forestillingsassociation liggende Helhedslov, som vidner om en Trang og Evne til at
% DEFECT (p. 24): a stray raised mark before "allerede" (layer reads
% "‘allerede"); not transcribed.
udfylde og fortsætte givne Forestillinger. Som allerede omtalt (§ 3), fastholdes
konkret Intuition saalænge som muligt, og der kan foregaa Forskydninger og
Artikulationer indenfor den, uden at der skrides til Dannelse af Domme. Foruden de
Exempler, jeg i tidligere Fremstillinger har anført, skal jeg her anføre et Exempel
paa Artikulation indenfor et Erindringsbillede. Mit Erindringsbillede om en
Afstemning ved Uddeling af et Legat indeholdt bl.~a. Billedet af et Bord, paa hvilket
der henlaa Stemmesedler. Men da Erindringsbilledet ved en senere Lejlighed
dukkede op igen, var det ikke Stemmesedler, der laa paa Bordet, men Breve.
Jeg kom da naturlig til at sammenligne de to Erindringsbilleder med, hvad jeg
ad anden Vej vidste, og fældede saa den Dom, at det sidste havde Ret: thi
Stemmesedlerne blev indsendte i lukkede Konvoluter. Erindringsbilledet var
altsaa uvilkaarlig undergaaet en Forskydning i rigtig Retning, og det var først
Forholdet, Modsigelsen mellem de to Erindringsbilleder, der nødvendiggjorde
Dommen.

Et allerede berørt Exempel paa Artikulation fremkommer ved Tilbøjeligheden
til at gøre en Flades Figur og dens Omrids til ét. Thi Figuren kan ikke
lokaliseres; den vedkommer hele Fladen. Fladefiguren er det umiddelbart Oplevede; men
Omridset er baade lettere at beskrive og har mere praktisk Betydning. At
Omridset trænger sig frem i Forhold til Fladefiguren, er en Antydning af en Analyse,
der senere kan finde Sted, men som ikke er absolut nødvendig, da Forskydningen
kan foregaa indenfor uvilkaarlig Anskuen. Som \textsc{Edgar Rubin}, der klart har belyst
dette Punkt, bemærker, staar vi her ved et Udslag af en Ejendommelighed ved det
Psykiske, der ogsaa gør sig gældende paa andre Omraader. Ligesaalidt som
Fladefiguren kan lokaliseres nogetsteds i Rummet indenfor den givne Helhed, fordi den
hører selve Helheden til, ligesaalidt kan Oplevelser af Rytme, Bevægelse og
Varighed henføres til noget enkelt Tidspunkt indenfor den hele
Oplevelse\footnote{\textit{Synsoplevede Figurer}, p. 179.}. Den begyndende
Analyse støder, som allerede Eleaterne saa, paa store Vanskeligheder
her. Selve Helheden staar i Fare, naar Analysen begynder, idet den er
inkommensurabel med enhver af de Enkeltheder, Analysen kan paapege. Tilsyneladende
usammensatte Sansefornemmelser er smaa Helheder, der ikke kommer til deres Ret ved
Analyse. Saaledes er Fornemmelsen af Ruhed, en Helhed af Elementarfornemmelser
af smaa Ujævnheder, der ikke mærkes hver for sig; ingen enkelt Fornemmelse af
en saadan Ujævnhed vilde give Helhedsfornemmelsen. --- Paa Grund af denne
Ejendommelighed ved Bevidstheden forstaas Trangen til at fastholde Anskuelsen
saa længe som muligt.

En særegen Art Forskydning har vi Exempel paa i de Lys- og Farvebilleder,
Tonerne kan fremkalde, hvad man har kaldt musikalske Fotismer, hvor den ved
Toner, særlig ved Klangfarven, fremkaldte Stemning umiddelbart faar Udtryk i
Figur- eller Farvevisioner. Det er en „instinktiv Oversættelse til visuelle Udtryk“
% --- p. 25 ---
\setcounter{footnote}{0}%
\opage{25}af hvad der har paavirket Høresansen\footnote{\textsc{Th. Flournoy}:
\textit{Des phénomènes de Synopsie}, p. 103.}. Naar herved Følelseselementer i Regelen
gør sig gældende, saa er det et Exempel paa, at konkret og praktisk Intuition gaar
over i hinanden, idet der ved de uvilkaarlige Helhedsdannelser paa Sansningens,
Erindringens og Fantasiens Omraader stedse, i større eller mindre Grad, vil
medvirke Trang og Følelse.

I den uvilkaarlige Trang til Helhedsdannelse har Forestillingen om Selvet
sin sidste Grund. Den skyldes et Helhedsbillede, der dannes af de mest konstante
% DEFECT (p. 25): a stray tick over the m of "Stemninger"; not transcribed.
Fornemmelser, Erindringer, Stemninger og Bestræbelser, et Helhedsbillede, der
stedse vil være artikuleret i en bestemt Retning efter den herskende Interesse. Især
saadanne Tilstande, der staar for os som Aarsager til Forandringer i os eller
udenfor os, øver her deres Indflydelse. Der virker her, hvad \textsc{Gadelius} har kaldt en
motorisk Helhedsfølelse\footnote{\textsc{Bror Gadelius}: \textit{Tro och öfvertro}.
I, p. 158.}. Til udtrykkelig Bevidsthed om os selv (eller om vort
Selv) hører Analyse. Og Analysen har her, som overalt, sine Vanskeligheder. Vi
holder os ogsaa her saa længe som muligt til, hvad umiddelbar og virkelig
Anskuen yder. Deraf kommer det Mystiske, der ofte gør sig gældende i Forestillingen
om os selv. Her som overalt er Mystik berettiget, naar den er Udtryk for det
Uudtømmelige i det Givne. Vor Forestilling om os selv er for det meste en
praktisk Intuition. Analysen kan kun drage et enkelt Træk frem ad Gangen; og naar
er alle Træk fremdragne, og hvilken Garanti kan vi have for, at de alle er fundne?
Vor Selverkendelse er derfor stedse ufuldkommen, ganske afset fra, at der,
saalænge Livet varer, kan komme nye Oplevelser og dermed muligvis fundamentale
Ændringer i den Totalsammenhæng, ud af hvilken Forestillingen om Selvet skal
fremgaa ved Analyse. Ogsaa her forstaas det, hvorfor vi klamrer os til Anskuelsen
saa længe som muligt. Vi beraaber os især paa vor Samvittighed, naar vi erfarer,
eller mener at erfare Analysens Utilstrækkelighed.

Naar den konkrete og den praktiske Intuition ikke længere formaar at rumme
Emnerne, enten paa Grund af disses Mangfoldighed eller paa Grund af deres
indbyrdes Modsætninger, er der Trang til, og maaske ogsaa Mulighed for en ny Art
Helhedsdannelse, nemlig Domsdannelsen. Idet jeg herom henviser til mine tidligere
Fremstillinger („Det psykologiske Grundlag for logiske Domme“ og „Den
menneskelige Tanke“), skal jeg her blot gøre nogle Bemærkninger om en Teori, der er
fremsat af \textsc{Heinrich Maier} i hans allerede ovenfor omtalte Værk „Psychologie des
emotionalen Denkens“. Overfor saare meget i dette Værk kan jeg udtale min
fuldkomne Samstemmen. Forskellen mellem \textsc{Maier}'s og min Fremstilling beror ofte
kun paa en forskellig Terminologi. Men i to Punkter kan jeg ikke være enig med
ham. Paa den ene Side udvider han, som allerede ovenfor omtalt, Begrebet Dom
til at gælde for alle Iagttagelser, Erindringer og Fantasier, fordi de ikke er
„ligegyldige“ overfor Modsætningen mellem Sandt og Usandt. Paa den anden Side
hævder han, at „den emotionale Tænken“ ikke fører til Erkendelse, men kun til
„Præsentation“ af, hvad der er Genstand for vor Følelse og vor Trang.
% SIGNATURE LINE at the foot of p. 25 ("D. K. D. Vidensk. Selsk. Skr., 7.
% Række, hist. og filos. Afd. III, 2." and "36") omitted as printer's apparatus.

% --- p. 26 ---
\setcounter{footnote}{0}%
\opage{26}Det første Punkt har jeg allerede berørt ovenfor (§ 3), hvor jeg har søgt at vise,
at \textsc{Maier} med Urette lægger Modsætningen mellem Sandt og Usandt over i den
primære Anskuen. Den spiller ingen Rolle der, fordi der raader en tryg Hvilen i,
hvad der sanses eller forestilles, en Tryghed, der ikke kender Tvivl og derfor ikke
skelner eller kan skelne mellem Sandt og Usandt. Adam og Eva var ikke
ligegyldige mod Forskellen mellem Godt og Ondt, før de spiste af Kundskabens Træ; de
kendte slet ikke denne Forskel. Og hin umiddelbare Tryghed kan ikke siges at
være ligegyldig mod en Forskel, der ikke existerer for den.

Jeg vender her tilbage til dette Punkt for at gøre opmærksom paa, at \textsc{Maier}
har to Kendemærker, ved hvilke han afgør, om en Dømmen gør sig gældende eller
ikke. Det ene er den allerede omtalte Modsætning mellem Sandt og Usandt. Det
andet bestaar deri, at der foretages en Overgang fra et selvstændigt Led indenfor
et Objekt til et andet selvstændigt Led. Han anerkender, at der kan gøre sig en
Forholdssætten (Beziehungstätigkeit) gældende, som endnu ikke er en Dømmen,
fordi Leddene ikke staar som selvstændige overfor hinanden\footnote{\textit{Psychologie
des emotionalen Denkens}, p. 123.}. Dette svarer til,
hvad jeg kalder artikuleret eller artikulerende Anskuen, et Fænomen, som jeg ikke
mener, at \textsc{Maier} lader komme til sin fulde Ret.

Det er efter min Opfattelse kun dette andet Kendemærke, Overgangen fra et
selvstændigt Led til et andet selvstændigt Led, der betegner Dommen i Modsætning
% DEFECT (p. 26): a stray low dot before "umiddelbare" (layer reads
% ".umiddelbare"); not transcribed.
til Anskuelsen. Og det forudsætter Analyse; ti i den umiddelbare Anskuen staar
Leddene ikke som selvstændige. De er endnu ikke udskilte af den uvilkaarlig
dannede Helhed. En saadan Analyse og Overgang kan iøvrigt meget godt foregaa,
uden at Forskellen mellem Sandt og Usandt gør sig gældende. En Dom kan
fældes ligesaa uvilkaarlig, som en Anskuelse kan dannes. Der kan raade umiddelbar
Tryghed ved vore Domme ligesaa vel som ved vore konkrete og praktiske
Intuitioner. Det er først, naar Tvivlen --- altsaa en vis speciel Art af Domme --- opstaar,
at hin Forskel kan komme frem.

\textsc{Maier}'s to Kendemærker kommer i Strid med hinanden, idet der er Domme,
i hvilke Modsætningen mellem Sandt og Usandt ikke gør sig gældende, og idet det
dog er denne Modsætnings Tilstedeværelse, der skal gøre det nødvendigt at
udstrække Dommens Begreb selv til Bevidsthedsvirksomheder, der ikke bestaar i en
Overgang fra et selvstændigt Led til et andet.

Hvad det andet Punkt angaar, kan jeg ikke se andet, end at \textsc{Maier}, trods den
Fortjeneste, han har indlagt sig ved Karakteristiken af den følelses- og villiesbestemte
Tænken, ikke har ladet „den emotionale Tænken“ komme til sin fulde Ret, som en
Funktion, der i Grunden er af ganske samme Art, som „den kognitive Tænken“,
skønt den arbejder under særegne psykiske Betingelser. Man maatte jo dog ogsaa
vente, at Tankevirksomheden væsentlig virker paa samme Maade, hvad enten den
foregaar overvejende paa Grundlag af Fornemmelser, Iagttagelser og Forestillinger,
eller den foregaar overvejende paa Grundlag af Trang, Lyst og Ulyst.

\textsc{Maier} bruger her det første af de ovenomtalte to Kendemærker, han har paa
% --- p. 27 ---
% STRAY MARKS on p. 27 (read by the layer as a right guillemet, "►", "s", "k" at the page
% head): marks in the scanned copy (ticks along the left margin); the head
% itself carries only the folios 27 / 283. Not text of the work.
\setcounter{footnote}{0}%
\opage{27}en Dom. Den emotionelle Tænken, siger han (p. 140), vil ikke være sand, og den
er derfor ingen Dom. Men er det nu rigtigt, at emotionel Tænken ikke vil være
sand? Ifølge \textsc{Maier} selv bestaar dens Betydning deri, at den bringer Genstanden
for vor Lyst og Ulyst eller for vor Trang klart og tydelig frem for os. Derfor taler
\textsc{Maier} ogsaa om Evidens i volitive Tankeakter, en Evidens, der bestaar i, at
Begæringstendensen kommer til adækvat Udtryk i Begæringsforestillinger. Og det er især,
naar en Tvivl først har gjort sig gældende, at denne Evidens tydelig fremtræder
(p. 42). Jeg kan ikke se andet, end at Begrebet Sandhed stikker Hovedet frem i
saadanne Udtryk som „adækvat“ og „Evidens“. Og Forholdet er her i Virkeligheden
det samme som ved „kognitive“ Funktioner, der jo ogsaa tydeligst gør sig
gældende efter en Tvivl. Hvad jeg sanser, erindrer eller forestiller mig, kan jo
blive mig mere eller mindre „adækvat“ eller „evident“ bevidst, d.~v.~s. min
Forestilling derom kan mere eller mindre fuldstændig dække Sansningens, Erindringens
eller Fantasiens Indhold, være Frugt af en mere eller mindre fuldstændig Analyse
af de umiddelbare Anskuelser. Og saaledes kan der ogsaa ved den Forestilling,
der angaar Maalet for en Trang eller en Drift, spørges, om den nu ogsaa er sand,
d.~v.~s. fuldstændig udtrykker Trangens eller Driftens virkelige Retning. Endnu
tydeligere er det ved en Beslutning: Er nu det, jeg har foresat mig, ogsaa det, jeg
inderst inde vil? Er der Sandhed i min Villen, d.~v.~s. dækker den min virkelige
Stræben? Er det ikke et illusorisk Maal, jeg har stillet mig, d.~v.~s. et Maal, jeg
i Grunden ikke tragter efter? Der er ved Bevidstheden om en Villen ligesaa vel
Illusioner mulige, som ved Sanseiagttagelser. Og der er en Analogi mellem Forsæt
og Iagttagelse: i Afgørelsens Øjeblik ser jeg mig selv handlende i et fremtidigt
Øjeblik, anteciperer (perciperer) en fremtidig Situation. Og det maa føjes til, at
selv om Forsættet udføres, er det ikke sagt, at det var Udtryk for, hvad der i
Virkeligheden var den inderste Stræben.

Der gives, som især danske Filosofer (\textsc{Poul Møller}, \textsc{Kierkegaard}) har
indskærpet, en personlig Sandhed, som bestaar deri, at de Maal, vi med Bevidsthed
sætter for vor Stræben, ogsaa stemmer med vor Personligheds inderste Væsen.

Istedetfor at skelne mellem kognitiv og emotionel Tænken bør der derfor
skelnes mellem to Arter af Emner, saadanne, der fremtræder som blotte, væsentlig
uinteresserede Iagttagelser, og saadanne, der optræder med Karakteren af Trang
og Stræben. Eftertanken virker paa samme Maade overfor begge Arter. Begge
Steder begynder vi med uvilkaarlig Tryghed, uden at være „ligegyldige for Sandt
og Usandt“, men med Mulighed af at hildes i Illusioner. Og begge Steder kan
der blive Plads for en Dømmen i Betydning af Overgang fra Led til Led, og derefter
maaske for en kritisk Prøvelse, der ophæver vor primære Tryghed. De
Forestillinger, der er indvævede --- snart mere som Betingelser, snart mere som
Produkter --- i vor Trang og Stræben\footnote{\textsc{Maier} tillægger Forskellen
mellem uvilkaarlig og vilkaarlig Villen stor Betydning (558 ff; 574); men han
sætter Grænsen paa et
% DEFECT (p. 27 note): a stray low mark stands between "et" and "andet".
andet Sted end jeg, idet han regner Drift og Ønske til
uvilkaarlig Villen og lader Overvejelse være Forudsætning for vilkaarlig Villen.
Spontane og reflexe Bevægelser
% NOTE continues at the foot of p. 28.
regner han ikke for psykiske Funktioner, og Instinktet regner han til Driften,
idet han paastaar, at der ogsaa i det stedse er en Forestilling om, hvad der
tilstræbes. Jeg tror ikke, at Iagttagelse bekræfter denne Inddeling. Der kan
gøre sig Trang gældende og faa Luft i uvilkaarlig Handlen, uden at nogen
Forestilling om et Maal medvirker.}, vore Drifter og Ønsker, vore Hensigter, Forsæt
% SIGNATURE MARK "36*" at the foot of p. 27 omitted as printer's apparatus.
% --- p. 28 ---
\setcounter{footnote}{0}%
\opage{28}og Beslutninger, kan ligesaavel blive Led i Domme og Genstand for Kritik som
de Forestillinger, der indvæves i vore Iagttagelser, Erindringer og Fantasier. Det
er en Fortjeneste af \textsc{Maier}, at han saa energisk har henledet Opmærksomheden
paa dem. Men han har derved ikke paavist nogen helt ny Art af Tænken. De
fundamentale Kategorier optræder i begge Arter, og naar Kritiken er vaagnet, vil
de andre Kategorigrupper ogsaa kunne gøre sig gældende. I Drøftelsen af
Værdikategorierne (Kap.~V) vil vi komme tilbage til dette Spørgsmaal.

\medskip
10. Eftertanke forudsætter et umiddelbart Givet i Form af Anskuen (konkret
eller praktisk Intuition), og dette umiddelbart Givne staar som en Helhed, der nu
kan blive Genstand for Analyse. I en saadan Helhed viser en Mangfoldighed af
Elementer sig forbundne paa ejendommelig Maade, og vi finder derfor allerede
her Syntesens Lov gældende.

Mod Berettigelsen af at finde Syntese i det umiddelbart Givne har \textsc{Hermann
Cohen} gjort den Indvending, at Syntese forudsætter en given Mangfoldighed, der
kan sammenfattes, og han har ment, at \textsc{Kant} ved at anvende Syntesebegrebet gør
sig skyldig i \textsc{Locke}'s og \textsc{Hume}'s Fejl, nemlig at betragte det oprindelig
Givne som et Kaos af Sansefornemmelser\footnote{\textsc{H. Cohen}: \textit{Logik der
reinen Erkenntniss.} (1902). p. 24.}.
% PRINTED AS IS (p. 28): "nogen anden Form"; the sense wants "noget".
Under en nogen anden Form fremtræder den
samme Indvending, naar \textsc{Bergson}\footnote{\textsc{Bergson}: \textit{Introduction
à la Métaphysique.} (Revue de Métaphysique et de Morale. 1903), p. 15---17.} hævder,
at det umiddelbart Givne ikke skal
% PRINTED AS IS (p. 28): "Modsætntng" for "Modsætning" (checked at 600 dpi;
% the text layer silently reads "Modsætning").
kunne udtrykkes ved et Begreb, der indeholder en Modsætntng til et andet Begreb,
og Begrebet Syntese indeholder jo en Modsætning til Begrebet Mangfoldighed.

Overfor denne Indvending maa det indrømmes, at Syntesebegrebet er overført
ved Analogi fra mere differentierede psykiske Funktioner til de mere primære.
Sproget har, som allerede bemærket, ikke dannet Udtryk for de mest primitive
psykiske Funktioner. Syntese fremtræder tydeligst, hvor den Opgave stilles at
forbinde og ordne faktiske og tydelig foreliggende Forskelligheder. Ved bevidst
stillede Opgaver har vi baade den forudgaaende Mangfoldighed og den ved psykisk
Arbejde vundne Sammenhæng for os som Genstand for Bevidsthed, og maaske
forskellige Mellemstadier af fremskridende Ordning og Helhed. Ad Analogiens Vej
anvender vi Begrebet Syntese, og med det Begrebet psykisk Arbejde, selv der, hvor
der ikke kan iagttages en forudgaaende Mangfoldighed og ikke heller opfattes
Mellemstadier i dens Sammenfatten, saaledes ved umiddelbar Anskuen, altsaa ved
Sansebilleders, Erindringsbilleders og Fantasibilleders Opdukken med det
Helhedspræg, der er Betingelse for Eftertankens Analyse. Det er baade Kontinuitet og
Analogi, der her ligger til Grund. Ti der kan være alle mulige Grader i den
Hastighed, hvormed det sammenfattende Arbejde øves, og Erfaringen viser, at jo
kortere den fysiologiske Tid er, des mere afløses Bevidsthed af Ubevidsthed. Og
% --- p. 29 ---
\setcounter{footnote}{0}%
\opage{29}alt tyder paa, at selv det, der psykologisk fremtræder som enkelt og afsluttet,
fysiologisk set repræsenteres ved en Mangfoldighed af Processer. Vi har derfor ikke
Ret til at mene, at psykisk Arbejde (til hvilket der naturligvis antages stedse at
svare et fysiologisk Arbejde) nogensinde helt falder bort, saalænge Bevidstheden
bestaar. Der er alle mulige Grader mellem potentiel og aktuel psykisk
Energi\footnote{Allerede i \textit{Den menneskelige Tanke} p. 10---19 har jeg drøftet
dette Spørgsmaal.}.
% STRAY MARK before the note numeral at the foot of p. 29, a mark in the copy.

For umiddelbar Anskuen (konkret og praktisk Intuition) er der ingen Problemer
til. De kommer først ved Eftertanke. Hvis det drejer sig om at komme
i en problemløs Tilstand, gælder det blot om at holde Eftertanken ude. Det har
vi alle en Trang til, og det er ogsaa berettiget, naar vi helt og fuldt skal gaa op i
en Oplevelse. Fantasiens og Kunstens Værdi forudsætter en saadan Opgaaen, ligesaa
forholder det sig med Livsforhold, vi gaar op i, og med Hengivelse til store
Ideer og Bestræbelser. Men Livet kræver stadig nyt psykisk Arbejde, idet Syntesen
støder paa ny Modstand. Og det vil være overordentlig vanskeligt at afgøre,
om man i de enkelte Tilfælde staar overfor noget aldeles umiddelbart, --- om ikke
Overlevering, Indskydelse (Suggestion) og Eftertanke medvirker. Der er en Tendens
til at betragte tilvante og ved Tradition helligede Tilstande som aldeles
umiddelbart givne. I alle Forsøg paa at skelne mellem Tro og Viden fremtræder en
saadan Tendens i større eller mindre Grad.

Ogsaa \textsc{Hume} henviste til et umiddelbart Givet. Men for ham var det ikke en
Helhed eller et Kontinuum, men en kaotisk Mangfoldighed. Men et Kaos er det
ligesaa vanskeligt at konstatere som et absolut Kontinuum. Tilstedeværelsen af en
kaotisk Forskelsrække kan kun godtgøres, naar indgaaende Forsøg paa at finde
Orden og Sammenhæng er endte med negativt Resultat. \textsc{Hume} saa, ligesom
\textsc{Bergson}, at det første Grundlag for Erkendelse ikke selv kan være rationelt.
Absolut Mangfoldighed er ligesaa irrationel som absolut Kontinuum. Baade Diskontinuitet
og Kontinuitet stiller Opgaver --- hin den at finde mere Sammenhæng, denne den
at finde mere Forskellighed. Ethvert Emne frembyder, saavel indenfor det selv,
som overfor andre Emner, baade Kontinuitet og Diskontinuitet. Vil vi tale om
et umiddelbart Givet, maa begge Kategorier anvendes.

Selve Begrebet „umiddelbart Givet“ er dannet af Eftertanken. Hvis der ikke
forelaa Andet end et absolut „umiddelbart Givet“, vilde intetsomhelst Begreb, og
altsaa heller ikke Begrebet „umiddelbart Givet“ kunne dannes. Det dannes netop
i Modsætning til Eftertanken --- og af selve Eftertanken, naar denne bliver sig sin
Virkemaade og sine Forudsætninger bevidst. Ethvert umiddelbart Givet er bestemt
ved dets Modsætning til en vis Grad og Art af Eftertanke. Hvad der i en
Henseende er umiddelbart, behøver ikke at være det i andre Henseender.

Det er Trangen til at bestemme Forholdene, Relationerne mellem forskellige
Emner og mellem samme Emnes Elementer, der fører ud over det umiddelbart
Givne. I et umiddelbart Givet (en Anskuen) er der maaske foreløbig Ligevægt
mellem Sammenfatten (Syntese) og Relation, mellem Kontinuitet og Diskontinuitet,
mellem Lighed og Forskel. Men Ligevægten kan ophæves, og der vil da svinges
% --- p. 30 ---
\setcounter{footnote}{0}%
\opage{30}mellem modsatte Tendenser. Overvægten af Syntese, Kontinuitet og Lighed fører
til at fortsætte i det uvilkaarlig indslaaede Spor, til at handle efter Principet
„Hvorfor ikke?“ Denne Tendens kan føre til at sammenstille, forbinde og
identificere i stedse større Omfang. Men naar Relation, Diskontinuitet og Forskel gør
sig stærkt gældende, søges nærmere Bestemmelser og Grænser, og der handles efter
Principet „Hvorfor?“ Der kræves da Angivelse af de Forhold, i hvilke Emner og
Elementer staar til hverandre, og Paavisning af, hvorledes der kan være
Sammenhæng trods Diskontinuiteten, og hvorledes Ligheds- og Forskelsgraderne forholder
sig til hverandre. Der stilles nu bestemte Problemer, nu, da Anskuelsens Ligevægt
er ophævet.

Men de forskellige Kategorier, der ved Problemdannelse træder i Modsætning
til hverandre, er dog parvis forbundne som Korrelatbegreber, og det ene Led i et
saadant Par kan ikke finde Anvendelse uden det andet. En bestemt Udformning
af det ene Led maa medføre en bestemt Udformning af det andet. \textsc{Wundt} har i
en interessant Afhandling drøftet Korrelatbegrebers Betydning og paavist den
Tendens, der, især i spekulative Systemer, ofte gør sig gældende til at holde sig
udelukkende til et enkelt af Leddene uden at tage Hensyn til det andet, eller maaske
endog saaledes, at dette andets Gyldighed helt nægtedes. \textsc{Hegel}'s Dialektik fremkom
netop derved, at han blev opmærksom paa Grundbegrebernes indbyrdes Sammenhæng,
ifølge hvilken det ene med Nødvendighed maa suppleres med det
andet\footnote{\textsc{Wundt}: \textit{Zur Geschichte und Theorie der abstrakten
Begriffe.} (Kleine Schriften, I. p. 226---258).}.

Et Pseudoproblem har ofte fremstillet sig her ved, at man har lagt den udviklede
Bevidstheds Skelnen mellem Subjekt og Objekt ind i den primitive Bevidsthed. En
saadan Skelnen bliver først mulig ved Tvivl, der igen forudsætter Skuffelse. Man
har endog lagt den skolastiske Distinktion mellem Bevidsthedsakt og
Bevidsthedsindhold ind i den umiddelbare Bevidsthed, og forklarer derpaa Troen
paa Realitet derved, at Indholdet faar Indflydelse paa Akten. Men der kan kun
skelnes mellem Akt og Indhold, efterat det har vist sig, at der kan være „Akter“
uden realt Indhold. Der gør sig fra først af ingen Forskel gældende mellem den
egne Virksomhed og det modtagne Indhold. Denne Forskel opstaar, naar man har
handlet efter, d.~v.~s. draget alle teoretiske og praktiske Konsekvenser af, sit
Bevidsthedsindhold, og denne Handlen saa medfører Skuffelse. Det er Handletrangen,
der baade begrunder den oprindelige Tillid til ethvert givet Bevidsthedsindhold og
senere kan føre til at betvivle og berigtige det.

Af særlig Interesse er det i erkendelsesteoretisk Henseende, at Kriteriet, som
kan gøre Ende paa Tvivlen, ligger i det Spørgsmaal, om Indholdet kan udgøre Del
af saa stor en Sammenhæng, som vi er istand til at tilvejebringe mellem vore
Iagttagelser og vore Erindringer. For at tage et simpelt Exempel. Naar jeg ikke kan
huske, om jeg har lukket min Gadedør, --- d.~v.~s. naar jeg ikke kan genkende de
Berørings- og Bevægelsesfornemmelser, Lukningen medfører, --- kan det hjælpe, om
jeg kan genkalde og genkende alle øvrige Omstændigheder (Grebet efter Nøglen i
Lommen, Indstikken af Nøglen i Laasen, Smækket af Døren, --- maaske ogsaa det,
% --- p. 31 ---
\opage{31}jeg i Øjeblikket tænkte paa, da jeg stod ved Døren o.~s.~v.). Det er Helheden, der
ophæver Tvivlen.

Men ved et Exempel som dette er vi allerede kommet ud over de fundamentale
Kategoriers Gruppe. Det viser sig allerede deri, at en lille og spinkel, men
rationel Traad af logisk Konsekvens kan betyde mere end et stort Komplex af
Enkeltheder. For at forstaa, hvorledes dette er muligt, maa vi undersøge de to
følgende Grupper af Kategorier, de formale og de reale.

% CHAPTER HEAD as printed on p. 31, two centred lines ("III." / title with
% final stop), set below a short rule that closes ch. II.
\section*{III.\ Totalitetsbegrebet og de formale Kategorier.}
\addcontentsline{toc}{section}{III.\ Totalitetsbegrebet og de formale Kategorier}

11. Begreb, Dom og Slutning betegner forskellige Former og Grader af logisk
% PRINTED AS IS (p. 31): "cn Gruppe" for "en Gruppe" (checked at 600 dpi).
Totalitet. Begrebet er en Tankehelhed, i hvilken cn Gruppe af Kendemærker er
forbundne paa en bestemt Maade. Dommen er en Tankehelhed, i hvilken en Gruppe
af Begreber er forbundne efter deres indbyrdes Forhold. Slutningen er en
Tankehelhed, i hvilken flere Domme er forbundne efter den Maade, paa hvilken de
indbyrdes forudsætter hverandre.

I Slutningen fremtræder den formale Totalitet tydeligst. Ethvert Led er paa den
inderligste Maade knyttet til de andre, nemlig gennem Identitet, hvad et enkelt
fælles Begreb angaar. Og i det Forhold, der saaledes viser sig at være mellem
forskellige Domme, fremtræder Rationaliteten. Alle Præmisser tilsammen indeholder
paa Forhaand Slutsætningen, og hver enkelt Præmis faar kun Betydning ved
at forbindes med andre Præmisser. En saadan inderlig Enhed af Tanker, ved
hvilken hver enkelt Tanke bærer de andre og bæres af dem, er Filosofiens store
Ideal. \textsc{Goethe}'s Mefistofeles spotter over den Maade, hvorpaa den formelle Logik
udsondrer og betoner de enkelte Led i denne Helhed. Men denne Spot er kun
berettiget, hvor det „aandige Baand“ mangler. Det, at en Organisme kan
dissekeres, hindrer ikke, at den, saalænge den er i Live, er en Helhed. Det er nu
heller ikke alle Fremstillinger af formel Logik, der lader Tankehelheden i dens
% TRYKFEJL (p. 80): deres læs den
forskellige Former komme lige meget til deres Ret. Sproget lægger her ofte
Vanskeligheder i Vejen, idet allerede det, at hvert enkelt Kendemærke, hvert enkelt
Begreb og hver enkelt Dom faar sit særlige sproglige Udtryk, giver dem en
tilsyneladende Selvstændighed overfor hverandre indbyrdes og en tilsyneladende
Ligegyldighed overfor den Sammenhæng, i hvilken de forekommer. Det vilde dog være
Uret mod Sproget blot at lægge Vægt paa denne Side af Sagen. Thi Sproget viser
% --- p. 32 ---
\setcounter{footnote}{0}%
\opage{32}netop sin nøje Sammenhæng med Tanken derved, at det enkelte Ord og den enkelte
Sætning kun bliver helt forstaaelige ved at henføres til en Helhed, enten den
sproglige Helhed, hele det System af Ord og Sætninger, hvortil de hører, eller ogsaa til
den hele konkrete Situation, hvori de fremkommer\footnote{\textit{Den menneskelige
Tanke}, p. 84. --- \textit{Det logiske Prædikat.} (Vidensk. Selsk. Overs. 1914), p.
242.}. Det skyldes altid en Abstraktion, ofte en utilbørlig Abstraktion, naar det
enkelte Ord eller den enkelte Sætning
rives ud af den sproglige eller reale Sammenhæng, i hvilken de fremkommer.

Det er Videnskabens Opgave at føre alle Emner, der fremtræder, ind i en stor
Sammenhæng, der har sit Forbillede i den Tankehelhed, vore Begreber, Domme og
Slutninger udgør. Den kritiske Filosofi hævder her, i Modsætning til Dogmatismen
i alle dens forskellige Former, at der kun kan opnaas en fremskridende
Gennemførelse af en stor Analogi, uden at vi har Ret til at gøre denne Analogi til
Identitet. Der stiller sig alligevel Opgaver nok indenfor denne Synsmaade. Thi
der kan gives flere forskellige Tankehelheder, hver gældende for sit
Erfaringsomraade, og Spørgsmaalet bliver, om disse forskellige Tankehelheder kan
sammenarbejdes til at udgøre en eneste Tankehelhed. Og indenfor enhver enkelt
Tankehelhed opstaar Spørgsmaalet, om den udtømmer det Erfaringsomraade, hvortil den
svarer, --- om Begrebet indeholder alt, hvad der „ligger i“ Anskuelsen, --- om
Dommens Prædikat udtømmer Subjektets Indhold, --- om Slutsætningen giver Alt, hvad
der indeholdes i Præmisserne.

Betingelse for en Begrebsdannelse er, at Emnerne er ordnede paa en bestemt
Maade, saa at deres indbyrdes Forhold kan fremtræde. Skal Begrebet Hest dannes,
maa vi have de forskellige Racer og Variationer, nulevende og uddøde Former,
ordnede i en Række, og Begrebet udtrykker da Forholdet mellem disse forskellige
Former, deres Ligheder og deres Forskelligheder. Domsdannelsen forudsætter igen
en Ordning af Begreber, særlig i en partiel Identitetsrække, og Slutningen
forudsætter en Ordning af Begreber i en af de rationelle Rækker.

Alle saadanne Rækkedannelser gaar ud paa at ordne Indholdet af en foreløbig
dannet Emnegruppe (en Emnetotalitet), saaledes at en nøjere Analyse kan
finde Sted. I ethvert enkelt Tilfælde vil der være givet mere, end den stillede
Opgave kræver. Hele Rækken $A < B < C < D$ faar jeg ikke Brug for, hvis det blot er
Forholdet mellem $B$ og $D$, jeg vil bestemme; jeg lader da $A$ ude af Betragtning, og
$C$ benytter jeg kun som et Mellemled, der kan udskydes. Den analytiske Metode
gaar ud fra Tanken om en Opgave, og det gælder saa at udfinde, hvad der i de
foreliggende Rækker kan bruges, og hvad der er overflødigt. En Opdagelse gøres
ofte ved, at der skelnes mellem de foreliggende Emner og derpaa vælges, hvad der
bruges. Som \textsc{Poincaré} udtrykker det: Inventer, c'est discerner, c'est
choisir\footnote{\textit{Science et Méthode}, p. 48.}.

\medskip
12. Den formelle Logik har saa godt som udelukkende taget Hensyn til Begreber,
der uddrages af Anskuelsen, og som kan ordnes i en Trinrække, saaledes
at det følgende Begreb stedse har mindre Indhold, men større Omfang end det fore%
% --- p. 33 ---
\setcounter{footnote}{0}%
\opage{33}gaaende. Det er Klassifikationen, der fra \textsc{Aristoteles}' Tid har behersket den. En
Forudsætning er det her, at de første Begreber i den opstigende Række, nemlig
Individualbegreberne, er fuldt færdige, førend der gaas over til Almenbegreberne.

Men naar man nærmere overvejer Forholdet mellem Individualbegrebet og det
umiddelbart Givne, Emnerne, det skal svare til, ser man Nødvendigheden af at
skelne mellem to Slags Individualbegreber, som jeg i min Psykologi har foreslaaet
at kalde konkrete og typiske Individualbegreber. Naar et Emne (eller en
Emnegruppe) undergaar Forandringer og gennemløber forskellige Stadier, kan man søge
at bestemme det indre Forhold mellem disse forskellige Stadier, den Lov, efter
hvilken det ene afløser det andet, altsaa det givne Emnes, den foreliggende Helheds
Udviklingslov. Om hvert enkelt Stadium kan man have et konkret
Individualbegreb, om Helheden; som den fremtræder repræsenteret ved Stadiernes Række,
kan man have et typisk Individualbegreb. Naar et typisk Individualbegreb er dannet,
foreligger der en Række konkrete Individualbegreber i et bestemt indbyrdes
Forhold som Led af det. Saaledes i det Begreb, en Biografi formaar at give os af en
menneskelig Personlighed; de forskellige Udviklingstrin, denne har gennemløbet,
% DEFECT (p. 33): a stray mark stands before "af" ("Karakteristiken ·af").
giver deres Bidrag i et bestemt indbyrdes Forhold til Karakteristiken af
Personligheden som Helhed. Jo rigere og tillige bestemtere det typiske
Individualbegreb er,
des bedre vil man ud fra det tilbagegaaende kunne forstaa de enkelte Stadier i
Emnets Udvikling. Det betydningsfulde ved de typiske Individualbegreber er, at
% TRYKFEJL (p. 80): Individualforestillinger læs Individualbegreber
de ikke forholder sig til de tilsvarende konkrete Individualforestillinger som Art til
Individer, men som Helhed til Dele. Helheden indeholder her Loven, i Kraft af
hvilken det ene Stadium afløser det andet. Og her kan Indhold og Omfang vokse
med hinanden, hvad der ikke var muligt ved de Begreber, der dannes i
Klassifikationens Tjeneste.

Der kan i den nyere Tid spores en Tendens til at sætte typiske
Individualbegreber i Almenbegrebernes Sted. Til Almenbegreber kan man stige op, men
man kan fra dem ikke uden Iagttagelsens Hjælp stige ned til de underordnede
Begreber. Analysen fører os jo ikke fra Del til Helhed, men kun fra Helhed til
Del, og Almenbegreber giver jo kun en Del af Individualbegrebernes Indhold.
Allerede hos \textsc{Bruno} og \textsc{Bacon} fremtræder den nævnte Tendens, idet de, som
allerede omtalt, fortolker \textsc{Platon}'s „Ideer“ som Begreber om Love. \textsc{Spinoza}
afviser Almenbegreberne (abstracta et universalia) og kræver istedetfor dem Love, der
i lige Grad lægger sig for Dagen i Helheden og i dens Dele; i disse Love finder han
Udtryk for Tilværelsens Ejendommelighed, for noget, som er fast og evigt, og dog
individuelt\footnote{Hæc fixa et æterna, qvamvis sint singularia, tamen ob eorum
ubique præsentiam ac latissimam potentiam, erunt nobis tanqvam universalia.
\textit{De emend. intell.}}. Hermed hænger det sammen, at han opfatter hele Naturen
som et stort Individ. Hans Naturbegreb er et typisk Individualbegreb. \textsc{Leibniz}
har direkte opstillet det typiske Individualbegreb, idet han finder
Individualitetens Væsen udtrykt i Loven for Rækken af dens Virksomheder (lex
continuationis seriei suarum operationum, Lettre à Arnauld 1690; --- cette loi de
l'ordre qui fait l'individualité de chaque
% SIGNATURE LINE at the foot of p. 33 ("D. K. D. Vidensk. Selsk. Skr., 7.
% Række, hist. og filos. Afd. III. 2." and "37") omitted as printer's apparatus.
% --- p. 34 ---
\setcounter{footnote}{0}%
\opage{34}substance particulière, Lettre à Basnage 1698). Senere saa \textsc{Lotze}, at hvor et
% PRINTED AS IS (p. 34): "Bildungsgegesetz" (for "Bildungsgesetz").
Begreb udtrykker Loven for en Udvikling (ein Bildungsgegesetz), gælder ikke Regelen
om det omvendte Forhold mellem Indhold og Omfang\footnote{\textit{Drei Bücher der
Logik}, p. 50.}.

Spørgsmaalet er nu, hvorvidt saadanne typiske Individualbegreber staar til
vor Raadighed. Paa alle Erfaringsomraader maa vi først møjsommelig danne dem
paa Grundlag af Iagttagelser, der aldrig kan godtgøres at være udtømmende. Den
Karakteristik af en Personlighed, et Folk eller en Periode, som historisk Forsken
kan føre til, og som saa derefter kan tjene til Belysning af de enkelte
Karaktertræk og Begivenheder, maa selv vindes paa Grundlag af Beretning og Iagttagelse
og staar stedse som en Hypotese. Etiken virker ved Analyse af et typisk
Individualbegreb, et Forbillede (et idealt Samfund, eller en ideal Personlighed), der
bestemmes ved den Grundværdi, der bærer den etiske Tankegang, og ved
Betingelserne for at hævde og gennemføre den. Etiske Domme udsiger, hvilke Egenskaber
og Handlinger der i Henhold til dette ved Grundværdien bestemte Forbillede kræves
af det empiriske Samfund og den empiriske Personlighed\footnote{Smlgn. \textit{Det
psykologiske Grundlag for logiske Domme} (Vid. Selsk. Skr., 6. Række, hist. og
filos. Afd. IV, 6), p. 385. --- I Kap. V vil vi komme tilbage til dette Punkt.}.
Erkendelsesteorien
bygger paa et typisk Individualbegreb om menneskelig Tænkning, bygget paa den
Maade, denne lægger sig for Dagen i det uvilkaarlige Tankeliv og i Videnskabens
Historie, og de enkelte Kategorier og Grundsætninger findes ved Analyse af det
saaledes dannede Begreb.

Men skulde der ikke være Mulighed for en Tænken, ved hvilken de typiske
Individualbegreber, der danner Analysens Grundlag, konstrueres af Tanken selv,
uden at Iagttagelsens Hjælp er nødvendig ved hvert enkelt Skridt? Den Helhed,
der betinger Forstaaelsen, vilde da være Tankens eget Værk. --- Dette Spørgsmaal
fører os til at drøfte Forholdet mellem Logik og Matematik.

\medskip
13. \textsc{Kant} har allerede gjort opmærksom paa, at Rum og Tid ikke er blotte
Almenbegreber. De enkelte Rum forholder sig ikke blot til Rummet overhovedet,
som et Individ forholder sig til en Art, f.~Ex. som den enkelte Hest forholder sig til
Arten Hest, men tillige som Del til Helhed. Rummet er sammensat af Rum. „Man
kan,“ siger \textsc{Kant}\footnote{\textit{Kritik der reinen
Vernunft}\textsuperscript{1}, p. 25.}, „kun forestille sig et eneste Rum, og naar
man taler om mange Rum, forstaar man derved kun Dele af samme Rum.“ Og som det
forholder sig med Rummet, forholder det sig ogsaa med Tiden. I disse to Begreber har vi
Exempler paa Noget, der --- for at bruge det gamle Udtryk --- gælder i lige Maade
for det Hele og for Delene.

Vi finder ogsaa hos \textsc{Kant} Oplysning om, hvorledes disse typiske
Individualbegreber bliver til. Det sker ved, hvad han kalder „Synthesis des
Gleichartigen“. Jeg kan vedblive at føje Punkt til Punkt, Linie til Linie, Flade til
Flade, Volumen til Volumen, efter samme Lov, som jeg er begyndt med at anvende. Og
Tallet er
% --- p. 35 ---
\setcounter{footnote}{0}%
\opage{35}ikke Andet end den ved Syntese frembragte Enhed af en i en ensartet
Anskuelse given Mangfoldighed\footnote{\textit{Kritik der reinen
Vernunft}\textsuperscript{1}, p. 242, cfr. p. 142 f.}.

Rum, Tid og Tal er altsaa typiske Individualbegreber, vi kan konstruere ved
at sammenfatte en Række af aldeles ensartede Forskelligheder. Det er Lovens
Enhed indenfor Rækken, der gør Konstruktion af saadanne Helheder mulig.

I Analogi med de typiske Individualforestillinger, der ligger til Grund for
Matematiken, har \textsc{Kant} fremdeles, da han søgte Naturvidenskabens
Forudsætninger, konstrueret et Begreb om Erfaring eller Natur som en Helhed af
lovmæssig sammenhængende Iagttagelser\footnote{\textit{Kritik der reinen
Vernunft}\textsuperscript{1}, p. 110; 156, cfr. 2. Aufl. p. 185; 263.}, og det
er kun paa Grundlag af denne Konstruktion, at han kan opstille et helt System
af Grundsætninger som gyldige paa Forhaand for alle Fænomener. Især fremtræder
dette tydelig ved hans Begrundelse af Aarsagssætningen. Erfaring i
\textsc{Kant}'s Betydning er et typisk Individualbegreb. Hver enkelt speciel
Iagttagelse fører kun til et konkret Individualbegreb. Indenfor Erfaringen
(eller „Naturen“), i denne Betydning af Ordet, forholder alle Emner sig som
Grund til Følge eller omvendt. Dette logiske Grundforhold gentages overalt,
ligesom det indenfor Rummets, Tidens og Tallets Verden er samme anskuelige
Forhold, der gentages paa ethvert Punkt. Al streng Videnskab bygger altsaa paa
typiske Individualbegreber, hvis Konstruktion (der ikke altid sker med fuld
Bevidsthed) er mulig ved, at det staar i Tankens Magt at gentage sig selv ved
stedse at følge samme Regel for Fremskriden.

Alle disse Konstruktioner bygger paa et Princip, ved hvilket matematisk
Tænkning skiller sig fra den formelt logiske Tænkning. Rent logisk har
Gentagelse ingen Betydning. Det fører ikke til noget at kombinere et Begreb med
sig selv; hverken Indhold eller Omfang ændres derved. Overfor en absolut
Identitetsrække staar den rene Logik lige saa afmægtig som overfor en kaotisk
Forskelsrække. Men Erkendelsesteorien har flere Opgaver end dem, der stilles ved
den formelle Logik; den skal ogsaa undersøge Forudsætningen for matematisk og
exakt naturvidenskabelig Erkendelse. Og her er det, at Muligheden af en positiv
Betydning af Gentagelsen viser sig. Ved at gentage en og samme Tankeakt, noget,
som vi kan gøre uden Hjælp af Iagttagelser, kan vi frembringe Rækker af
ensartede Led, som saa kan sammenfattes til Helheder af forskellige Grader. Det
var netop det, \textsc{Kant} mente med „Synthesis des Gleichartigen“. Medens
den rene Logik stanser ved en kaotisk Forskelsrække, tager Matematiken her sin
Begyndelse ved at gøre opmærksom paa, at selv om alle Led er lige forskellige
indbyrdes, saa at ingen Ordning er mulig, har Leddene dog det fælles, at de
alle skyldes gentagne Tankeakter, og at disse Tankeakters Emner sammenfattes
til et Hele. Og medens den formelle Logik stanser ved en absolut
Identitetsrække, gør matematisk Tænkning opmærksom paa, at det dog er en Række
af forskellige Tankeakter, der har frembragt de identiske Begreber. Fra begge
Sider ses nu Muligheden af en identisk varierende Forskelsrække. En saadan
ligger tydelig til Grund for Talbegrebet, men ikke
% SIGNATURE MARK "37*" at the foot of p. 35 omitted as printer's apparatus.
% --- p. 36 ---
\setcounter{footnote}{0}%
% TRYKFEJL (p. 80): fra læs for
\opage{36}mindre fra Rækken af Steder i Rummet, af Øjeblikke i Tiden og af Trin
paa en Gradeskala, naar man lader de kvalitative Forskelle mellem Steder
indbyrdes, Øjeblikke indbyrdes og Grader indbyrdes ude af Betragtning.

I \textit{Den menneskelige Tanke} (p. 181) har jeg vel fremhævet Betydningen
af, at det er muligt ved Begreberne Tid, Tal, Grad og Sted at eliminere
Kvalitetsforskellene og derved bane Vej for en rationel Behandling, der ad
Analogiens Vej ogsaa kan komme Erkendelsen til Gode overfor andre
Kvalitetsforskelligheder. Men jeg har der endnu ikke klart set, at det er det
samme Princip, Gentagelsens positive Betydning, der gør sig gældende ved alle
hine fire Kategorier. Kun ved Tallets Genesis har jeg fremhævet denne
Betydning, og det er jo ogsaa der, den historisk først er blevet paapeget og
tillige fremtræder tydeligst.

\textsc{Leibniz} er den første, der har set, at det er Gentagelsens positive
Betydning, der danner Grænsen mellem Logik og Matematik\footnote{Si idem secum
ipso sumatur, nihil constituitur novum, seu $A + A = A$. Equidem in numeris
$4 + 4$ facit 8, seu bini nummi binis additi faciant qvatuor numeros, sed tunc
bini additi sunt alii a prioribus; si iidem essent, nihil novi prodiret.
\textit{Non inelegans specimen demonstrandi in abstractis.} (Opera
philosophica. \textsc{Ed. Erdmann}, p. 95).}, kun at han ikke har set, at selv
om det, der tælles, stedse er det samme (som ved Tordenskjolds Soldater, eller
ved Muddermaskinens Skuffer), saa foreligger der dog ogsaa i saadanne
Tilfælde en Række Tankeakter, der kan sammenfattes. \textsc{Kant} lægger, som
vi har set, „Syntesen af det Ensartede“ til Grund for sin Erkendelsesteori og
har konstrueret sit Begreb om „Erfaring“ eller „Natur“ i Analogi med
Begreberne Rum, Tid og Tal. Man har ment, at \textsc{Kant} her skulde være
paavirket af \textsc{Maupertuis}, der til Forklaring af
% TRYKFEJL (p. 80): Erfaringer læs Erfaringen
Matematikens Mulighed henviser til, at i Erfaringer raader en vis Mulighed for
Gentagelse (réplicabilité), idet visse simple Naturting kommer igen uden
væsentlig Forskel, og fra saadanne Erfaringer skulde da Matematiken hente sit
Begreb om en Helhed, der fremkom ved simpel Gentagelse af en
Ener\footnote{\textsc{Maupertuis} (Histoire de l'Académie de Berlin. 1756.
p. 394) finder Gentagelighed (réplicabilité) som faktisk Egenskab ved Tal og
Udstrækning, og som den eneste Egenskab, ved hvilken de adskiller sig fra
% TRYKFEJL (p. 80): samtlige læs sanselige
andre samtlige Egenskaber, men han gaar ikke ind paa den Proces, ved hvilken
Tal og Udstrækning bliver mulig. \textsc{Cassirer}
% PRINTED AS IS (p. 36, note 2): "Erkennnissproblem" (three n) for
% Erkenntnisproblem; and "725 f)" with no stop after f.
(\textit{Das Erkennnissproblem in der Philosophie und Wissenschaft der neueren
Zeit,} II, p. 334; 495; 725 f), der har henledet Opmærksomheden paa
\textsc{Maupertuis}'s Afhandling, mener, at \textsc{Kant} har fundet et
Tilknytningspunkt i denne. Men \textsc{Kant} opfatter Gentagelsen som en
Aktivitet, ved hvilken rene Begreber frembringes. Og foreløbig anvender (som
\textsc{Cassirer} selv oplyser) \textsc{Kant} denne Opfattelse kun paa
Talbegrebet, først senere (efter Opdagelsen af Syntesebegrebet) ogsaa paa
Udstrækningen, medens \textsc{Maupertuis} taler om begge paa én Gang.}. Men
\textsc{Kant} lægger Vægten paa Tankens egen Evne til at gentage samme Akt. Og
først da \textsc{Kant} ad
% DEFECT (p. 36): a stray raised dot stands before "anden" at the head of
% the line ("·anden Vej"); a mark in the copy, not text. Not transcribed.
anden Vej\footnote{Se herom \textit{Kontinuiteten i \textsc{Kant}'s
filosofiske Udviklingsgang} (Vid. Selsk. Skr., 6. Række, hist. og filos.
Afd. IV, 1.) p. 27 f; 52---54.} havde fundet Syntesebegrebet, kunde han se den
store erkendelsesteoretiske Betydning af Gentagelsens positive Karakter, som
„Synthesis des Gleichartigen“ i Kraft af „Apperceptionens gennemgaaende
Identitet“\footnote{\textit{Kritik der reinen Vernunft}\textsuperscript{2},
p. 133.}. Det er mærkværdigt, og det er tillige et vigtigt Udgangspunkt for en
Kritik af \textsc{Kant}'s Erkendelsesteori, at medens han ved de typiske
Individualbegreber, der danner Grundlaget
% --- p. 37 ---
\setcounter{footnote}{0}%
\opage{37}for den rene Matematik, udtrykkelig og vidtløftig behandler
Spørgsmaalet om Berettigelsen af deres Anvendelse paa givne, af os ikke
konstruerede Emner, opkaster han ikke noget tilsvarende Spørgsmaal med Hensyn
til det af ham konstruerede Begreb „Erfaring“ eller „Natur“. Dette kommer
vistnok af, at \textsc{Kant} ikke selv har været helt klar over, at hvad han
kaldte „syntetiske Domme a priori“, var et med de Domme, der skyldes en
„Synthesis des Gleichartigen“. Og dette hænger igen sammen med hans uheldige
Distinktion mellem „Æstetik“ og „Analytik“.

I det 19. Aarhundrede har \textsc{Boole} indgaaende drøftet Forholdet mellem
Logik og Matematik fra Synspunktet om Gentagelsens positive Betydning.
Ligesaalidt som \textsc{Kant} har han kendt \textsc{Leibniz}'s Afhandling om
dette Emne (den blev først trykt 1840, og \textsc{Boole} har ikke kendt
\textsc{Erdmann}'s Udgave). Han udtrykker Forskellen
saaledes\footnote{\textit{Investigation of the laws of thought.} (1854), p. 48
ff.}, at rent logisk er $x = x^2$, men matematisk kun for $x = 1$ eller $= 0$.
Og i den nyeste Tid har \textsc{Henri Poincaré} kastet Lys over dette
Forhold. Hverken Erfaring eller rent logisk Analyse kan, hævder han, forklare
den matematiske Induktion, d.~v.~s. Bevisførelsen par récurrence, hvis Princip
er, at hvad der gælder for $n$, gælder ogsaa for $n + 1$. Dette Princip svarer
ifølge \textsc{Poincaré} til, hvad \textsc{Kant} kalder en
apriorisk-syntetisk Dom, og ved det bliver paa en Gang en Udvidelse af vore
Begrebers Indhold og af deres Omfang mulig. Selve Principet har sin Grund i
Muligheden af at gentage en Akt, naar denne Akt overhovedet er
mulig\footnote{\textit{Science et Hypothèse}, p. 23. --- Smlgn. dog
\textsc{Paul Tannéry}'s Bemærkninger i \textit{Bulletin de la Société
française de Philosophie,} III, p. 126 f.}.

Det er dog ikke alle Gentagelser af en Tankeakt, der er lige frugtbare. Den
Akt at blive sig selv bevidst, gøre sit Jeg (sine Tilstande, sit Arbejde, sine
Kaar) til Genstand for sig, kan formelt stadig gentages. Det Jeg, der bliver
sig selv bevidst, kan selv blive Genstand for en ny Selvbevidsthedsakt, og
saaledes fremdeles. En
% FORMULA (p. 37): the sign between the terms is printed as a left curly brace
% (checked at 300 dpi), set \lbrace; four spaced dots follow S_4.
saadan Række ($S_1 \lbrace S_2 \lbrace S_3 \lbrace S_4$~.~.~.~.) er allerede omtalt ovenfor
i Anledning af Muligheden af en erkendelsesteoretisk Prøvelse af
Erkendelsesteorien. Her vender vi tilbage til den for at paavise dens Forskel
fra de Rækker, ved hvilke Begreberne Tid, Tal, Grad og Sted konstrueres.
Medens ved disse fire Rækker det, som gentages, stedse er af samme Kvalitet,
saa at det følgende stedse er en blot Fortsættelse af det foregaaende, er i
Objektiveringsrækken hvert nyt Led, hvert nyt $S$, netop som Virkning af
Objektiveringen, mere eller mindre forskelligartet i Forhold til det
foregaaende. Det at blive sig selv bevidst er aldrig uden Virkning. Det kan
betyde Kritik eller Forstaaelse, men det kan ogsaa betyde subjektiv Hildethed i
Modsætning til Opgaaen i en Sag eller en Stræben; det kan betyde større
Rigdom, men det kan ogsaa betyde Fattigdom. Det er Aandslivets historiske
Karakter, der udelukker, at Objektiveringsaktens Gentagelse kan faa den
Betydning af videnskabelig Konstruktion, som de fire omtalte Rækker besidder.
Objektiveringsprocessen kan blive af stor psykologisk og historisk Betydning,
men den yder intet nyt erkendelsesteoretisk Grundlag, saaledes som
Romantikens Filosofi mente.

% --- p. 38 ---
\setcounter{footnote}{0}%
\medskip
\opage{38}14. Allerede den formale Tænkning har sit Sandhedskriterium. I den
rene Logik og Matematik er alt sandt, hvad der ikke indeholder nogen
Modsigelse. Men dette negative Kriterium forudsætter en forudgaaende positiv
Udfoldelse af Anskuelser og Forestillinger, en uvilkaarlig, maaske ubevidst
Dannelsesproces. Herved virker den i Bevidsthedens Natur saa dybt grundede
Trang og Tendens til Supplering og Expansion, til helt og fuldt at gennemleve
en Tilstand og til at fortsætte den saalænge som muligt. De saakaldte Love for
Forestillingers Association er kun særlige Former for denne Tendens, som i
\textsc{Hume}'s Erkendelsespsykologi spiller saa stor en Rolle. Ligheds- og
Forskelsforhold gør sig her samtidig gældende paa en mærkelig Maade. Thi hvad
der udløser den nævnte Expansion, vil især være Iagttagelser og Oplevelser,
som ved deres Styrkegrad eller ved deres Følelsesvirkning sætter Bevidstheden i
stærk Bevægelse. Der gør sig da en dobbelt Trang gældende: til at fastholde
Emnet og dvæle ved det, --- men ogsaa til at faa fuldt Afløb for den udløste
Energi ved at gaa ud over Emnet. Denne dobbelte Trang udløses ved ikke blot at
fortsætte Tilstanden saa længe som muligt, men ved ogsaa at gaa over til mere
eller mindre beslægtede Tilstande. Her kommer de forskellige Lighedsgrader,
ikke mindst Analogien, til at gøre sig gældende. Fantasi (i snevrere
Betydning) er en speciel Form for denne Art psykiske Processer. En Stansning i
den uvilkaarlige Udfoldelse kommer først, naar paa en eller anden Maade to
Emner eller Led i de uvilkaarlig dannede Rækker sammenholdes uden at kunne
forenes, idet det ene river ned, hvad det andet bygger op. Da knager den
uvilkaarlig opbyggede Helhed i sine Fuger.

Skønt det negative Kriterium saaledes kan træde i skarp Modsætning til den
positive Udfoldelse og Helhedsdannelse, som selve dets Anvendelse
forudsætter\footnote{Positive Domme gaar derfor forud for negative.
\textit{Den menneskelige Tanke}, p. 76 f.}, er det dog selv et Vidnesbyrd om
Tankens syntetiske, helhedsdannende Natur. Bevidsthed om Modsigelse
forudsætter, at Leddene, der modsiger hinanden, sammenholdes til en forsøgt
Helhedsdannelse. Hvis Emnerne hvert for sig hørte til sin Helhed, vilde
Modsigelsen ikke komme for Dagen, idetmindste ikke før de to Helheder var
blevne analyserede, og deres Dele forsøgtes kombinerede med hverandre. Syntesen
viser sig netop tydelig i en Selvmodsigelse, som man bliver sig bevidst,
ligesom overhovedet alle Delthedstilstande (Tvivl, Sorg, Anger) tydeligst
godtgør Bevidsthedens syntetiske Karakter. Det negative Kriterium er altsaa
Udtryk for en Ejendommelighed ved selve Tanken. Vor Aand er allerede her
Forbillede for alt, hvad der skal kunne blive Genstand for den, eller som
\textsc{Kant} i en ældre Optegnelse\footnote{\textit{Lose Blätter aus
\textsc{Kant}'s Nachlass,} I, p. 19. --- Jeg har allerede anført dette Sted i
\textit{Kontinuiteten i \textsc{Kant}'s Udviklingsgang,} p. 15.} har udtrykt
det: „Ich bin selbst das Original aller Objecte“.

Den formale Tænken har dog ogsaa et positivt Kriterium. Den skyder ikke blot
det modsigende ud som uholdbart, men den frembringer selv nye Forbindelser,
nye umiddelbare Forhold gennem Substitution og Elimination. Snart gaar
Substitutionen forud for Eliminationen, snart falder de sammen. Har jeg to
Domme $A = B$
% --- p. 39 ---
\setcounter{footnote}{0}%
\opage{39}og $B = C$, kan jeg sætte Udtrykket for $B$ i den ene Dom ind
istedetfor $B$ i den anden Dom, og faar da $A = C$, idet $B$ saa kan falde
bort. Har jeg derimod Rækken $A = B = C$, skyder jeg $B$ ud, under
Forudsætning af, at det $B$, jeg kommer til fra $A$, er det samme som det, jeg
kommer til fra $C$. Det er en og samme Proces, udtrykt paa to Maader.

Det negative Kriterium giver Svar paa det Spørgsmaal, der stiltiende ligger i
den trygge, positive Udfoldelse: Hvorfor ikke? Hvorfor ikke blive ved som
hidtil? Det positive Kriterium giver Svar paa Spørgsmaalet: Hvorfor? Hvorfor
opstille denne nye Forbindelse? --- Rationaliteten gør sig gældende ad begge
Veje\footnote{\textit{Den menneskelige Tanke}, p. 204---207.}. Der kan drages
Slutning, dannes Overgang fra Grund til Følge paa begge Maader.

Grund og Følge (f.~Ex. en Rækkedannelse og de nye Forhold, der kan udledes af
den) udgør en logisk Totalitet. Men en fuldkommen logisk Totalitet faas kun,
naar Følgen ikke blot kan udledes af Grunden, men naar Slutningen kan vendes
om, saa at der er fuld
% ID-EST MARK (p. 39): the reversed c "ɔ:" typed as U+0254 + colon, as in
% the Brøchner transcriptions; the preamble needs a rendering for U+0254.
(ɔ: tvesidig) logisk Ækvivalens mellem Leddene i den logiske Totalitet. Men
hvor Slutningen er naaet ved Elimination (hvad enten denne indtræder efter
Substitution eller ikke), kan den ikke vendes om. Af Rækken $A = C$ kan jeg
ikke udlede Rækken $A = B = C$. Saalænge Præmisserne er mere indholdsrige end
Konklusionen, --- altsaa saalænge Tanken arbejder paa Grundlag af et
umiddelbart Givet, der er en større Helhed end den selv er, --- saalænge kan
den logiske Ækvivalens, der er det logiske Ideal, ikke naas. Naar man paa
Forhaand udelukker alle Led i de foreliggende Erfaringer, der ikke indgaar som
Led i Slutsætningen, kan man naturligvis naa Ækvivalens. Saaledes stod det
klart for den moderne Naturvidenskabs Grundlæggere, at man maatte se bort fra
Sansekvaliteterne, hvis man vilde have et rationelt System, der kunde staa som
intellektuelt Udtryk for Forandringerne i Verden. Der staar da her det
Spørgsmaal tilbage, hvorledes Forholdet er mellem det intellektuelle Resultat,
der saaledes er vundet, og de Kvalitetsrækker, hvis særegne Egenskaber man har
set bort fra. Dette Spørgsmaal kommer vi tilbage til i en anden Sammenhæng.

Den Sandhed, Tanken kan finde, er en ny Helhed, der afløser den uvilkaarlig
dannede Helhed, der var Grundlaget for Tankearbejdet. Og denne Substitutions
Berettigelse beror paa, at den Forudsætning, at Sandheden maa udgøre en
Helhed, kommer mere til sin Ret, end i den umiddelbart givne Helhed. Denne
Forudsætning ligger til Grund ved ethvert Spørgsmaal, der opkastes, og enhver
Opgave, der stilles. Indenfor den Helhed, Sandheden forudsættes at udgøre, kan
der intet mangle, som Sammenhængen kræver; intet kan staa ubestemt; intet Led
kan stride mod noget andet Led; --- ja Leddene eller Delene maa gensidig
bestemme hverandre og kan udledes af hverandre i Kraft af den Lov, der
forbinder dem indenfor Helheden.

\medskip
15. Kun ad Analysens Vej bliver videnskabelig Sandhed til, men selve
Analysen forudsætter, som jeg har forsøgt at vise i det foregaaende, en
uvilkaarlig Syntese. Som jeg har udtrykt det i min Afhandling om Intuition: vi
lever umiddelbart
% --- p. 40 ---
\setcounter{footnote}{0}%
\opage{40}og uvilkaarlig i konkrete og praktiske Intuitioner, men Videnskabens
Vej gaar gennem analytiske Intuitioner, d.~v.~s. gennem Bevidsthedsakter, der
konstaterer Identitet eller Forskel mellem foreliggende Emner. Ethvert enkelt
Led i en Bevisførelse bestaar i en analytisk Intuition\footnote{I den nævnte
Afhandling har jeg nævnt \textsc{Descartes}, som den, der tydelig har opstillet
Begrebet analytisk Intuition. Men det er ikke forsvundet efter hans Tid. Det
findes baade hos \textsc{Locke} (\textit{Essay} IV, 2, 1) og hos \textsc{Hume}
(\textit{Treatise} I, 3, 1. 3).}. Vi bliver os da ved dens Hjælp bevidst,
hvilke Elementer der indeholdes i de foreliggende Emner. Den er Nerven i
enhver sammenhængende Tankegang, betegner enhver af Overgangene indenfor den.
Gennem en Række af analytiske Intuitioner kan der saa maaske opstaa en
syntetisk Intuition, der giver os en sammenfattende Overskuen over Tankevejene
og deres Resultater. Men selve Analysen, Bestemmelsen af Identiteter og af
Forskelligheder ved de enkelte Tankegange forudsætter stadig en Sammenfatten af
de Led, der ses som identiske eller som forskellige. Der finder altsaa under
Erkendelsens Udfoldelse en stadig Vekselvirkning Sted mellem Syntese og
Analyse.

Det var derfor en umulig og unødvendig Distinktion, \textsc{Kant} vilde
indføre, naar han skarpt skelnede mellem syntetiske og analytiske Domme.
Enhver Dom maa være syntetisk, thi der forbindes i den noget, der ikke før har
været forbundet for Tanken, selv om det nu viser sig at være forbundet i de
Emner, Tanken sysler med. Og enhver Dom er analytisk, idet dens Prædikat maa
være indeholdt i det Emne, der drøftes, dersom ellers Dommen er rigtig.
\textsc{Kant} tænkte ved analytiske Domme paa saadanne Domme, i hvilke
Subjektet (Udgangsforestillingen, Emnet) udgør hele den Totalitet, som
analyseres. Men det kan dog her ingen principiel Forskel gøre, om det er et
enkelt Emne eller en Flerhed af Emner, jeg analyserer. Begrænsningen til et
enkelt Emne, en enkelt Udgangsforestilling vil altid være mere eller mindre
kunstig. Naar jeg siger „Manden er god“, og danner denne Dom paa første
Haand, finder jeg Anledningen til Prædikatet „god“ i en Iagttagelse af Mandens
Handlemaade i forskellige Situationer, hverken i det abstrakte Begreb eller
det konkrete Anskuelsesbillede af Manden. Jeg bringer Manden i forskellige
Relationer. Og saaledes ogsaa, for at bruge \textsc{Kant}'s Yndlingsexempel,
naar jeg siger „Legemer er tunge“. Jeg har da set de enkelte Legemer som Led i
et helt System af Legemer, og det er ved Analyse af Forholdene indenfor dette
System, at Prædikatet „tung“ opdukker. Det er altid typiske
Individualforestillinger, der analyseres i de saakaldte „syntetiske“ Domme.
Derfor maatte, som vi ovenfor har set, \textsc{Kant} konstruere det typiske
Individualbegreb „Erfaring“ eller „Natur“ for at kunne udlede sine
„apriorisk-syntetiske“ Domme.

De saakaldte analytiske Domme skal efter \textsc{Kant} intet nyt bringe. Men
hvorfor danner vi da saadanne Domme? Hvilken Trang tilfredsstiller de? Domme
er oprindelig Prædikatsdomme. Det er Prædikatet, som bringer det nye. Derfor
kan Sproget have Sætninger, i hvilke intet logisk Subjekt udtrykkes. ---
\textsc{Kant} indrømmer da ogsaa, at selv de „analytiske“ Domme bringer noget
nyt. Thi han siger i anden Udgave af Fornuftkritiken, at i saadanne Domme
indeholdes Prædikatet „versteckter
% --- p. 41 ---
\setcounter{footnote}{0}%
\opage{41}Weise“ eller „verworn“ i Subjektet. I første Udgave udtrykkes det
samme ved, at i analytiske Domme bliver mit Begreb om Subjektet „aus
einandergesezt und mir selbst verständlich gemacht“\footnote{\textit{Kritik
der reinen Vernunft.} 2. Udg. Einleitung §~4. --- Første Udgave: Einleitung.
Von dem Unterschiede analytischer und synthetischer Urteile.}. Men saa er her
jo noget nyt: det skjulte bliver draget frem, det forvirrede klaret, og jeg
forstaar nu først mit eget Begreb. Det er kun et Spørgsmaal, om dette kan
naas, uden at man gaar ud over det enkelte Emne og bringer det i Sammenhæng med
andre Emner. Der gives egentlig slet ikke rent analytiske Domme. End ikke
Dommen $A = A$ er rent analytisk. Denne Dom er en Definition af, hvad et
Begreb er, og tillige en Paastand om, hvad der er Forudsætning for enhver
Bevisførelse, enhver analytisk Intuition. Den Helhed, der ligger til Grund for
dens Opstilling, er det Krav, som Bevisførelse stiller til de Forestillinger,
der skal indgaa som Led i den. Sætningen ligger altsaa „versteckter Weise“ i
enhver Slutning, ethvert Forsøg paa Bevis, ikke i $A$ i og for sig\footnote{I
den senere Tid er Sondringen mellem syntetiske og analytiske Domme gentagne
Gange bleven drøftet i det franske filosofiske Selskab. Bulletin II (1902):
\textsc{Couturat}: \textit{Sur le rapport de la Logique et de la Métaphysique
de \textsc{Leibniz}.} Bulletin III (1903): \textsc{Paul Tannéry}: \textit{La
% PRINTED AS IS (p. 41, note 2): "cantienne" (broken "can-/tienne" at the
% line end) for kantienne.
valeur de la classification cantienne des jugements en analytiques et
synthétiques.} \textsc{Couturat} vil opfatte alle Domme som analytiske;
\textsc{P. Tannéry} hævder, at enhver Dom alt efter Synspunktet kan betragtes
som analytisk eller som syntetisk.}.

Som vi ovenfor (§~13) har set, formuleres matematiske Grundsætninger paa
Grundlag af en Intuition eller „Induktion“, der stedse kan udvides i Kraft af
samme Lov eller Tendens, ifølge hvilken den er bleven til. \textsc{Poincaré}
har, forekommer det mig, klart gendrevet moderne Forsøg paa at opfatte de
matematiske Grundsætninger som „analytiske“ i \textsc{Kant}'s Betydning af
Ordet, og han har fortolket \textsc{Kant}'s „syntetiske Sætninger a priori“
% PRINTED AS IS (p. 41): "recurrence" unaccented (p. 37 has "récurrence").
som grundede paa Analyse af „matematisk Induktion“ (loi de recurrence).

Hvor langt denne Analyse skal gaa, naar det gælder Opstilling af de første
Definitioner i Geometrien, er et Spørgsmaal, der er blevet besvaret
forskellig. Efter den strenge platonisk-euklidiske Opfattelse maa man her gaa
til Grænsen, --- til Punktet som Grænse for Linien, til Linien som Grænse for
Fladen, til Fladen som Grænse for Legemet, og de matematiske Sætninger gælder
kun for disse ideale Emner. Tangenten berører Cirkelen kun i et eneste Punkt,
og mellem to Punkter kan der kun drages en eneste ret Linie. I Oldtiden
% PRINTED AS IS (p. 41): "Protogoras" for Protagoras (so spelt on p. 42).
indvendte \textsc{Protogoras}\footnote{\textsc{Diels}: \textit{Fragmente der
Vorsokratiker}\textsuperscript{1}, p. 520.} mod en saadan Opfattelse, at
enhver virkelig Tangent berører Cirklen i mere end ét Punkt, og henviste i det
hele til de sanselige Linier istedetfor til de ideale. I nyeste Tid har
\textsc{Hjelmslev} givet en Fremstilling af den elementære Geometri, der ikke
opstiller de ideale Definitioner, men gaar ud fra de Linier og Figurer, som kan
tegnes, og i disse har en Tangent stedse en hel Strækning fælles med Cirklen,
og der kan være mere end én ret Linie mellem to Punkter. Han lader sig nøje
med at sige, at der i
% SPACED (p. 41): "det mindste" and "indeholdt" are letterspaced in print.
\emph{det mindste} er én ret Linie, som gaar gennem to givne Punkter, og at
Fodpunktet af den vinkelrette Linie paa Tangenten maa være \emph{indeholdt} i
den
% SIGNATURE LINE at the foot of p. 41 ("D. K. D. Vidensk. Selsk. Skr., 7.
% Række, hist. og filos. Afd. III. 2." and "38") omitted as printer's apparatus.
% --- p. 42 ---
\setcounter{footnote}{0}%
\opage{42}Strækning, som Cirkel og Tangent har fælles\footnote{\textsc{Hjelmslev}:
\textit{Om den rette Linies Bestemmelse ved to Punkter.} (Vid. Selsk. Overs.
1916), p. 183.}. Medens den platonisk-euklidiske Opfattelse gaar til Grænsen i
sin Bestemmelse af Begreberne Punkt og Linie, --- en Grænse, der i vor
Anskuen aldrig naas, --- holder en Opfattelse som \textsc{Hjelmslev}'s sig til,
at vi faktisk bevæger os paa Anskuelsens Omraade, og at der ved alle de
Tangenter, vi faar med at gøre, stedse vil være flere fælles Punkter for den og
Cirklen. Hovedsagen er, at Fodpunktet er indeholdt i den for Cirklen og
Tangenten fælles Strækning\footnote{\textsc{Hjelmslev}: \textit{Elementær
Geometri,} I, p. 31.}. --- Til syvende og sidst bliver Forskellen mellem de to
Opfattelser ikke stor, men den er af erkendelsesteoretisk Interesse, idet den
kaster Lys over Tankearbejdets Forhold til Anskuelsen. Begge Opfattelser
fastholder, at vi faktisk ikke naar ud over Anskuelsen; men den ene lægger Vægt
paa Tilnærmelsen til den ideale Maalestok, medens den anden lægger Vægt paa, at
Idealet har sit Sted i selve det Anskuede, skønt Anskuelsen indeholder Mere end
Idealet\footnote{Maaske kunde \textsc{Hobbes}' Definition af et Punkt være af
Interesse i denne Sammenhæng. \textit{De Corpore} XV, 2: --- per punctum non
intelligi id, qvod qvantitatem nullam habet~.~.~. sed id, cujus qvantitas non
consideratur, hoc est, cujus neque qvantitas neque pars ulla inter
demonstrandum computatur.}. Om \textsc{Protagoras} tænkte sig Sagen paa denne
Maade, eller han blot benyttede sin Opfattelse i skeptisk Øjemed, maa staa hen.

% CHAPTER HEAD (p. 42) as printed: a short rule, then two centred lines
% ("IV." / title with final stop) over a short rule. The Indhold (p. 80)
% gives the same title.
\section*{IV.\ Totalitetsbegrebet og de reale Kategorier.}
\addcontentsline{toc}{section}{IV.\ Totalitetsbegrebet og de reale Kategorier}

16. De Helheder, der dannes ved Hjælp af de formale Kategorier, er paa en Gang
Udtryk for den menneskelige Tankes Natur og Forbilleder for den Virksomhed, ved
hvilken Emner, Tanken ikke selv har frembragt, bearbejdes. Tanken kan i sin
bearbejdende Virksomhed overfor Emnerne ikke fornægte sig selv. Derfor
forudsætter de reale Kategorier de formale, ligesom disse forudsætter de
fundamentale. Dette fremgaar tydelig deraf, at de formale (logisk-matematiske)
Grundbegreber indeholder Grundlaget for al Forstaaelse, hvad det saa er for
Emner, den angaar. At forstaa noget i strengere Betydning er ikke blot at
genkende det, men at se det som Følge af noget, som allerede kendes. Der
forudsættes altsaa et Forhold mellem Grund og Følge. Dette Forhold kan igen
være af den Art, at der kun er ensidig Ækvivalens, d.~v.~s. at Forholdet ikke
kan vendes om, eller at der gælder tvesidig Ækvivalens, saa at man kan gøre
Følgen til Grund og Grunden til Følge. Derfor ser vi ogsaa, at det ved den
vigtigste af de reale Kategorier, som vi foreløbig holder
% --- p. 43 ---
% p. 43: stray marks at the left edge of the scan beside the head and down the
% margin (the layer reads "k", "►", "b") are marks in the scanned copy, not
% text of the work.
\setcounter{footnote}{0}%
\opage{43}os til, Kausaliteten, drejer sig om at paavise ensidig eller tvesidig
Ækvivalens mellem to Emner. Forudsætningen er, at der af de Emner, Talen er om,
kan dannes rationale Rækker.

Hvor saadanne Rækker ikke kan dannes, og derfor intet Ækvivalensforhold, hverken
ensidigt eller tvesidigt, kan dannes, er streng Videnskab ikke mulig. Der kan
kun være Tale om Ordning, Inddeling og Sammenligning. Grunden til, at rationelle
Rækker ikke kan dannes, ligger tilsidst i Emnernes skiftende Kvaliteter. Det er
ikke Kvaliteterne i og for sig, heller ikke Tidsforholdet i og for sig, der gør
Vanskeligheder\footnote{Det følgende er en Berigtigelse af Fremstillingen i „Den
menneskelige Tanke“, p. 174 f.}.

Kvaliteterne kunde tænkes existerende en Gang for alle, realiserede i alle
mulige Omfang og i alle mulige Grader, dannende et stort, hvilende System. I
\textsc{Spinoza}'s System er denne Mulighed greben, som iøvrigt allerede græsk
Tænken lige fra \textsc{Parmenides} stilede hen imod. \textsc{Spinoza} reducerer
først alle Kvalitetsforskelle til den ene Forskel mellem Aandeligt og Materielt,
og disse to Kvaliteter opfattes som koordinerede Attributer for Tilværelsen. I
et saadant System kan der kun være Tale om formale Helheder. Alle Forhold
indenfor hvert enkelt Attribut (\textsc{Spinoza} antager, at der er uendelig
mange foruden de to, vi kender) er logisk-matematiske. Forholdet mellem Grund og
Følge bliver det eneste gældende, --- og Følgerne behøver end ikke at drages; de
er evig indeholdte i deres Grunde, tilsidst i en eneste Grund.

Tidsforholdet i og for sig kunde ogsaa tænkes rent rationelt. Ifølge
Tidsbegrebet i dets Renhed ligger hvert Øjeblik mellem to andre Øjeblikke, og da
Overgangen mellem hvilkesomhelst to Øjeblikke sker ganske paa samme Maade som
mellem to hvilkesomhelst andre Øjeblikke, faar vi en identisk varierende
% DEFECT (p. 43): a stray tick over the n of "indenfor" (the layer reads
% "inderifor"); not transcribed.
Forskelsrække, indenfor hvilken de enkelte Led ikke frembyder anden Forskel end
den, der bestaar i deres forskellige Plads i Rækken. Tidstotaliteter
(Øjebliksgrupper) kan ordnes saaledes, at det ene Tidsrum, ikke blot bestandig
ligger mellem to andre Tidsrum, men tillige saaledes, at der dannes en
fremskridende Forskelsrække, i hvilken ethvert Led er større (eller mindre) end
det foregaaende og mindre (eller større) end det følgende. Det rene Tidsbegreb
tillader altsaa meget vel Rationalitet\footnote{\textsc{Kant} lader ikke det
rene Tidsbegreb komme til dets Ret som Grundlag for en egen lille rationel
Verden, der i sin Udformning dog behøver Analogi med Rummet og Anvendelse af
Størrelsesbegrebet. Han paralleliserer det for meget med Rummet og er tilbøjelig
til at sætte det i nøje Forhold til Tallet. I \textit{Kritik der reinen
Vernunft}\textsuperscript{1}, p. 142 f. defineres Tallet som den Enhed i
Sammenfatten af en ensartet Mangfoldighed, der betinges ved, at jeg frembringer
Tiden ved Analyse af Anskuelsen. I \textit{Prolegomena} § 10 lægges Tiden til
Grund for Aritmetiken, ligesom Rummet for Geometrien.}.

Hverken Kvalitetsforskelle for sig eller Tidsforskelle for sig strider mod
Rationaliteten. Men det er Tidsindholdets Skiften, Kvalitetsforandringerne, der
medfører Problemet. Hvad Indhold angaar, kan man ikke uden videre slutte fra det
ene Øjeblik eller det ene Tidsrum til det andet. Er det muligt paa Grundlag af
Kvalitetsforandringer som givne at bygge strenge Slutninger (med ensidig eller
tvesidig Ækvivalens)?
% SIGNATURE MARK "38*" at the foot of p. 43 omitted as printer's apparatus.

% --- p. 44 ---
\setcounter{footnote}{0}%
\opage{44}Ser vi paa Videnskabens Historie, viser det sig, at Tænkningen er
gaaet den Vej at gennemarbejde visse vigtige, stadig genkommende
Kvalitetsforskelle saaledes, at de kan afgive Skemaer for de skiftende
Kvaliteter, og saaledes, at der kan dannes Rækker, indenfor hvilke hvert Led er
bestemt ved sin Plads i Forhold til de andre Led. Saadanne Skematismer er
dannede, foruden af Tidsbegrebet, af Begreberne Tal, Grad og
Sted\footnote{\textit{Den menneskelige Tanke}, p. 180---201.}. Det er dem, den
exakte Erfaringsvidenskab skylder sin Tilblivelse. Kvalitetsforskellighederne
(baade de samtidige og de successive) ordnes i Rækker af den Art, at Leddene
forholder sig til hverandre som Leddene i et eller flere af hine Skemaer.

Men Forholdet mellem Skema (f.~Ex. den rene Tid eller det rene Tal) og Emner
(f.~Ex. den kvalitativt udfyldte Tid eller de kvalitativt forskellige Fænomener)
er ikke Identitet, men Analogi. For Iagttagelsen er og bliver ethvert Øjeblik
opfyldt paa en Maade, der er kvalitativt forskellig fra den Maade, paa hvilken
andre Øjeblikke er opfyldte, og medens Overgang fra et Tal til et andet stedse
kan foregaa paa samme Maade, er og bliver Overgangen fra en Kvalitet til en
anden forskellig fra Overgangen mellem to andre Kvaliteter. Tidsindholdet kan
ændres ved Overgangen fra et Øjeblik til et andet, skønt Tidens Lov er den
samme. Og medens f.~Ex. Overgangen fra en Brydningsgrad af Lyset til en anden er
rent kvantitativ, og Overgangen danner en fremskridende Forskelsrække, danner
Overgangen
% TRYKFEJL (p. 80): Farveaktens læs Farvesansens
% PRINTED AS IS (p. 44): "variererende" (for varierende; not in the Trykfejl).
indenfor Farveaktens Kvaliteter kun en regelmæssig variererende Forskelsrække.
Og dog kan Analogien gennemføres, idet der til en vis Brydningsgrad (en vis
kvantitativ Forskel i Lysbrydningen) stedse viser sig at svare en bestemt
Farvekvalitet (indenfor Spektret). Ad Analogiens Vej er Farverækken bleven
rationel.

\textsc{Kant} har klart set dette, hvad et enkelt Forhold angaar. Idet han
sammenlignede det rent logisk-matematiske Forhold mellem Grund og Følge, der
slet ikke er noget Tidsforhold, men helt igennem hviler paa Identitet, med de
faktiske Forandringer, vi søger at forklare som Aarsager til eller Virkninger af
hinanden, fandt han kun Analogi imellem dem, en Analogi, han, som vi ovenfor har
set, begrundede som gyldig ud fra Begrebet om Erfarings Mulighed. Mellem formal
og real Videnskab er der ifølge \textsc{Kant} Analogi, ikke Identitet.
Kvaliteternes og Successionernes Verden gør sig stadig gældende paany, selv om
Videnskaben dekreterer, at Kvalitet kan reduceres til Kvantitet og Succession
til Identitet. Mellem de to Rækker, der dannes af de rene Begreber (Tid, Tal,
Grad, Sted) paa den ene Side, de kvalitative Forandringer paa den anden Side,
kan kun paavises ensidig, ikke tvesidig Ækvivalens. Man kan ved at betragte de
kvalitative Forandringers Række, som om den kun frembød Overgang i Tid, Tal,
Grad og Sted, foregribe og forudsige fremtidige kvalitative Forandringer. Men
man kan ikke bevise, at ingen andre Forudsætninger end de, der bestemmes ved
hine fire formale Begreber, kunde føre til saadanne Resultater. Forskellen
mellem et Analogiforhold og et Identitetsforhold er netop det, at ved en vel
valgt Analogi kan et Omraade tjene til intellektuel Belysning af et andet, men
det kan ikke, som hvor der er et Identitetsforhold, godtgøres, at ingen anden
Belysning var mulig.

% --- p. 45 ---
\setcounter{footnote}{0}%
\opage{45}Kausalitet kan kun defineres som en Succession, hvis Led forholder
sig, som Grund forholder sig til Følge. Af de to Elementer i Kausaliteten, det
logiske og det temporale, kan vi ikke ved Kombination danne et nyt, positivt
Begreb. Vi ved egentlig ikke, hvad Aarsag er, hvis vi ikke er tilfreds med den
givne Definition. Et Aarsagsforhold er en Totalitet, som bestaar ved den stadige
Analogi mellem de logiske og de temporale Elementer, en Analogi, der --- som
Arbejdshypotese --- opfordrer til at overvinde de Vanskeligheder for
Forstaaelse, Kvalitetsforandringen medfører, gennem Omsætning til rent logiske
og matematiske Forhold. Saalænge Erkendelsestrangen er vaagen, vil der blive
arbejdet paa ad denne Vej at frembringe reale Sidestykker til de rent logiske og
matematiske Totaliteter. Men om dette er muligt overfor alle foreliggende Emner,
--- om Begreberne „Erfaring“ eller „Natur“, saaledes som \textsc{Kant} havde
konstrueret dem, blev mere end ideale Hypoteser, --- det Spørgsmaal opkastede
han ikke. Og kun derfor kunde han mene at have løst \textsc{Hume}'s Problem.

Det almindelige Princip, som i tidligere Tid kaldtes den tilstrækkelige Grunds
Princip, men som nu for os fremtræder i sin formale Form som Begrundelsens
Princip og i sin reale Form som Aarsagsprincipet, er selv en Følge af
Nødvendigheden af at tænke alle Emner som Led i en Helhed. Der gør sig her en
Trang til Supplering af det enkelte Emne gældende. Allerede \textsc{Telesio},
senere \textsc{Herbart} og \textsc{Hamilton}\footnote{%
% BROKEN SORT? (p. 45 note 1): the last e of „nyere“ prints without its
% crossbar and reads as c („nyerc“); taken as a broken e, not a c sort
% (second reading 2026-09-30, 500 dpi).
\textit{Den nyere
Filosofis Historie}, I, p. 93 f;
% TRYKFEJL (p. 80): 340 læs 388
II, p. 248, 340.} og i nyeste Tid \textsc{Bradley}\footnote{\textit{Essays on
Truth and Reality} (1914), p. 312: Where \textit{A} is not real by itself, but
implies and belongs to an ideal whole, you want a reason for \textit{A}, for you
want to
% PRINTED AS IS (p. 45, note 2): "known" for know (300 dpi); a stray raised dot
% after "want" not transcribed.
known the How of this unity . . . But the „axiom of ground“ is a result and not
a premise . . . The „axiom“ holds only so far as a thing is not complete in
itself . . . The demand for the making good of such imperfection, . . . the
completion of the thing in idea so as to satisfy us theoretically, is what we
mean by the search for „why“ and „how“.} har klart set dette.

\medskip
17. Fra to forskellige Sider kan den her fremsatte Teori kritiseres. Man kan
nægte Berettigelsen og Nødvendigheden af Eftertankens Bearbejdelse af Emnerne,
og man kan hævde, at denne Bearbejdelse fører til en saadan Gennemførelse af
logiske og matematiske Synspunkter, at der ikke er Grund til at tale om blot
Analogi. I første Tilfælde falder Analogien bort, fordi man holder sig til
Emnerne i deres umiddelbare Form, uden Tydning; i andet Tilfælde falder den
bort, fordi Emnerne er helt omsatte til Led i identiske Domme eller Principer.
--- Som Typer for den første Opfattelse kan anføres \textsc{Henri Bergson}, for
den anden \textsc{Hermann Cohen}.

For det første Standpunkt er Tankegangen omtrent følgende. Vi maa gøre os
fortrolige med den Tanke, at hvad der er os umiddelbart givet, er kvalitativ
Forandring, og vi maa holde os til dette umiddelbart Givne, med Bortseen fra
alle Abstraktioner og Reflektioner. Det gælder om at leve sig ind i
Virkeligheden, og den er ikke i vore Tydninger og Tanker, men i det, der
umiddelbart forkynder sig for os. Det store Apparat, som Videnskaben konstruerer
til Forklaring af Virkeligheden, er meget hensigtsmæssigt til at faa Tingene til
at tjene vore praktiske For%
% --- p. 46 ---
\setcounter{footnote}{0}%
\opage{46}maal og til at meddele os til hverandre; men det fører bort fra det
virkelig Givne. Dette Apparat lægger i den Grad Beslag paa os af praktiske
Grunde, at det kan være meget vanskeligt at skyde det tilside og at frigøre sig
fra alle de Abstraktioner og Refleksioner, der dels fører til det, dels
foranlediges ved det. Men lykkes det, falder alle Problemer bort\footnote{Smlgn.
min Bog \textit{Henri Bergson's Filosofi}, p. 16 f; 58---60. --- \textsc{William
James} staar paa dette Punkt \textsc{Bergson} nær. Se \textit{A pluralistic
Universe} (1909), p. 242---247; 260---263.}.

Denne Opfattelse gaar ud fra, at der en Gang for alle kan skelnes mellem Emne og
Eftertanke. Men denne Forskel varierer i det uendelige, baade i Art og Grad. Og
man kan allerede i det umiddelbart Givne spore de Love og Former, efter hvilke
Eftertanken arbejder. Totalitetspræget er allerede tilstede i det, og
Eftertanken udfolder kun videre, hvad der allerede er antydet i det. (Smlgn.
ovenfor § 10). En saadan videre Udfoldelse er nødvendig paa Grund af de
Modsigelser og Modsætninger, allerede det umiddelbart Givne kan indeholde. Der
gives jo ikke et enkelt umiddelbart Givet, en enkelt konkret eller praktisk
Intuition, men mange forskellige Intuitioner, hver for sig med et Helhedspræg
--- og hvorledes kan disse forskellige Intuitioner forenes til en mere
omfattende Intuition? Uden Eftertanke og Sammenligning kan det ikke ske.

I min Kritik af \textsc{Bergson} i „Henri Bergson's Filosofi“ holdt jeg mig til
Forholdet mellem Intuition og Eftertanke i Almindelighed. Men ogsaa Forholdet
mellem de mange forskellige Intuitioner indbyrdes fører ud over det umiddelbart
Givne, selv om man indrømmer, at vi kan naa tilbage til dette.

Fremdeles trænger sig nødvendigvis det Spørgsmaal frem, med hvilken Ret man
kalder det umiddelbart Givne for „Virkelighed“. Forskellen mellem Fantasi, eller
Drøm, og Virkelighed existerer ikke for konkrete og praktiske Intuitioner. Det
Problem, denne Forskel stiller, existerer der ligesaa lidt som noget andet
Problem. En begrundet Forskel mellem Virkelighed og Drøm forudsætter Erkendelse
af fast, lovmæssig Sammenhæng mellem de Iagttagelser, der staar til Raadighed,
--- selv om denne Erkendelse er nok saa elementær og barnlig. Og det er nu netop
Trangen til Kendemærker for Virkelighed, der fører til Loverkendelse.
\textsc{Henri Poincaré} havde derfor Ret i den Sætning: „Le monde Bergsonien n'a
pas des lois“. (Dernières
% DEFECT (p. 46): a stray raised tick after the comma in "Pensées, p. 31" (the
% layer reads "/"); not transcribed.
Pensées, p. 31).

I skarp Modsætning til \textsc{Bergson}'s Intuitionsteori fremtræder
Marburgerskolens Teori om, at Kausalitetsproblemet er fuldstændig løst, dersom
det lykkes at gennemarbejde de Forskelligheder, det Givne frembyder, saaledes at
alle Forhold indenfor det Givne kan udtrykkes i identiske Domme eller Ligninger.
Det matematiske Funktionsbegreb indeholder da Løsningen af Kausalitetsproblemet.
En matematisk Funktion betegner gensidig Afhængighed mellem to foranderlige
Størrelser og udtrykker saaledes paa én Gang deres Forskellighed og deres
Sammenhæng. Naar $y = f(x)$, falder $x$ ikke bort, fordi $y$ indtræder, men
bevares i $y$. Naar et Forhold mellem to foranderlige Fænomener kan udtrykkes
paa denne Maade, falder det Udvortes i den populære Opfattelse af
Aarsagsforholdet bort, og der bliver ikke mere
% --- p. 47 ---
\setcounter{footnote}{0}%
\opage{47}Tale om en ydre Indgriben eller Handling af en Ting paa en anden.
Tingenes Væsen og Forhold har en Gang for alle fundet deres Udtryk i
Ligningen\footnote{H. \textsc{Cohen}: \textit{Logik der reinen Erkenntniss}, p.
185; 239---247. --- Til et lignende Resultat kommer \textsc{Richard Avenarius}
ud fra sit pragmatiske Standpunkt i sin Teori om absolut „Deproblematisering“.
(Se „Moderne Filosofer“, p. 103---104).}.

% PRINTED AS IS (p. 47): "Tydning;" followed by capital "Ligningen" (checked at
% 300 dpi; UNSURE: rejected a full stop with a stray mark above it).
Imidlertid skelner \textsc{Cohen} dog (p. 248) mellem Ligningen og dens Tydning;
Ligningen danner kun det metodologiske Grundlag. Den tydes som Udtryk for
Konstans, Bestaaen (Erhaltung). Men Ligningen $y = f(x)$ kan betyde saare mange
Ting, finde Anvendelse paa højst forskellige Omraader. Naar den skal være det
metodologiske Grundlag for „Bestaaen“ eller „fysisk Struktur“, saa har vi her
netop hvad \textsc{Kant} kaldte Analogi. Der er kun Analogi mellem matematiske
og fysiske Forhold, ligesom der kun er Analogi mellem logisk Begrundelse og
kausal Relation. Det er paa denne Analogis Berettigelse, den matematiske Fysik
støtter sig, eller rettere: denne Analogis Berettigelse er den store Hypotese,
al moderne Naturvidenskab søger at godtgøre i det Enkelte.

\textsc{Cohen} anerkender fremdeles (p. 412 ff), at Problemer kun opstaar, hvor
et nyt Emne bliver til for os. Det gælder da at omsætte dette nye kvalitative
Indhold i kvantitativ Form. Sansefornemmelsen stiller Krav, men kun Tanken kan
afgøre Kravets Betydning. Sansefornemmelsen skal ikke høres, men forhøres. Men
Resultatet af disse Forhør kan ikke fastslaas en Gang for alle. Vore Begreber
kan derfor aldrig blive fuldkomne. Herved er indrømmet saa meget, at
Gennemførelsen af den rent logiske Opfattelse af Virkeligheden stadig staar som
et Ideal eller en Hypotese. Men selv hvor Tydningen maa siges at være lykkedes,
som naar Varmelæren er bleven en Del af Bevægelseslæren, ophæves saa derved
Forskellen mellem Fornemmelseskvaliteten og de matematiske Formler, der
udtrykker Varmen som fysisk Emne? Fornemmelsens Krykke, siger \textsc{Cohen} (p.
425), er knust derved, at Varmelæren er en Del af Bevægelseslæren. Men
Fornemmelsen behøver i sin umiddelbare Kvalitet slet ingen Krykke, og selv efter
at Varmelæren er bleven en Del af Bevægelseslæren, staar det Spørgsmaal uløst:
hvorfor og hvorledes gør det Fænomen, den fysiske Varmelære beskriver ved Hjælp
af sine Ligninger, sig gældende for sansende Væsener paa en bestemt, kvalitativ
Maade? Kan man af den fysiske Varmelæres Ligninger udlede de Varmekvaliteter,
som er umiddelbart givne for os? Man kan erklære disse Kvaliteter for
subjektive, men derved skaffes de ikke ud af Verden. Hvad der gælder for
Temperatursansen, gælder ogsaa for alle andre Sanser. Og det er ikke blot
Forholdet til menneskelige Sanser, der stiller Problemet paa ny. Trods Sætningen
om Energiens Bestaaen er Lys, Varme, mekanisk Bevægelse, Elektricitet o.~s.~v.
dog forskellige Naturkræfter, og den objektive Forskel mellem dem kan ikke
udledes af den almindelige Sætning, at der til et vist Kvantum af den ene svarer
et vist Kvantum af den anden. Den ene har faktisk andre Virkninger end den
anden, og i hvert enkelt Tilfælde foreligger der en Overgang i en bestemt
Retning. Selv om Sætningen $A = B$ gælder for et vist Synspunkt, kan det dog i
det enkelte Tilfælde være af afgørende Betydning, om det er $A$ eller $B$,
% --- p. 48 ---
\setcounter{footnote}{0}%
\opage{48}der foreligger. Spørgsmaalet bliver tilsidst, om Virkeligheden kan
tænkes at bestaa i de blotte Ligninger. \textsc{Cohen} er ikke langt fra denne
Konsekvens. Kun Tanken, Begrebet, siger han, giver sand Væren (p. 180). Al sand
Filosofi begynder efter ham med \textsc{Parmenides}. Dog synes han ikke at ville
drage den Konsekvens, at kun de rene Identiteters Verden existerer.

Endnu mindre end \textsc{Kant}'s Analogilære formaar \textsc{Cohen}'s
Identitetslære at løse \textsc{Hume}'s Problem. Men \textsc{Cohen} anerkender
slet ikke Berettigelsen af dette Problem, som jo lægger Vægten paa Kvalitets- og
Tidsforskellighederne. Han afviser (p. 230) \textsc{Hume}'s hele Standpunkt som
umoralsk! Det har saa ogsaa været umoralsk af \textsc{Kant} at lade sig vække af
dogmatisk Slummer ved \textsc{Hume}'s Problem. Og konsekvent er det da, at man
rolig lægger sig til at slumre igen!

\medskip
18. Aarsagsbegrebet fremtræder ifølge den Opfattelse, jeg her har gjort
gældende, som en Udvidelse af den Tankehelhed, der kan naas ad logisk-matematisk
Vej, en Udvidelse, der sker ved, at kvalitative og successive Forskelligheder ad
Analogiens Vej bringes ind under logiske og matematiske Synspunkter. Der dannes
da en real Tankehelhed, som giver den mest fuldkomne Virkelighedserkendelse, der
kan vindes, uden at dog Virkeligheden paa eleatisk-platonisk Vis identificeres
med Logik og Matematik. At det kun er ad Tænkningens Vej, vi kan begrunde vor
Tro paa at staa overfor en Virkelighed, kan ikke fortolkes saaledes, at denne
Virkelighed kun skulde bestaa i selve Tænkningens Former. Forholdet mellem
Iagttagelse og Tydning bliver aldrig Identitet. \textsc{Hume}'s Problem bliver
staaende\footnote{I \textsc{Anton Thomsen}'s Bog om \textsc{Hume} gjordes der
Forsøg paa at vise, at Æren for at have opstillet Aarsagsproblemet ikke
tilkommer \textsc{Hume}, men \textsc{Hobbes}. Det er dog kun et enkelt Sted hos
\textsc{Hobbes}, man her kan beraabe sig paa (De Corpore I, 3, 20). Det var
fortjenstfuldt at drage dette Sted frem, men hvad \textsc{Hobbes} her har
udtalt, faar ingen Indflydelse paa hans dogmatiske Tankegang. Og ogsaa før
\textsc{Hobbes} er der Antydninger af Forskellen mellem Grund og Aarsag.
Saaledes hos \textsc{Bruno} (De la causa, principio, et uno. Secondo dialogo:
Forskellen mellem principio og causa. Ed. \textsc{Lagarde}, p. 230), og endnu
tydeligere hos det 14. Aarhundredes Nominalister, især hos \textsc{Nicolas
d'Autrecourt}. (Se \textsc{Lappe}: Die Philosophie des Nicolaus von Autrecourt.
Bonn 1905. --- H. \textsc{Rashdall}: Nicholas de Ultricuria, a medieval Hume.
Aristotelian Society. Proceedings 1906).}.

Aarsag og Virkning, der foreløbig staar som to forskellige Ting, kommer i
inderligere Forhold til hinanden som Led i en real Tankehelhed, der er dannet
med den formale Tankehelhed som Forbillede, og i Kraft af den Helhedstendens,
der ytrer sig i alt Bevidsthedsliv, især i alt Tankeliv. Men ogsaa fra andre
Sider lægger Helhedsbegrebet eller Helhedstendensen sig for Dagen i
Aarsagsforholdet, og det kan tjene til Belysning af vort Emne at fremdrage dem.

% ENUMERATOR (p. 48): italic Latin "a.", decided by the series b. (p. 49), c.
% (p. 50).
\textit{a.} Den Helhed, som Aarsag og Virkning udgør, er selv Led i en større
Helhed. Ethvert enkelt, isoleret Aarsagsforhold (om der gives et saadant) er
blindt og tilfældigt; Betingelsen for, at to Begivenheder paa dette bestemte
Sted og til denne bestemte Tid kan staa i Aarsagsforhold, maa søges i en større
Sammenhæng, tilsidst i den Lovmæssighed, der ligger til Grund for hele den
Natursammenhæng, vi kender. Ved at se to Begivenheder i Aarsagssammenhæng ser
man dem tillige som indbefat%
% --- p. 49 ---
\setcounter{footnote}{0}%
\opage{49}tede i en Helhed, der er mere omfattende end den, de selv tilsammen
udgør. Dette træder især frem, naar man lægger Vægt paa Aarsagsforholdets
Retning, en Side ved det, man er tilbøjelig til at overse. \textsc{Leibniz}
henledede bestemt Opmærksomheden paa Retningen og opstillede Sætningen om
Retningens Bestaaen\footnote{\textit{Opera philosophica}, ed. \textsc{Erdmann},
p. 132 f.}. Senere er Retningsbegrebet traadt i Baggrunden, skønt \textsc{Kant}
lagde meget stor Vægt paa, at Aarsagsforholdet er forskelligt fra et rent og
purt Successionsforhold derved, at Leddene ikke kan byttes om. Hvor der er
tvesidig Ækvivalens mellem Aarsag og Virkning, kan man baade experimentalt og
tankemæssigt gaa frem og tilbage. Men i virkelig Iagttagelse, i det enkelte
bestemte Tilfælde gaas der enten frem eller tilbage; der er Ækvivalensen altid
ensidig, og det er af afgørende Betydning, hvilket af Leddene der er Aarsag, og
hvilket der er Virkning. Opdagelsen af Naturkræfternes Ækvivalens har
ingenlunde, som \textsc{Avenarius}, \textsc{Riehl} og \textsc{Cohen} mener, løst
\textsc{Hume}'s Problem\footnote{Smlgn. \textit{Den menneskelige Tanke}, p.
282---287.}. --- I det lange Løb er maaske endda Ækvivalensen faktisk ensidig;
% PRINTED AS IS (p. 49): the small caps run on through "Lov" (checked at 300 dpi).
det er idetmindste den Slutning, nogle drager af \textsc{Carnot}'s \textsc{Lov}.

Dersom det kun er paa Baggrund af en mere omfattende Helhed, at Retningen kan
konstateres, og derved det bestemte, konkrete Aarsagsforhold ses i sin
Ejendommelighed, --- og dersom Virkelighedskriteriet bestaar i den kausale
Sammenhæng, den bestemte Orden og Retning, i hvilken de enkelte Emner fremstaar,
--- saa bliver „Virkeligheden“, den reale Sandhed, et Ideal, der til hver given
Tid kun tilnærmelsesvis kan naas. Dette fører dog ikke, som man har
ment\footnote{Smlgn. \textsc{Bertrand Russell}'s Kritik af det dynamiske
Sandhedsbegreb i hans Afhandling \textit{On the Nature of Truth} (Proceedings of
the Aristotelian Society, 1907, p. 28---49). Selv hylder \textsc{Russell} en
„analytisk Realisme“, der antager en Verden af absolute Existenser, dels
universelle (i Lighed med \textsc{Platon}'s Ideer), dels sanselige (\textit{Le
Réalisme analytique}. Bulletin de la Société Française de Philosophie, 1911).
Kun derved mener han at have fast Grund under Fødderne.}, til fuldkommen
Skepsis; det er kun et Vidnesbyrd, det mest afgørende af alle Vidnesbyrd, om, at
Sandheden udvikler sig historisk.

\textit{b.} Ligesom Grunden altid er et Indbegreb af Forudsætninger, saaledes er
Aarsagen altid et Indbegreb af positive og negative Betingelser. Da nu enhver af
disse Bestanddele af „Aarsagen“ har en Tendens til at virke i en bestemt
Retning, og da Forholdet mellem disse Tendenser kan være højst forskelligt,
bliver der Mulighed for meget forskellige Ordninger og Kombinationer indenfor
det, der betegnes som „Aarsagen“. Aarsagen som Helhed kan derfor fremtræde paa
højst forskellig Maade, og det vil være vanskeligt, om ikke umuligt, at
formulere en Lov, der udtømmer de virkelige Forhold. „En Naturlov kan“, siger
\textsc{Ernst Mach}, „ikke indeholde mere end en omfattende og fortættet
Beretning om Kendsgerninger. Ja, den indeholder endog stedse mindre end selve
Kendsgerningen, fordi den ikke gengiver hele Kendsgerningen, men kun den for os
vigtige Side af den, medens der med Hensigt eller nødtvungent ses bort fra
Fuldstændighed“\footnote{\textit{Oekonomische Natur der physikalischen
Forschung}(Populär-wissenschaftliche Vorlesungen\textsuperscript{4}) p. 224 f.}.
\textsc{Bergson} gaar endog saa vidt at kalde Aarsagslovene for blotte
Etiketter. Dette Misforhold mellem Aarsagslov og
% SIGNATURE LINE at the foot of p. 49 ("D. K. D. Vidensk. Selsk. Skr., 7.
% Række, hist. og filos. Afd. III. 2." and "39") omitted as printer's
% apparatus.
% --- p. 50 ---
\setcounter{footnote}{0}%
\opage{50}virkelig Skeen fremtræder ikke mindst paa det aandelige Omraade. En
Villiesvirksomhed, hvis Formaal omfattes med fuld Bevidsthed, udfolder sig ikke,
uden at en Mængde uvilkaarlige Tendenser af spontan, reflex og instinktiv Art
virker med, foruden at Forestillinger og Følelser, der ikke direkte vedkommer
Formaalet, gør sig gældende. Og Forholdet mellem det bevidste Formaal og de
halvt eller helt ubevidste Tendenser kan i de enkelte Tilfælde variere i høj
Grad, saa at det intet Under er, om Aarsagsforklaringen af en menneskelig
Handling ofte er saa overordentlig vanskelig\footnote{Smlgn. min Afhandling
\textit{Begrebet Villie} (Mindre Arbejder. Tredie Række), p. 40 f.}.

\textit{c.} Hvad der saaledes gælder indenfor den enkelte „Aarsag“, fremtræder
paa analog Maade, naar vi betragter forskellige Aarsagsrækker i deres indbyrdes
Sammenhæng. Der kan ikke paavises nogen Aarsagsrække, som forløber aldeles
isoleret. Forskellige Aarsagsrækker griber ind i hverandre, og hvad der
foreligger i de enkelte Tilfælde, vil stedse være et Komplex af Aarsagsrækker,
hvis Retninger mødes paa harmonisk eller disharmonisk Maade. ---

Der stiller sig altsaa uafbrudt for Forskningen den Opgave \textit{a}) at finde
den store Sammenhæng, indenfor hvilken det enkelte Aarsagsforhold kan finde sin
Begrundelse, --- \textit{b}) at bestemme Forholdet mellem Elementerne indenfor
samme Led i den enkelte Aarsagsrække, --- og \textit{c}) at udrede Forholdet
mellem de forskellige Aarsagsrækker, der tilsammen bestemmer et Resultat.

Real Helhed er derfor (som iøvrigt alle Kategorier er det) en „Idé“, i kantisk
Betydning af Ordet, det vil sige en Tanke, som stedse stiller nye Opgaver og
ikke kan realiseres i nogen Anskuen eller i nogen afsluttet Tankegang. Det er en
Arbejdstanke, der kun løser foreliggende Opgaver saaledes, at den tillige gør
opmærksom paa nye Opgaver.

\medskip
19. Dersom det er rigtigt, hvad der i det Foregaaende er søgt begrundet, at den
menneskelige Tanke i Kraft af sin egen Natur søger at forme Helheder, fordi den
selv uvilkaarlig arbejder efter en Helhedslov, saa vilde det være en lykkelig
Ting, om der gaves Emner, der fra første Færd, forud for al tankemæssig
Bearbejdelse, fremtraadte med Helhedens Præg, og i hvilke det var en Helhedslov,
der bestemte de enkelte Elementers Vekselvirkning. Det var da her ikke
Tankearbejdet, der frembragte Helheden, men det var Emnet selv, der mødte frem
som en given Helhed, der havde Loven skreven i sig selv og ikke behøvede at
hente den fra Tanken. Nu er, som vi ovenfor har set, alle Emner givne som
Helheder; men det Tankearbejde, vi hidtil har undersøgt, gik dog væsentlig ud
paa at opløse det Givne i dets Elementer og vinde Forstaaelsen gennem et
Tankehele, der tilvejebragtes gennem Analogi med den logisk-matematiske
Totalitet, Tanken paa egen Haand kan frembringe. Saadanne Emner optræder paa det
organiske, det psykiske og det sociale Omraade, og Spørgsmaalet bliver da om
Forskningens Forhold til dem. Er her helt nye Kategorier nødvendige, og i saa
Fald, i hvilket Forhold staar de til de Kategorier, der i det Foregaaende er
omtalte?

% --- p. 51 ---
% MARGIN MARKS (p. 51): small arrow-like marks in the left margin of the
% scanned copy; not text of the work, not transcribed.
\opage{51}Naar Totalitet, saaledes som det er sket i „Den menneskelige Tanke“ (p. 218---227),
opstilles som en speciel Kategori, som et enkelt Led i Kategoriernes Række,
saa er det dog ikke blot paa Grund af Biologiens, Psykologiens og Sociologiens
Stilling indenfor Videnskaben, men ogsaa paa Grund af den Sammenhæng mellem
Aarsagsbegrebet og Helhedsbegrebet, der ovenfor er paavist, en Sammenhæng, der
ikke blot fremtræder paa de nævnte tre Omraader. Overalt hvor Forskningen rykker
frem, opdager den Helheder, der hidtil havde skjult sig. Ellers vilde jo ogsaa
Tankens egen Helhedskarakter, dens Ejendommelighed som Syntese, være et højst
uhensigtsmæssigt Organ. Men det har vist sig at være en ren Abstraktion, naar
man opfatter Aarsag som noget Enkelt, Aarsag og Virkning som to forskellige Ting,
og Aarsagsrækkerne som liggende aldeles udenfor hverandre. Naar Helhedsbegrebet
særlig er karakteristisk for de tre nævnte Videnskaber, saa kommer det dog ikke
her frem som en deus ex machina.

Dog gør det en betydelig Forskel, om Helheden kun viser sig som et Maal,
der naas, naar Forskningen kan finde Sammenhæng mellem forskellige Aarsagsrækker,
eller fremtræder som et foreliggende Emne, der saa stiller Forskningen den
Opgave at finde dets indre og ydre Betingelser. I sidste Tilfælde viser det sig i
ganske speciel Forstand, hvad der mere eller mindre gælder overfor ethvert Emne,
at Forskningen gaar analytisk frem. Selv om et Emne i nok saa høj Grad
fremtræder som en Helhed, er Forstaaelsen dog ikke given derved. Der maa analyseres,
skelnes mellem Dele og Egenskaber, for at Loven for deres Sammenhæng kan findes.
De faktisk foreliggende Helheder er ligesaa mange Opgaver, først for Beskrivelse
og Klassificering, saa for Undersøgelse med Hensyn til Betingelserne for deres
Bestaaen og Udvikling. Selve Begrebet Helhed kan aldrig bruges til Forklaring.
Det er, som enhver Kategori, baade en Form for Tankens Arbejde og et Ideal, der
søges udført i det Enkelte. Men ved de umiddelbart givne Helheder er der, som
Erfaringen lige til den nyeste Tid viser, en Fristelse til at finde Forklaringen i selve
Helhedsbegrebet, bruge dette som en Nøgle, der lukker op overalt. Alle teologiske,
vitalistiske og spiritualistiske Forklaringer beror paa et saadant Misbrug. I
Gudsbegrebet, saaledes som det er udviklet i de højere Folkereligioner, er al Kausalitet
i Verden koncentreret som i en uendelig Enhed eller Helhed, og det staar da som
en Selvfølge, at ethvert enkelt Fænomen har sin Forklaring ved det. \textsc{Galilei} gjorde
allerede opmærksom paa, at netop fordi Alt har sin Grund i Gud, bringer det ingen
Forstaaelse at beraabe sig paa Gud som Aarsag til et enkelt Fænomen i dets
Ejendommelighed til given Tid og Sted. Vitalismen betragter Organismen som et Væsen,
der har sin Aarsag i sig selv, og dette er den egentlige Aarsag til alt, hvad der
foregaar indenfor den. Paa det psykologiske Omraade svarer hertil, at man lader
Sjælen være Aarsag til sine egne Tilstande, istedetfor at finde Sjælelivets Love
% PRINTED AS IS (p. 51): „Forboldet“ for „Forholdet“ (b for h; second
% reading 2026-09-30, clear at 500 dpi).
gennem Studiet af Forboldet mellem de forskellige sjælelige Ytringer. Paa det
sociologiske Omraade fremtræder en lignende Tendens, naar man lader Samfundet som
Helhed være Aarsag til specielle sociale Fænomener.

I \textsc{Leibniz}' Individualitetsbegreb ligger en Anvisning til paa en Gang at aner%
% SIGNATURE MARK "39*" at the foot of p. 51 omitted as printer's apparatus.
% --- p. 52 ---
\setcounter{footnote}{0}%
\opage{52}kende Ejendommeligheden ved en given Helhed og at finde dens Væsen i den Lov,
der erfaringsmæssig udtrykker Sammenhængen mellem dens forskellige Tilstande
og Forandringer\footnote{Smlgn. \textit{Den menneskelige Tanke}, p. 131---135.}.
Som vi ovenfor har set, var det for \textsc{Leibniz} Erkendelsens Ideal
at føre alle Domme tilbage til Identiteter, og det kunde synes at være i Modstrid
hermed, at han lægger saa stor Vægt paa individuelle Helheder og deres
Forskelligheder. Men Sagen er, at Identiteten, der for \textsc{Platon} stod som Genstand for
Kontemplation, en Kontemplation, der nærmede sig til Mystik, for \textsc{Leibniz} var en
Tankeform og et Tankemiddel, der benyttes ved hvert enkelt lille Skridt, hver lille
Overgang, Tanken gør i den Analyse, der fører til at finde Loven for den individuelle
Helhed (cette loi de l'ordre qui fait l'individualité de chaque substance
particulière; --- lex continuationis operationum). Denne Tanke hos \textsc{Leibniz} hænger nøje
sammen med Grundtanken i hans Infinitesimalregning, idet det er gennem uendelig smaa
Overgange mellem Tilstandene, at et Væsens Kontinuitet udtrykker sig, og idet et Væsens
Ejendommelighed netop bestaar i den Lov, efter hvilken disse Overgange foregaar.

\textsc{Leibniz}'s Tanke kan lede os i vor Betragtning af Organismen, Personligheden
og Samfundet fra et erkendelsesteoretisk Synspunkt.

\medskip
20. Paa det biologiske Omraade staar Mekanicisme og Vitalisme overfor
hinanden.

% SMALL CAPS (p. 52): the pronoun "Hin" (= Mekanicismen) opens the paragraph in
% small caps, as the names are set (checked at 300 dpi).
\textsc{Hin} ser i de organiske Fænomener blot en Række fysiske og kemiske Processer,
hvis Love i Grunden er fundne, saa at det blot gælder Anvendelse af dem
paa specielle Tilfælde. Organismen er en Sum af mekanisk (fysisk-kemisk) virkende
Elementer. Den videnskabelige Biologi begynder med \textsc{Lavoisier}'s og \textsc{Laplace}'s
Opdagelse (1780) af, at dyrisk Varme kan forklares ad fysisk-kemisk Vej. Siden er en
saadan Forklaring gennemført Skridt for Skridt for forskellige Funktioners
Vedkommende, bl. A. for Forplantningens. Der er endnu Huller tilbage, saaledes den
kemiske Beskaffenhed af Enzymerne og af de Stoffer, der betinger Arveligheden af
en Egenskab. Men et særegent Livsprincip er overflødigt, og en Organisme kan
defineres som en kemisk virkende Maskine, der besidder den Ejendommelighed at
kunne udvikle, opholde og reproducere sig automatisk. At en Del er saaledes
tilpasset, at den tjener Helheden, er for den mekanistiske Opfattelse kun et uklart
Udtryk for den Kendsgerning, at en Art kun kan leve, naar den besidder den
Mekanisme, Selvopholdelse og Reproduktion forudsætter. Overordentlig mange
Former er gaaede til Grunde af Mangel paa en saadan Mekanisme. Der er Intet, der
tyder paa, at kunstig Frembringelse af levende Materie skulde ligge ude over
Videnskabens Muligheder\footnote{%
% PRINTED AS IS (p. 52, note 2): "developement" and "Congres" (checked at 300 dpi).
\textsc{Jacques Loeb}: \textit{The recent
developement of Biology}. (Congres of Arts and Sciences. St. Louis 1904.
Vol. V.). --- \textit{The Mechanistic Conception of Life}. Chicago 1912. ---
\textsc{Loeb} anerkender Betydningen af \textsc{Arrhenius}' Tanke, at meget smaa
Spirer drives gennem Rummet ved Straaletryk, og at de, naar de støder paa
Legemer af passende Temperatur og forsynede med Vand, Salte og Ilt, fremkalder
nye organiske Udviklinger. Men han mener, at Videnskaben alligevel maa arbejde
paa at frembringe levende Materie, eller i hvert Tilfælde bevise, hvorfor det er
umuligt.}.

% --- p. 53 ---
\setcounter{footnote}{0}%
\opage{53}Selv om der lægges nok saa megen Vægt paa den strenge Gennemførelse af
fysiske og kemiske Metoder indenfor Biologien, maa det dog, som selv \textsc{Loeb}
indrømmer, foreløbig hævdes, at ikke Alt paa det organiske Omraade er forklaret ad
denne Vej. Paa dette Punkt tager derfor ganske naturlig den modstaaende
Opfattelse, Vitalismen, fat. Ved Vitalisme maa man, hvis man vil have et klart
Begreb og tillige tager Hensyn til Biologiens Historie, forstaa den Opfattelse, der ikke
blot nægter, at det organiske Liv kan forklares ad fysisk og kemisk Vej, men tillige
finder Forklaringen i en særlig Livskraft, der er forskellig fra alle andre
Naturkræfter\footnote{Se min Afhandling \textit{Om Vitalisme} (Mindre Arbejder.
Første Række), p. 44---46. De klassiske Vitalister er \textsc{van Helmont} og
\textsc{Stahl} i d. 17. Aarh., Montpellierskolen, især \textsc{Barthez} i d. 18.
Aarh. (Med Urette er \textsc{Verworn} p. 46 nævnt som Vitalist). --- Jeg forstaar
ikke, hvorledes \textsc{Driesch} (\textit{History and Theory of Vitalism}. 1914.
p. 59 f) kan betragte \textsc{Blumenbach} som den mest karakteristiske
Repræsentant for den gamle Vitalisme. \textsc{Blumenbach} betragter, som
\textsc{Driesch} selv fremhæver, den formende Tendens (nisus formativus) ikke som
Aarsag, men som et Resultat, hvis Aarsag ikke kendes. \textsc{Blumenbach}
(\textit{Über den Bildungstrieb}. 1789, p. 25 f) siger udtrykkelig, at Ordet
„Bildungstrieb“ ligesaa lidt indeholder nogen Forklaring, som Ordet „Tiltrækning“
hos \textsc{Newton} gør det, og han minder om \textsc{Newton}'s udtrykkelige
Erklæring i den Henseende. Disse Ord betegner Fænomenet, der skal forklares, ikke
Aarsagen. Naar \textsc{Driesch} derfor siger, at hvis alle Vitalister havde
udtrykt sig ligesaa klart som \textsc{Blumenbach}, vilde \textsc{Lotze}'s og
\textsc{Claude Bernard}'s Kritik have været unødvendig, saa er dette ganske
rigtigt; ti der vilde da ingen Vitalisme have existeret!}.

Vitalisme er efter denne Definition baade et negativt og et positivt Begreb.
Den beror først og fremmest paa den Forudsætning, at en mekanisk Oprindelse af
Livet ikke er mulig. Men er det muligt at føre et Bevis herfor? Naar Biologiens
Fremskridt, hvad Aarsagsforklaring angaar, stedse har bestaaet i, at der paavistes
bestemte fysiske og kemiske Aarsager til organiske Processer, saa vil det ligge
nærmest at formode, at denne Fremskriden kan fortsættes, saalænge der overhovedet
forskes, selv om en Afslutning her ikke skulde kunne naas. Vi staar maaske her
ved et irrationelt Forhold, ligesom ved Forholdet mellem Sansekvaliteten og den
matematiske Fysiks Ligninger. (Smlgn. § 11 og § 15). Al Videnskab arbejder egentlig
paa Analyse af givne Helheder. Et Hul i de fysiske og kemiske Processers Forløb
indenfor Organismen, betegner kun en ny Opgave, Nødvendigheden af flere
Decimaler. Saaledes fastholder ogsaa \textsc{Hugo de Vries}, at Mutationsteorien, der antager
pludselig Tilblivelse af nye organiske Typer, netop stiller den Opgave at udfinde
den kemiske Aarsag til denne Afvigelse fra den forhen raadende Type.

Positive Beviser for Vitalismen mener \textsc{Driesch}, den Forsker i vore Dage, der
udførligst har hævdet Vitalismen, at have fundet, dels i Embryologien, dels i
Instinktet. Morfologien er Vitalismens egentlige Omraade. De organiske Former med
deres Helhedspræg synes ikke at kunne udledes som Resultat af blot Sammenspil
af uorganiske Elementer. \textsc{Driesch} henviser, hvad Embryologien angaar, til, at i
de tidligste Stadier af Fosterudviklingen kan Cellernes indbyrdes Ordning helt
forandres ved mekanisk Indgreb, og dog kan Cellekomplexet trods Splittelsen udvikle
sig paa aldeles normal Maade. Men, som
\textsc{Haldane}\footnote{\textit{Mechanism, Life and Personality}. (1913), p. 27.}
bemærker, kan der fra et
% --- p. 54 ---
\setcounter{footnote}{0}%
\opage{54}fysiologisk Synspunkt ikke være Tvivl om, at der herved virker specielle fysiske
og kemiske Aarsager, selv om de endnu ikke er paaviste. Og i hvert Tilfælde er
det tilsyneladende positive Bevis kun en speciel Form for det negative, der har vist
sig umuligt at føre. --- Ligesom i Formdannelsen skal nu ifølge \textsc{Driesch}
Organismen virke som Helhed i Instinkthandlingerne, lede og bestemme Processernes
Retning og Sammenhæng, snart ved at hæmme en enkelt Proces, snart ved at ændre
Retningen, snart ved at omordne Forholdet mellem de enkelte Processer. En saadan
hæmmende, ændrende og ordnende Indgriben strider ifølge \textsc{Driesch} (der her
benytter en allerede af \textsc{Descartes} til Fordel for Sjælens Indgriben i Materien udtalt
Tanke) ikke mod Sætningen om Energiens Bestaaen. Men heller ikke herved
kommer vi ud over en negativ Paastand. En Retningsændring uden materiel Aarsag
vilde ganske vist ikke stride mod Energiens Bestaaen, men den vilde stride mod
den strengt fysiske Opfattelse af Inertisætningen.

\textsc{Driesch}'s Behandling af Forholdet mellem det Mekaniske og det Organiske er
i denne Sammenhæng af særlig Interesse, fordi han som Logiker lægger stærk Vægt
paa Helhedens Kategori\footnote{%
% PRINTED AS IS (p. 54, note 1): "af" for "of" in the English title, which is set
% roman, not italic (checked at 300 dpi).
\textsc{Driesch}: The Problem af Individuality.
(1914), p. 41.}. Men det er i det foregaaende forhaabentlig
tilstrækkelig paavist, at man af Tankens almene Former ikke kan udlede specielle
Naturforklaringer. Totalitet som Kategori fører til et bevidst og ubevidst Arbejde paa at
forme Iagttagelser til Helheder. Men hvor, som paa det organiske Omraade,
Fænomener allerede for selve Iagttagelsen staar som Helheder, der maa Opgaven blive
den, ad Analysens Vej at finde de Love, der betinger saadanne Helheders Bestaaen,
ikke den at appellere til Helheden selv som selvfrembringende Aarsag. Paa dette
% PRINTED AS IS (p. 54): "J. C. Haldane" (the physiologist is J. S. Haldane), checked at 300 dpi.
Punkt staar en anden allerede omtalt filosoferende Biolog J. C. \textsc{Haldane} mere
kritisk. Han anvender ikke Helhedsbegrebet kausalt, men ser de Helheder, der
fremtræder i Naturen, som Typer paa en højere Form for Realitet (a higher or more
concrete aspect of reality). Derimod mener han, at hvor en saadan „højere“ Type
fremtræder, stilles der overhovedet ikke noget kausalt Problem: den er som en
Aabenbaring af Tilværelsens Væsen, som ikke afbryder den mekaniske Sammenhæng
og ikke berøver det fysisk-kemiske Verdensbegreb dets Nytte som
Arbejdshypotese\footnote{\textit{Mechanism, Life and Personality}, p. 98, 137.
--- For \textsc{Haldane} er det biologiske Synspunkt „højere“ end det
fysisk-kemiske, det psykologiske igen „højere“ end det biologiske. Materie,
Organisme og Sjæl betegner forskellige Trin af Tilnærmelse til sand Virkelighed
(p. 117, 123).}. Som saa mange Hegelianere ser \textsc{Haldane} ikke, at det i Virkeligheden
er et Værdibegreb, han her opererer med, paa samme Maade, som naar man f. Ex.
siger, at et Kunstværks æstetiske Værdi er uafhængig af den Maade, paa hvilken
det er blevet til. I næste Kapitel kommer vi til at drøfte Forholdet mellem Totalitet
og Værdi.

Naar det saaledes viser sig, at det positive Bevis for Vitalismen kun er en
Form for det negative, spørger man med en vis Nysgerrighed, hvorledes man
egentlig skulde tænke sig den nye Kraft, der ifølge Vitalismen træder frem i de organiske
Fænomener for at udfylde de Huller, som den mekaniske Naturforklaring stadig
% STRAY MARKS (p. 54): small dots and ticks in the right margin and below the notes
% (the layer reads "*", "'", "„"); marks in the scanned copy, not transcribed.
% --- p. 55 ---
\setcounter{footnote}{0}%
\opage{55}skulde efterlade. \textsc{Driesch} oplyser herom, at den ikke er en rumlig Aarsag, skønt
den virker paa (into) det, som er i Rummet. Hvis man ved Aarsag vil forstaa rumlig
Aarsag, saa er den her virkende ordnende Faktor, som \textsc{Driesch} vil kalde Entelechi,
ikke Aarsag. Hvor den lægger sig for Dagen som ledende, ikke blot formende,
kalder \textsc{Driesch} den et Psychoïd. Psychoïdet er ikke af psykisk Art, men er et
Analogon til det Psykiske, og \textsc{Driesch} finder Vitalismens Hovedvanskelighed i denne
Antagelse af ubevidste Analoga til Sjælelivet. Psychoïdet skal dog ikke blot være
nødvendigt for at forstaa visse biologiske Fænomener, men skal tillige forklare,
hvorledes det er muligt, at der kan opstaa psykisk Væren (Bevidsthed).

For mit Vedkommende maa jeg med Hensyn til dette sidste Punkt bemærke,
at der vel maa skelnes mellem de to Ting, der hos \textsc{Driesch} gaar sammen i Et.
Jeg antager det for uundgaaeligt at forudsætte psykiske Analoga paa lavere Trin af
materiel Existens, hvis man ikke vil antage, at Bevidsthedslivet opstaar med et
Spring, og hvis man ikke vil paatage sig at bevise dets Opstaaen af rent materielle
Aarsager. Men man har derfor ikke Ret til at lade et saadant Analogon optræde
som deus ex machina i Biologien. ---

Et med Vitalisme beslægtet Begreb er Begrebet
Orthogenesis\footnote{\textsc{Bergson}'s „élan vital“ er i \textsc{Driesch}'s
Øjne for ubestemt et Udtryk. Han foretrækker Begrebet Individualitet, fordi der
her er Tale om et Grundbegreb, en Kategori. (\textit{Science and Philosophy of the
Organism}, II, p. 305, 313). Iøvrigt finder han mange Berøringspunkter mellem sin
og \textsc{Bergson}'s Opfattelse (ib. p. 292 Note). Om \textsc{Bergson}'s Vitalisme
henviser jeg til min Bog \textit{Henri Bergson's Filosofi}, p. 35---42.}, som er opstillet
til Forklaring af, hvorledes Udviklingen kan gaa frem fra lavere til højere Former.
Denne Forklaring finder man saa i en oprindelig Trang eller Tendens til Udvikling
i en bestemt Retning. Der er i denne Sammenhæng ingen Grund til at gaa ind paa
dette Begreb, da Kritiken af det vil være ganske tilsvarende til Kritiken af
Vitalismen. Begrebet Orthogenesis er for Udviklingsbegrebet, hvad Begrebet Livskraft er
for Organismens Begreb i det hele. ---

Resultatet af vor Drøftelse bliver, at Umuligheden af en mekanisk Oprindelse
af Livet ikke kan godtgøres, --- hvorved dog den Mulighed, at Livet ingen absolut
Oprindelse har, ikke maa tabes af Syne. Dersom det skulde kunne godtgøres, at
Livet, Organismen med dens Helhedspræg, kan forklares ad mekanisk Vej, saa
mister Begrebet Helhed derfor ikke sin Betydning. Tvertimod viser Helheden sig
da netop at have en dyb Grund i Tilværelsens Elementer, siden disse i Kraft af
deres egne Love kan forenes til at danne Helheder. Det er da de samvirkende
Aarsagsrækker, der gør saadanne komplexe Fænomener som Organismer mulige.
Helhedsbegrebet stiller, som enhver Kategori, en Opgave, et Spørgsmaal, men siger i
sig selv Intet om, hvorledes dette Spørgsmaal skal besvares.

Den specielle biologiske Forsker behøver ikke at lægge hverken Mekanicisme
eller Vitalisme til Grund. Det er hans Sag at anvende fysiske og kemiske
Synspunkter indenfor Organismen saa langt, som det overhovedet er muligt. Og nogen
absolut Sikkerhed for, at han staar ved Mulighedernes Grænse, kan han aldrig faa. ---

Det har ved hele denne Drøftelse været Hovedopgaven at faa det almindelige
% --- p. 56 ---
\opage{56}Problem frem og vise dets Sammenhæng med den menneskelige Tankes Natur,
saaledes som den lægger sig for Dagen paa de forskellige Omraader, hvorfra den henter
sine Emner. Det er let nok at spotte over Vitalismen paa en mere eller mindre
smagfuld og vittig Maade. Men man tjener Tænkningen bedre ved klart og bestemt
at pege paa det Problem, der her atter og atter gør sig gældende, selv om der er
dem, der lader, som om det ikke existerer.

\medskip
21. Personligheden, Selvet, er et andet Exempel paa en Helhed, der træder os
umiddelbart imøde for derefter at gøres til Genstand for Analyse. Forholdet mellem
den givne Helhed og den videnskabelige Undersøgelse er her et Sidestykke til, hvad
der gjorde sig gældende for det organiske Livs Vedkommende. Her fremtræder en
Modsætning mellem en Opfattelse, der betragter Personligheden (og Bevidsthedslivet
i det hele) som et blot Resultat af et Sammenspil af Elementer, og en anden
Opfattelse, der finder Forklaringen til de enkelte Bevidsthedsfænomener i Selvet,
Personligheden som Helhed. Den første Opfattelse repræsenteres dygtigst og paa mest
karakteristisk Maade af den engelske „Associationspsykologi“ (\textsc{Hume}, \textsc{James Mill},
\textsc{Stuart Mill}); den anden hyldes ikke blot af den spiritualistiske Psykologi, der
lader Sjælen som en særegen „Substans“ ligge til Grund for Bevidsthedsfænomenerne
og gribe ind i deres Gang, men ogsaa af nyere psykologisk Opfattelser, der lader
„Jeget“ eller „Selvet“ øve en lignende Funktion. Ogsaa paa dette Omraade er det
Loverkendelsen, forsaavidt en saadan er mulig, der fører ud over de stridende
Opfattelser. Den eneste Lov, der kan spores overalt paa det psykiske Omraade, er
Syntesens Lov. Naar Begrebet Syntese opfattes som et Grundbegreb i Psykologien,
menes derved ikke Begrebet om en Kraft, men netop Begrebet om en Lov eller en
Form for alt Bevidsthedsliv. Hvad psykologisk Analyse kan finde, vil stedse vise
sig at være specielle Former for denne Grundform. Det er ikke overflødigt at
bemærke, at Spørgsmaalet om, hvorvidt Syntesens Grundlov kan paavises overalt i
Bevidsthedslivet, er uafhængigt af Besvarelsen af det Spørgsmaal, om
% STRAY MARK (p. 56): a tick after the line-end hyphen of "Bevidstheds-" (layer
% reads "-'"); a mark in the copy, not transcribed.
Bevidsthedslivet overhovedet kan forklares som Produkt af fysiske og kemiske Elementer.
Hvad enten man besvarer dette Spørgsmaal med Ja eller Nej, bliver Psykologiens
egen Opgave dog den samme; Forholdet mellem Emne og Tydning indenfor
Psykologien forandres ikke derved. Selve de Kategorier, hvormed den fysiske og kemiske
Forsken arbejder, er specielle Former for Syntese, der ligesom alle andre
Kategorier har udviklet sig under Tankens Stræben efter at tyde foreliggende Emner.
Tanken selv er baade fra et psykologisk og et erkendelsesteoretisk Synspunkt, sit
eget sidste Problem, hvorledes saa alle andre Problemer er blevne besvarede.

Forholdet mellem den psykologiske og den erkendelsesteoretiske Synsmaade er
drøftet ovenfor (§ 5---6). Vi gør i nærværende Sammenhæng Psykologiens Stade til
Genstand for erkendelsesteoretisk Undersøgelse, ligesom Psykologien fra sin Side
undersøger, hvorledes det er muligt, at en videnskabelig Bevidsthed kan udvikle
% DEFECT (p. 56): a line-end hyphen stands after "udvikle" ("udvikle-" / "sig."),
% though two words are meant; set as two words, hyphen not transcribed.
% UNSURE: rejected "udviklesig" (the layer's joined reading).
sig. Hvad der erkendelsesteoretisk er den eneste Vej til Sandhed, det er
psykologisk set kun en af mange mulige Retninger, Bevidsthedslivet kan slaa ind paa, og
% --- p. 57 ---
\setcounter{footnote}{0}%
\opage{57}det bliver da at undersøge, hvorledes den forholder sig til de andre Retninger.
Naar Psykologien gøres til Genstand for Erkendelsesteorien, er det for at undersøge
Berettigelsen af den Vej til Sandhed, Psykologi som Videnskab følger. Og det
vigtigste Spørgsmaal er her det om Bevidsthedslivets Helhedskarakter.

Baade subjektiv og objektiv Psykologi, d. v. s. baade Selviagttagelse og
Iagttagelse af de ydre Fænomener, til hvilke Bevidsthedsliv paa en eller anden Maade
antages knyttet, fører til at lægge Vægt paa Bevidsthedslivet som en lille Verden,
hvis Elementer staar i stadig Vekselvirkning indbyrdes efter bestemte Love, og hvis
Nutid bestemmes ved dens Fortid. I Forholdet mellem Fortid og Nutid, altsaa i
den historiske Karakter af Bevidsthedslivet, lægger Helhedsloven, Syntesens Lov sig
særlig tydelig for Dagen. Den historiske Karakter har nu, som vi senere skal
se, ogsaa alle materielle Naturfænomener. Der maa ogsaa ved dem stedse tages
Hensyn til Fortiden. Det gælder om Organismer, Kloder og Stjernesystemer, og det
gælder endog om saa livløse Ting som Fabrikata. To Søm kan se aldeles ens
ud; men hvis det ene har været hamret flere Gange, vil dets molekulære Tilstand
være forskellig fra det andets, der ikke er undergaaet en saadan Behandling.
Statiken er overhovedet, naar den isoleres fra Dynamiken, „en ufuldstændig og
ufrugtbar Videnskab, en Gren, der er skilt fra
Stammen“\footnote{\textsc{Painlevé}: \textit{Mécanique}. (De la Méthode dans les
Sciences, 1909). p. 370, 400.}. Men i Bevidsthedslivet
fremtræder denne historiske Karakter paa den mest betydningsfulde Maade.

Denne historiske Karakter er given med Syntesebegrebets Stilling i Psykologien.
De foregaaende Elementers og Tilstandes bevidste og ubevidste Eftervirkninger
forbindes stadig med de følgende Elementer og Tilstande, og derved bestemmes disses
Grad og Beskaffenhed. Fechners Lov og Kontrastloven peger i denne Retning.
Særlig tydelig fremtræder dette Forhold ved Bevægelsesfornemmelser. For at vor
egen Bevægelse skal fornemmes, maa de tidligere Paavirkninger fra Muskelfibres,
Leds og Hudens Tilstande og Stillinger stedse kombineres med de følgende, og kun
formedelst denne stadige Sammenfatten opstaar Bevægelsesfornemmelser. Hvad der
fysiologisk set er en successiv Mangfoldighed, fremtræder psykologisk som en
kontinuerlig Helhed (se min Psykologi\textsuperscript{6} V A, 6. Slutn.). Man har med Rette i
Henhold hertil talt om en „motorisk Helhedsfornemmelse“ som grundlæggende for
Opfattelsen af os selv som villende og handlende
Personer\footnote{\textsc{Bror Gadelius}, se ovenfor p. 25.}. Som Sidestykke til
saadanne elementære Fænomener, i hvilke dog Bevidsthedslivets Grundlov lægger
sig for Dagen, kan vi henvise til Beslutningen --- i dette Ords strengeste
Betydning, i Modsætning til Drift og til Forsæt; --- i den bestemmes en Afgørelse ved
alle Elementer i Personligheden og ved dennes Historie i Fortiden. I en saadan
Beslutning fremtræder baade den reale og den formale Side ved Personligheden,
baade Indhold og Form, baade „Sjæl“ og „Aand“, paa udpræget Maade.

Naar man ofte lægger Vægt paa Selvbevidstheden, Selvgenkendelsen som
karakteristisk Udtryk for Personligheden, saa maa bemærkes, at den Selvidentificering,
% PRINTED AS IS (p. 57): "idenfor" for indenfor (checked at 300 dpi).
disse Udtryk betegner, netop forudsætter en omfattende Syntese, idenfor hvilken en
% SIGNATURE LINE at the foot of p. 57 ("D. K. D. Vidensk. Selsk. Skr., 7.
% Række, hist. og filos. Afd. III. 2." and "40") omitted as printer's apparatus.
% --- p. 58 ---
\setcounter{footnote}{0}%
\opage{58}Analyse, en særlig Fremdragen af de Elementer, der betinger Selvgenkendelsen, kan
finde Sted. Og selve de Elementer, der identificeres i Selvgenkendelsen, maa
sammenfattes, for at denne kan blive mulig. Som \textsc{Kant} paa sit Sprog udtrykker det,
Apperceptionens analytiske Enhed (ɔ: Forestillingen om Bevidsthedens Identitet i
en Mangfoldighed af givne Forestillinger) er kun mulig under Forudsætning af dens
syntetiske Enhed. Man maa erindre, siger han fremdeles, at Bevidsthed ikke er
en [enkelt] Forestilling, men en Form for Forestillen. Vi bevæger os her i en Kreds,
da vi stedse maa gøre Brug af Bevidstheden for at kunne udsige Noget om
den\footnote{\textit{Kritik der reinen Vernunft}\textsuperscript{2}, p. 133; 404.}.

% PRINTED AS IS (p. 58): "Helhedsbegreaet" for Helhedsbegrebet (checked at 300 dpi;
% the layer reads "Helhedsbegreael").
Ogsaa fra en anden Side fremtræder Helhedsbegreaet paa det psykologiske
Omraade. De skarpe Forskelle mellem psykiske Tilstande og mellem psykiske
Elementer, de Forskelle, der ligger til Grund for de Distinktioner og Inddelinger,
den analytiske Psykologi foretager, fremtræder først klart under Bevidsthedslivets
mere fremskredne Udvikling. De mere primitive Tilstande er mere ensartede og
tillige mere koncentrerede. Udviklingen bestaar her, som paa andre Omraader, i
fremskridende Differentiation. For at forstaa en afgørende Forskel eller
Modsætning paa det psykiske Omraade maa man derfor se tilbage paa den historiske
Udvikling, der har ført til den. Det samme gælder naturligvis om de Koncentrationer,
der igen (f. Ex. i Beslutning efter Overvejelse) kan afløse Differentiationer. Men
det er af største Vigtighed at mærke sig denne historiske Karakter, for at man
hverken skal komme til at lægge for megen eller for lille en Vægt paa psykologiske
Inddelinger. Megen Kritik, som er rettet mod „Tredelingen“ i Psykologien, bliver
genstandsløs, naar det her antydede Synspunkt (som er lagt til Grund i min
Fremstilling af Psykologien, se dennes Kap. IV) fastholdes. Jeg ser her bort fra, at man
skematisk kan benytte Tredeling og dog foretage en Reduktion af de tre Arter af
Elementer til en (for mit Vedkommende til Villieselementerne).

Med Bevidsthedslivets historiske Karakter hænger det nøje sammen, at
psykologiske Problemer for en meget stor Del beror paa tilsyneladende Afbrydelser af
% PRINTED AS IS (p. 58): "Bevidstlivet" (elsewhere Bevidsthedslivet), checked at 300 dpi.
den direkte Sammenhæng i Bevidstlivet. Det gælder i saadanne Tilfælde at finde
en dybere liggende Sammenhæng i Kraft af de Love, der har vist sig at gælde
for Bevidsthedslivet. Et vigtigt Exempel herpaa er Læren om Motivforskydning,
der grunder sig paa den Indflydelse, Forestillinger faktisk øver paa Følelser; i denne
Indflydelse ligger en Mulighed for, at naar en Forestilling fører til en anden
Forestilling, vil ogsaa de Følelser, der var knyttet til den første Forestilling, kunne blive
knyttet til den anden og vedblive at være det, selv om den første Forestilling helt
falder bort. Der foreligger her en Mulighed for aandelig Udvikling, som \textsc{Spinoza}
og de engelske Psykologer\footnote{Jeg benytter Lejligheden til at rette en Fejl
i „Den nyere Filosofis Historie“, II, p. 372 f. Det her omtalte „Appendix B“ i
\textsc{James Mill}'s „Fragment on Mackintosh“ er Aftryk af et Kapitel hos
\textsc{Mackintosh}. Æren for den klare Fremstilling af Motivforskydningens Teori
tilkommer altsaa ikke \textsc{James Mill}, men \textsc{Mackintosh}, skønt
\textsc{James Mill} i Principet hylder den samme Teori.} har fremdraget. Paa
Erkendelsens Omraade haves et stort Exempel i Arbejdet paa at overvinde de
% PRINTED AS IS (p. 58): "kaotisk Forskelsrækker" (the sense wants kaotiske), checked at 300 dpi.
Hindringer for kontinuerligt Tankeliv, som kaotisk Forskelsrækker frembyder.

% --- p. 59 ---
% STRAY MARKS on p. 59 (read by the layer as "K", "*", "/"): ticks at the
% left edge and margin of the scanned copy; the head carries only the folios
% 59 / 315. Not text of the work.
\setcounter{footnote}{0}%
\opage{59}Der gives dog ogsaa psykologiske Problemer af anden Art, nemlig saadanne,
i hvilke det gælder at forstaa komplexe psykiske Fænomeners Natur og Tilblivelse.
Det gælder at beskrive disse specielle psykiske Helheder, analysere dem og finde
Betingelserne for deres Opstaaen og Bestaaen. Der kan f.~Ex. spørges om den
religiøse Tros, Samvittighedens og Kunstinteressens Væsen og Betingelser. I min
Religionsfilosofi har jeg undersøgt Religionen fra dette Synspunkt, og i „Den store
Humor“ har jeg givet et andet Exempel paa en saadan Undersøgelse. ---

Jeg har hidtil væsentlig holdt mig til subjektiv Psykologi, der støtter sig paa
Selviagttagelse, men iøvrigt meget godt kan benytte sig af alt, hvad der kan
oplyses ad experimental, fysiologisk og historisk Vej. En rent objektiv Psykologi
vilde forsøge at se helt bort fra Selviagttagelse og holde sig til foreliggende ydre
Iagttagelser, især af fysiologisk Art. Det vil selvfølgelig bero paa en Illusion, naar
man mener ad denne Vej at naa til en Psykologi, ti man vil uden det Kendskab
til Bevidsthedslivet, Selviagttagelsen giver, slet ikke kunne vide, hvad de fysiologiske
Fænomener, man ser for sig, betyder. Et saadant Forsøg er alligevel gjort af
\textsc{Bechterew} i hans „Objective Psychologie“. Jeg ser her helt bort fra, at en
kritisk Gennemgang af dette Værk paa mange Punkter viser Selviagttagelsen at være
skjult tilstede. Men hvad der har Interesse i denne Sammenhæng, er, at
\textsc{Bechterew} selv fra sit rent fysiologiske Standpunkt mener at kunne paavise
en „personlig Sfære“, der især træder frem i sin Selvstændighed, naar Organismens
Tilstand ikke bestemmes ved de aktuelle Paavirkningers Beskaffenhed, men ved
„Sporene“ af tidligere Paavirkninger. Den personlige Sfære er Indbegrebet af
saadanne Spor, som udgør Organismens „Nevropsyki“, og hvis Reaktioner vi kaldte
Handlinger eller Gerninger\footnote{\textit{Objective Psychologie}. Deutsche Übers.
(1913), p. 432, 435.}. Ogsaa fra et rent objektivt Synspunkt fremtræder altsaa
Personligheden som en centraliseret, historisk udviklet Helhed, som det er
Forskningens Opgave at analysere.

\medskip
22. Paa det sociologiske Omraade staar, ligesom paa det biologiske og det
psykologiske, en „mekanisk“ Opfattelse ligeoverfor en „organisk“. Spørgsmaalet
er, om en Helhed --- her et Samfund eller en Stat --- kan forklares blot ved ydre
Forbindelse af Elementer, der i og for sig kunde existere uden saadan Forbindelse.

En rent udvortes Forbindelse synes at foreligge, naar forhen uafhængige Individer
indgaar en Overenskomst, en Kontrakt, i hvilken de gaar ind paa at overholde
visse Regler, saa at et Samliv og en Samvirksomhed bliver mulig. Det er da ogsaa
som en Slags Kontrakt, at den mekaniske Sociologi, saaledes som den fremtræder i
den gamle Naturret og i mange Former af moderne Individualisme, tænker sig
Samfundets Opstaaen. Oppositionen mod et bestaaende, i bestemte Former
% PRINTED AS IS (p. 59): "krystaliseret" with one l (p. 62 has "Krystallisation").
krystaliseret Samfund benytter ofte denne Synsmaade. I skarp Modsætning til
Kontraktteorien hævdede i det nittende Aarhundrede den saakaldte historiske Skole
en „organisk“ Opfattelse af Samfundet, ifølge hvilken dette fremtræder som en
Helhed, ofte en
% SIGNATURE MARK "40*" at the foot of p. 59 omitted as printer's apparatus.
% --- p. 60 ---
\setcounter{footnote}{0}%
\opage{60}mystisk Helhed, der, langtfra at være en blot Sum af selvstændige Individer, netop
bestemmer Individerne i deres inderste Natur og i deres Handlemaade.

\textsc{Hobbes} er den mest udprægede Repræsentant for den første Opfattelse,
\textsc{Comte} og \textsc{Hegel} for den anden.

Naar den første Opfattelse benytter Synspunktet af en Kontrakt til at forklare
Samfundets Opstaaen, benytter den et Synspunkt, der kun har nogen Mening, naar
der allerede bestaar et Samfund og en med Magt hævdet Retsorden. Der
forudsættes altsaa, hvad der skulde godtgøres. Allerede de gamle Naturretslærere
har indset dette, naar de siger, at den samfundsstiftende Kontrakt er en
stiltiende Forudsætning (a supposed covenant, som \textsc{Hobbes} udtrykker sig),
altsaa ikke indgaaet med klar Bevidsthed, men „liggende i“ den uvilkaarlige Maade,
paa hvilken Forholdet mellem Mennesker kan ordne sig. Samfundet og dets Ordning
er altsaa en Frugt af den Retning eller det Sammenspil af Retninger, Livet
allerede uvilkaarlig havde taget hos Individerne, før der bestemt kunde være Tale
om Samfundsformer og Institutioner. Retningsbegrebet trænger sig ogsaa her frem.
Ligesom det er en Illusion, naar man taler om Atomer, Kræfter eller Bevægelser,
uden at tage Hensyn til de bestemte Retninger, i hvilke Atomer bevæger sig og
virker, saaledes er det ogsaa en Illusion, naar man taler om Individer uden
Hensyn til den Maade, paa hvilken Livet har ført dem til Samliv og Samvirken med
Andre, allerede før de med Bevidsthed kan ordne deres Forhold. Der har da aldrig
% PRINTED AS IS (p. 60): the quote opened by „ closes with a left-pointing
% guillemet instead of “ -- so „“ are unbalanced by one here (250 dpi).
været en Tid, da Menneskene levede i ren „Naturtilstand\guillemotleft. Det fremtræder paa
karakteristisk Maade i den Sætning hos \textsc{Hobbes}, at i Naturtilstanden gives
der ingen Sønner (filium in statu naturali intelligi non posse, De cive II, 10),
da allerede Forholdet mellem Moder og Barn, naar det skal have retslig eller
moralsk Betydning, maa ses under Synspunktet af en stiltiende Kontrakt, som
Barnet indgaar ved at modtage Moderens Hjælp. Bedre kunde det ikke udtrykkes, at
Samfundet er tilstede i Menneskelivets første Ytringer. \textsc{Hobbes} har da
ogsaa klart set, at han for at gennemføre sin sociale Mekanik maatte tænke sig
Menneskene skyde op af Jorden ligesom Paddehatte, hvert for sig, og saa derefter
indgaa Forbindelser. Den Analogi, \textsc{Hobbes} i Virkeligheden støtter sig
til, er den mellem Individer og Stater. Forskellige Stater kan jo blive til
ganske uafhængig af hverandre, og i Forholdet mellem dem kan der derfor være Tale
om en Naturtilstand, der ifølge \textsc{Hobbes} væsentlig er en
Krigstilstand\footnote{In all times, Kings, and Persons of Sovereign authority,
because of their Independency, are in continual jealousies, and in the state and
posture of gladiators, having their weapons pointing, and their eyes fixed on one
another. \textit{Leviathan}, I. Chap. 13.}. Han tænker sig Forholdet mellem
Individer i Lighed med Forholdet mellem suveræne Stater. Men Spørgsmaalet vender
da tilbage, naar vi vender Opmærksomheden mod enhver enkelt af disse Staters
Tilblivelse.

Karakteristisk for den naturretlige Opfattelse er, at man vil forklare
uvilkaarlige, tildels ubevidste sociale Tilstande ved Hjælp af saadanne sociale
Processer, der foregaar med mere eller mindre klar Bevidsthed hos de enkelte
Individer, der deltager i dem. Det er en analog Fremgangsmaade som den, der fører
mange Psykologer
% --- p. 61 ---
% STRAY MARKS on p. 61 (read by the layer as a low quote and "*"): ticks at
% the left edge of the scanned copy; the head carries only the folios 61 / 317.
% Not text of the work.
\opage{61}til at forklare Sjælelivets primitive og uvilkaarlige Ytringer ved Hjælp af
Elementer, der er fundne ved Analyse af højt udviklede og differentierede
Tilstande. ---

Den „organiske“ Opfattelse hævder, at Samfundslivet stedse er tilstede i en
eller anden Form og i et eller andet Omfang, saasnart der overhovedet existerer et
Menneskeliv. I den primitive Familie, i hvert Tilfælde i Forholdet mellem Moder
og Barn, og i Livet i Flok, Fænomener, der forekommer overalt, hvor Mennesker
lever, er der allerede givet et Samfund. Efter denne Opfattelse er Samfundet det
primære og Individet, saa selvstændigt det end kan udvikle sig, socialt betinget og
bestemt efter hele sin Natur og alle sine Livsvilkaar. Enhver menneskelig
Livsytring og Livsform finder sin sidste Forklaring i social Indflydelse. Ved denne
% BROKEN SORT? (p. 61): the e of „kommer“ prints without its crossbar and
% reads as c („kommcr“); taken as a broken e (second reading 2026-09-30).
Opfattelse kommer det Ubevidste og det Uvilkaarlige, alt hvad der gaar forud for
klar Bevidsthed, Slutning og Beregning, til sin Ret. Det ses, at de enkelte
Individer ikke opstaar eller udvikler sig isolerede, men er, hvad de er, ved hele
den Sammenhæng, indenfor hvilken de er blevne til. Her er Helheden ikke et
mekanisk Aggregat, men en oprindelig Realitet, hvis Historie er den egentlige
Historie.

Den „organiske“ Opfattelse passer bedst paa Samfundslivet i dets mere primitive
Former, medens den „mekaniske“ Opfattelse kan beraabe sig paa Erfaringer fra
senere sociale Udviklingstrin. Som Sir \textsc{Henry Maine} i en Række
betydningsfulde Fremstillinger (først i „Ancient Law“ 1861) har vist, kan
Forskellen mellem mere primitive og mere udviklede Samfundsformer betegnes som en
Forskel mellem Tilstand (status) og Pagt (contractus). Senere er denne Modsætning
nærmere gennemført af \textsc{Spencer}, \textsc{Durkheim} og \textsc{Tönnies}.
Særligt har \textsc{Tönnies} (Gemeinschaft und Gesellschaft 1887) vist, at denne
Modsætning hænger sammen med Modsætninger mellem den uvilkaarlige Karakter,
Villieslivet har paa sine tidligere Trin, hvor det umiddelbart udspringer af
Naturen (som „Wesenswille“), og den bevidste, ved Valg mellem bestemte Muligheder
% PRINTED AS IS (p. 61): „Wilkür“ (for Willkür), as printed.
betingede Villen („Wilkür“). I Modsætning til Instinkt og uvilkaarlig Hengivelse
og Tilslutning træder nu Sammenligning og Overvejelse af Maal og Midler, ligesom i
Samfundet Familiebaand afløses af Overenskomst, Sæd og Skik af Lovregler og
Interessebrydning, og ligesom Troen afløses af Videnskaben. Det er to sociale
Helhedstyper, som kan komme igen paa ellers højst forskellige Kulturtrin og afløse
hinanden paa rytmisk Maade.

Men indenfor begge Typer, den kontraktbestemte saavel som den naturbestemte,
gælder det, at Individerne stedse hører til et Hele, medens paa den anden Side
Samfundslivet bestaar i og gennem de enkelte Individer, hvis Initiativer ofte
bliver af afgørende Betydning for social Udvikling.

Overalt hvor Totalitetsbegrebet optræder, gælder det, at det ikke kan bruges
til Forklaring af enkelte bestemte Fænomener. Ligesaalidt som Livskraften
forklarer noget enkelt organisk Fænomen, og ligesaalidt som Sjælens Begreb
forklarer noget enkelt psykisk Fænomen, ligesaalidt forklarer Samfundsbegrebet
noget enkelt socialt Fænomen. Hvad det gælder om, paa det sociale, saavel som paa
det psykologiske og det biologiske Omraade, er at finde den Lov, eller de Love,
gennem hvilke en Helhedssammenhæng lægger sig for Dagen. Der gives en sociologisk
% --- p. 62 ---
\setcounter{footnote}{0}%
\opage{62}Mystik, ligesom der gives en psykologisk og en biologisk Mystik. Jeg finder
f.~Ex. en saadan Mystik i \textsc{Durkheim}'s Sætning, at et Faktum ikke er
socialt, fordi det genfindes hos alle Samfundets Individer, men omvendt:
% PRINTED AS IS (p. 62): the French as printed -- "re" for se, "repète" for
% répète, "induvidus" for individus, "qu'it" for qu'il.
il re repète chez les induvidus parce qu'it s'impose à
eux\footnote{\textit{Les règles de la méthode sociologique}, p. 14.}. Det er altsaa
i Delene, fordi det er i Helheden, ikke omvendt. Men en social Form maa have
udviklet sig, før den kan paanøde sig Individerne. Og den udvikler sig kun gennem
den stadige Vekselvirkning mellem Individerne, en Vekselvirkning, i hvilken de
enkelte Individer hvert for sig er baade givende og modtagende. Paa intet Punkt
forsvinder den individuelle Selvstændighed ganske. Hvert enkelt Individ er jo en
lille Helhed, saaledes som baade Biologien og Psykologien maa indskærpe.
Selvstændigheden ytrer sig baade i Modtagelighedens Art og Grad og gennem
Initiativer, der kan udformes gennem Vekselvirkning med andre Individer. Under
denne Vekselvirknings Forløb kan det nye Indskud, der udspringer i et enkelt
Individs Indre, undergaa Forandringer dels gennem uvilkaarlige Forskydninger og
dels ved bevidste Tilføjelser og Kombinationer, dels ved Brydning med Initiativer
fra andre Sider. \textsc{Grimm} lærte, at Folkedigtningen var udsprungen af
Folkets Sjæl, men \textsc{A.~W. Schlegel} hævdede overfor denne mystisk-poetiske
Opfattelse, at den kun skyldes den Omstændighed, at vi ikke kender Forfatterne til
de enkelte Digte\footnote{\textsc{C.~T. Gooch}: \textit{History and Historians of
the nineteenth century} (1913), p. 57.}. Og \textsc{Burckhardt} har udført dette
videre: „Folkeepos, Folkesang og Folkemelodi synes ikke at skyldes enkelte store
Individer; istedetfor disse træder et helt Folk, som vi ad hoc forestiller os i en
særlig lykkelig-naiv Kulturtilstand. Men denne Substitution beror hovedsagelig
paa Overleveringens Mangelfuldhed. Den episke Sanger, hvis Navn vi ikke mere
kender (eller kun kender som Kollektaneum), har været stor i det Øjeblik, da han
for første Gang gav en Gren af sit Folks Sagnkreds en varig Form; i dette Øjeblik
koncentrerede Folkeaanden sig i hans paa magisk Maade~.~.~.~. Den koncentrerede
Folkeaand taler ud af ham; ellers vilde Sangen ikke leve gennem
Tiderne“\footnote{\textsc{Jacob Burckhardt}: \textit{Weltgeschichtliche
Betrachtungen} (1905), p. 224 f.}. --- Alt nyt maa stamme fra de enkelte, tilsidst
fra et enkelt Individ, i hvis Bevidsthed en ejendommelig Krystallisation har
fundet Sted. Det bliver saa derefter en ny Opgave at undersøge, hvorledes denne
Krystallisation har kunnet finde Sted, og her vil man igen staa overfor
Vekselvirkningen mellem Individerne og mellem dem og de bestaaende Overleveringer.
Det er Mytologi, naar den historiske Skole (især den preussiske) betragter de
enkelte Individer blot som sekundære. De nye Begyndelser kan, før de fører til
social Organisation, kun studeres i de enkelte Individers Bevidsthed, saavidt den
er tilgængelig. Her kan opstaa Mutationer af højst forskellig Levetid. Her
foregaar en stadig Brydning mellem det Nye og det Overleverede, allerede i den
enkeltes Sind. Saadanne Processer kan ikke forstaas, naar man gør Samfundet til en
mystisk virkende Kraft og hylder en social Vitalisme.

Et interessant Exempel paa en Brydning af denne Art er \textsc{Erik Gustaf
Geijer}'s Forhold til den historisk-romantiske Skole, han med Begejstring havde
tilhørt.
% --- p. 63 ---
\setcounter{footnote}{0}%
\opage{63}Efter hans eget Udsagn var det Historien selv, der førte ham bort fra den
historiske Skole. Historien viste ham de individuelle Initiativers Betydning. Han
fik Øje for Muligheden af, at en Generation eller en Gruppe af Mennesker i Tillid
til Fremtiden og dens Muligheder saa at sige kunde gaa fra Arv og Gæld (gøre sig
urarfva).\footnote{%
% DEFECT (p. 63 note 1): a stray low mark under the comma after "I" (layer
% reads "I,("); not transcribed. "før „Affaldet“" as printed.
\textit{Samlade Skrifter}, I, 7, p. 137. --- Jeg er bleven
opmærksom paa dette Sted gennem Dr. \textsc{Hilma Borelius}' Afhandling om Geijer
før „Affaldet“. --- Smlgn. iøvrigt \textit{Personlighetsprincipen i Filosofin},
p. 39---49, hvor jeg drøfter Problemet i sin Almindelighed.} Overleveringen skydes
til Side, der slaas ind paa nye Baner, og der kommer tilsyneladende Brud paa
Kontinuiteten. Men alle saadanne tilsyneladende Brud opfordrer kun Forskeren til
at finde en dybere liggende Sammenhæng end den, man hidtil var standset ved. Et
Spring eller et Brud betyder for Videnskaben kun en Opgave, hvad enten denne kan
løses eller ikke. Man har med Rette sagt om \textsc{Alexis de Tocqueville}, at han
i sin Bog „L'ancien Régime“ har gjort for Historien, hvad \textsc{Lyell} gjorde
for Geologien, idet han gendrev Katastrofeteorien. Revolutionen kom nu ikke som et
Under, men staar som forberedt gennem det absolute Monarkis aarhundredlange
Tendens til at koncentrere al Magt paa et enkelt Punkt; det revolutionære
Regimente var kun en Fortsættelse af denne Koncentration og Revolutionen blev
mulig ved, at det kun gjaldt at trænge ind paa hint ene Punkt.

Jo mere en dybere liggende Sammenhæng findes, desto vanskeligere bliver det at
inddele Historien i Perioder. Den sædvanlige Inddeling af Europas Historie i
Oldtid, Middelalder og den nyere Tid skal stamme fra det syttende Aarhundredes
Filologer, for hvem de Aarhundreder, der gik forud for Humanismen, kom til staa
% PRINTED AS IS (p. 63): "kom til staa / at staa" -- dittography across the
% line-end, as printed.
at staa som en afsluttet Periode\footnote{\textsc{Joh. Steenstrup}:
\textit{Historieskrivningen} (1915), p. 150.}. For en senere Tids Forsken kan
Middelalderen synes at opløse sig i en hel Række af Renaissancer ligefra det
niende til det femtende Aarhundrede, og det kan blive et Spørgsmaal, om vi ikke
lever i Middelalderen endnu\footnote{\textsc{François Picavet}: \textit{Histoire
comparée des philosophies médiévales}. Chap. 4.}. Og en fremragende
Historiker\footnote{\textsc{Freeman}: \textit{Methods of Historical Study} (1886),
p. 20.} har endog erklæret, at han ikke anerkender Modsætningen mellem Oldtidens og
den nyere Tids Historie.

\medskip
23. Heller ikke paa det sociologiske Omraade, indeholder da Totalitetsbegrebet
i sig selv nogen Forklaring. Det gælder at finde den Lov, der holder Elementerne
sammen, og i Kraft af hvilken Helheden bestaar.

Sociologi er en historisk Videnskab. Samfundslivets Fremtid kender vi ikke;
dets Nutid kan vi ikke overse; det er da Samfundslivets Udvikling i Fortiden, der
er det Givne, Emnet. Historieforskningen gaar ud paa Beskrivelse og Forstaaelse af
Begivenheder, Institutioner og Personligheder. Den anvender herved samme Metode
som Naturvidenskaben. Dens Virkelighedskriterium er det samme. Fortiden er ikke
umiddelbart given, kun Beretningen om den eller Mindesmærker fra den; men
Naturforskeren har jo heller ikke Naturen umiddelbart given, kun Iagttagelser af
Naturfænomener. Ligesom Naturforskeren ved sammenlignende Iagttagelser af
% --- p. 64 ---
\setcounter{footnote}{0}%
\opage{64}Naturfænomener skal danne sig et Helhedsbillede af en Natursammenhæng, saaledes
skal Historikeren gennem Sammenligning af Beretninger og Mindesmærker danne sig en
sammenhængende Opfattelse af Fortiden. Overalt er det den lovmæssige Sammenhæng,
der er Kendemærket paa Virkelighed. Ligesom Modsigelse mellem Iagttagelser skaffes
afvejen ved Paavisning af, at de forskellige Iagttagelser skyldes forskellige
Synspunkter, Sanseorganer eller Iagttagelsesmidler, saa at hvad der fra et
Synspunkt fremtræder paa en vis Maade, fra et andet Synspunkt konsekvent maa
fremtræde paa en anden Maade, --- saaledes kan Modsigelse mellem forskellige
Beretninger forklares ved, at samme Begivenhed af forskellige Berettere efter
deres Forudsætninger maatte opfattes paa forskellig Maade. Den engelske Historiker
\textsc{Gardiner} har særlig haft Evne til at se den Betydning, i hvilken
% PRINTED AS IS (p. 64): the final e of "komme" is set as a turned e (ə);
% transcribed as e (250 dpi).
historiske Personligheder maatte komme til at staa for deres
Modstandere\footnote{\textsc{Gooch}: \textit{History and Historians}, p. 362 f.}.
Det er ikke nok at berigtige Iagttagelser eller Beretninger; der maa ogsaa vises,
hvorledes de Iagttagelser og Beretninger, der trænger til Berigtigelse, har kunnet
opstaa. Den fuldkomne Sandhed maa kunne vise, hvorledes det, der fremtræder paa en
vis Maade for et Subjekt, maa fremtræde paa en anden Maade for et andet Subjekt,
--- hvorved Forskeren naturligvis ikke maa glemme, at han ogsaa selv er et
Subjekt. Enhver Definition af, hvad Sandhed er, har kun Betydning ved at stille den
Opgave at finde nye Emner og nye Synspunkter og stedse nøjere Sammenhæng baade
mellem Synspunkterne indbyrdes og mellem Emnerne indbyrdes\footnote{I \textit{Den
menneskelige Tanke}, p. 303---305, har jeg søgt at vise, hvorledes denne
Opfattelse er forenelig med den Sætning, at Sandheden kun kan være én.}.

Der er i denne Henseende kun en Gradsforskel mellem Naturvidenskab og
Historie, en Gradsforskel, der beror paa den Nøjagtighed og Fuldstændighed,
hvormed Virkelighedskriteriet i det Enkelte kan anvendes. For Historieforskeren
staar den exakte Metode, Naturforskeren kan anvende, ikke til Raadighed. Men det
betinger ingen principiel Forskel.

Derimod har man ment at finde en Artsforskel mellem Naturforsken og
Historieforsken deri, at Naturvidenskaben gaar ud paa almene Love, medens
Historievidenskabens Maal er Opfattelse og Forstaaelse af individuelle Fænomener
(Begivenheder, Institutioner, Personligheder). \textsc{Heinrich Rickert} har paa en
meget interessant Maade gjort denne Forskel mellem Historie (Kulturvidenskab) og
Naturvidenskab gældende. Han erklærer udtrykkelig, at han herved lægger Vægten paa
de Formaal, de forskellige Videnskaber stiller sig, ikke paa de Midler, de virker
med\footnote{%
% PRINTED AS IS (p. 64 note 3): "naturwissenschafftlichen", ff doubled (250 dpi).
\textit{Die Grenzen der naturwissenschafftlichen
Begriffsbildung}\textsuperscript{2}, p. 302.}.

Er det da nu Naturvidenskabens sidste Maal at paavise almindelige Love for,
hvad der sker i Verden? Interessen for at finde almindelige Love har naturlig
udviklet sig af Trangen til saa nøjagtig som muligt at skelne mellem Virkelighed
% TRYKFEJL (p. 80): første læs faste
og Indbildning, idet Virkelighedskriteriet netop er den første, ubrydelige
Sammenhæng mellem forskellige Emner. Men Udfinden af almene Love er saa bleven
Genstand for selvstændig Interesse; der kan være en Interesse i at finde almindelige
% --- p. 65 ---
\setcounter{footnote}{0}%
\opage{65}Love, selv om Indsigten i dem ikke faar praktisk Betydning. Den enkelte Forsker
kan stanse ved en Lov og se sit Livs Opgave i at paavise og bekræfte den. Men jo
flere Love der findes, desto mere afslører der sig en stor Helhed for os, som vi
kalder Naturen, og i hvilken de enkelte Loves Sammenspil (eller som \textsc{H.~C.
Ørsted} sagde: Samfoldighed) er det, der gør Helheden til Helhed. Naturen i denne
Betydning, er et Unicum. Tilværelsen er kun én, og den udgør en individuel Helhed,
idet Individualiteten ifølge \textsc{Leibniz}'s Definition bestaar i en Lov for
Kontinuiteten gennem Variationerne. Naar \textsc{Rickert} siger, at det Enkelte og
Individuelle kun interesserer Naturforskeren, forsaavidt det kan bidrage til at
forstaa den materielle Verden som en Helhed (p. 60), saa indrømmer han selv, at
Love tilsidst skal tjene til at forstaa den store Individualitet, vi kalder
Naturen: Totam Naturam unum esse individuum (\textsc{Spinoza}: Ethica II. Lemma 7).
Naturforskningen studerer altsaa Tilværelsens Karakter, eller dens Biografi.
Historieforskningen studerer Menneskehedens, især den europæiske Menneskeheds
Biografi. En højere Opgave, som maaske ligger ud over Erkendelsens Omraade, vilde
være at sammenfatte disse to Biografier til én. Men hvorledes det nu end forholder
sig, --- skal der overhovedet tales om et sidste Maal for Videnskaben, saa hører
dette Maal under Begrebet Totalitet, ikke under det abstrakte Begreb om det
Almene.

\textsc{Rickert} vil hertil svare, at det er de enkelte, begrænsede
Individualiteter, han tænker paa. Historien søger at karakterisere enkelte
Begivenheder eller Personligheder, der er enestaaende og ikke kommer igen; den har
med „das Einmalige“ at gøre. Naturvidenskaben søger derimod i sine Tal og Formler
at bestemme de Forhold, der stedse gentager sig. % DEFECT (p. 65): a stray low mark after the comma of "enkelte," (layer
% reads ",,"); not transcribed.
De enkelte, individuelle Fænomener
er for Naturforskeren kun Exempler (p. 354 f). Herimod har jeg i en Kritik af
\textsc{Rickert}\footnote{\textit{Göttinger gelehrte Anzeigen}. 1906. Jfr.
\textit{Den menneskelige Tanke}, p. 214 f; 224 f. --- Min Kritik angik første
Udgave af
% PRINTED AS IS (p. 65 note 1): "Rickerte's" for Rickert's (220 dpi).
\textsc{Rickerte}'s Bog.} gjort gældende, at der heller ikke i den
materielle Natur hersker absolut Gentagelse, naar man blot analyserer de enkelte
Tilfælde; „das Einmalige“ raader baade i Naturen og i Historien. I den nye Udgave
af sit Værk (1913) (p. 468) indrømmer \textsc{Rickert} i sit Svar paa min
Indvending: „Es findet in der Natur wie in der Geschichte jede Begebenheit nur
einmal statt, und es gibt in Wirklichkeit keine Wiederholungen“. Jeg kan ikke se
andet, end at den principielle Modsætning mellem Naturvidenskab og Historie dermed
falder bort. Maalet maa jo dog være at forstaa Emnerne, og naar Emnerne stedse er
individuelle, maa Erkendelse af det Individuelle være Maalet for enhver
Videnskab. Der vil da, efterat de almene Love er fundne, stedse staa den Opgave
tilbage at forklare, hvorfra det skriver sig, der ikke gentages. „Das Einmalige“
stiller stedse et Problem, hvad enten dette kan løses eller ikke. Det viser sig jo
% DEFECT (p. 65): a stray raised tick over the comma after "ogsaa"; not transcribed.
ogsaa, at Naturvidenskaben stiller sig den Opgave at paavise og forklare
specielle Udviklingsproblemer (af Stjernesystemer, Kloder, Bjerge, Floder,
Organismer, Arter), og det er en fortvivlet Udvej, \textsc{Rickert} griber (p. 449,
454), naar han erklærer, at Udviklingslæren hører til Historiens, ikke til Natur%
% SIGNATURE LINE at the foot of p. 65 ("D. K. D. Vidensk. Selsk. Skr., 7.
% Række, hist. og filos. Afd. III, 2." and "41") omitted as printer's apparatus.
% --- p. 66 ---
\setcounter{footnote}{0}%
% DEFECT (p. 66): a stray dot after the note mark of "Omraade"; not transcribed.
\opage{66}videnskabens Omraade\footnote{Omvendt regner han Sociologien til
Naturvidenskaben (p. 611). Men Sociologi er, som \textsc{Tönnies}
(\textit{Soziologie und Geschichte}. Die Geisteswissenschaften, 1. Jahrg.) viser,
en sammenlignende Historieforsken. Som Resultat af en saadan sammenlignende
Videnskab er \textsc{Tönnies}'s Værk \textit{Gemeinschaft und Gesellschaft} at
opfatte. Det godtgør en typisk Modsætning, ikke en almen Lov, som \textsc{Rickert}
mener (p. 469).}. De forskellige Udviklingshypoteser er netop Forsøg paa at forstaa
individuelle Helheders Tilbliven og Bestaaen ved Hjælp af naturvidenskabelige Love.
Det er en lignende Opgave, Historikeren stiller sig, naar han undersøger
Betingelserne for et historisk Fænomens Indtræden.

\textsc{Georg Simmel}, hvem \textsc{Rickert} (p. 269) betragter som sin
Forgænger, lærer dog netop, at den Totalsammenhæng, som betinger historisk
Forstaaelse, staar over Modsætningen mellem det Almene og det Individuelle, som
man hidtil har betragtet som et uovervindeligt Enten-Eller%
% p. 66: the cue here is printed „2)“ but the note itself is marked „1)“
% in the print (printer's slip; second reading 2026-09-30).
\footnote{\textsc{Georg
Simmel}: \textit{Die Probleme der Geschichtsphilosophie}.
% PRINTED AS IS (p. 66 note 2): "erkenntnisstheoretische", as printed.
Eine erkenntnisstheoretische Studie. 2. Aufl. (1905), p. 34.}.
% PRINTED AS IS (p. 66): "mere til til sin Ret" -- "til" doubled, as printed.
\textsc{Simmel}'s Standpunkt vilde være kommet mere til til sin Ret, hvis han
indenfor det Individuelle igen havde skelnet mellem det typiske (Karakteristiken)
og det konkrete (de enkelte Ytringer og Situationer).

Videnskabens Maal er hverken det Almene eller det Individuelle, men det er
at forstaa Virkeligheden, der stedse er individuel, ved Hjælp af de almene Love,
disse Love, som jo allerede afgør, hvad der er Virkelighed, og hvad der er Drøm,
og som tillige afgiver den eneste Mulighed for at forstaa selve vore Drømme. Vore
fuldkomneste Tanker er hverken Almenbegreber eller konkrete Individualbegreber,
men typiske Individualbegreber.

% CHAPTER HEAD (p. 66) as printed: a short rule, then two centred lines
% ("V." / title with final stop) over a short rule. The Indhold (p. 80) reads
% "Totalitet og Værdikategorierne"; the head on the page governs.
\section*{V.\ Totalitet og Værdi.}
\addcontentsline{toc}{section}{V.\ Totalitet og Værdi}

24. Der er i Nutiden en filosofisk Retning, som vil gøre hele Filosofien til
Værdilære. Filosofiens Arbejdsomraade udgøres da af Værdiproblemerne,
Spørgsmaalene om Sandhed, om Skønhed, om Godhed, om Hellighed. Den gaar ud fra
den Kendsgerning, at disse Ideer gør Fordring paa Gyldighed, og dens Opgave er
da at undersøge, hvori denne Gyldighed bestaar, og hvorvidt den kan retfærdiggøres.
Allerede den blotte Konstatering af en Kendsgerning hviler paa en Vurdering, nemlig
paa Forudsætningen om Sandhedens Værdi. Sandheden er selv en Værdi, der kræver
Anerkendelse. Den teoretiske Tænkning, som vil Erkendelse, er at betragte som et
specielt Tilfælde af en Stræben efter at virkeliggøre Værdier. Forkastes Sandheds%
% --- p. 67 ---
\setcounter{footnote}{0}%
% SCAN MARKS (pp. 67--74): small ticks and dots stand in the outer margins of
% the scanned copy on most pages of this batch; not text, not transcribed.
\opage{67}værdiens Gyldighed, bliver ingen Dom mulig. Det drejer sig altsaa
ogsaa i Videnskaben ikke blot om Tanker og Forestillinger men ogsaa, og tilsidst
om Værdier og Vurderinger.

Denne Synsmaade er mest udførlig og energisk gjort gældende af \textsc{Heinrich
Rickert} i det allerede omtalte Skrift om Grænserne for den naturvidenskabelige
Begrebsdannelse. (Se især p. 587---613). Jeg er enig deri, at Værdisynspunktet
ikke er et helt nyt Synspunkt, der først skulde komme frem i Økonomi og Æstetik,
i Etik og Religionsfilosofi; det gør sig i Grunden gældende, fra først af og
uafbrudt, ved alle Tanker saavel som ved alle Ønsker. I min Fremstilling af det
menneskelige Tankeliv er jeg gaaet den Vej, at jeg, efterat have konstateret, at
de tre Elementer Emne, Form og Værdi gør sig gældende i enhver Tankeakt,
forudsætter en intellektuel Interesse, en Anerkendelse af Sandhedens Værdi som
den stadige Forudsætning og derefter undersøger de Tankeformer, som denne
Forudsætning fører til at anvende. Men denne Undersøgelse støder saa paa sin Vej
paa selve Begrebet Værdi, ikke blot, fordi det danner den stadige Forudsætning
for Søgen efter Sandhed, men ogsaa, fordi der er andre Værdier end de
intellektuelle. Derfor maa Værdibegrebet tages op i sin Almindelighed som et
Grundbegreb, en Kategori. Spørgsmaalet er, om Værdikvaliteten kan blive Genstand
for rationel Behandling, saaledes som Sansekvalitet og Existentialkvalitet har
vist sig at kunne\footnote{\textit{Den menneskelige Tanke}, p. 238 f.}.

Som ved alle Kvaliteter begynder vi tilsyneladende rent subjektivt. Værdidomme
udspringer af Lyst og Ulyst, af Tilfredsstillelse og Savn. Men Lyst og Ulyst er
Symptomer paa en Trang, en Bevægelse eller en Tendens i en vis Retning, og de
fremkommer, naar Tendensen enten begunstiges, saa at den kan fortsættes med fuld
Kraft, eller hæmmes og stanses. Vi støder her atter paa Begrebet Retning. Hvad
det saa er, vi iagttager, er det altid i Bevægelse i en eller anden Retning. Men
naar der er Tale om Værdier, forudsættes der, at der er en Helhed med Tendens
til at hævde sig overfor Omverdenen, og det kommer da an paa, om denne Tendens
fremmes eller hæmmes. Og i strengeste Forstand kan der kun være Tale om Værdier
fra personlige Væseners Standpunkt, fordi kun der Tendensen og dens Brydning med
Forholdene mærkes og føles. Kun ad Analogiens Vej overføres Værdibegrebet til
andre Helheder (Samfund, Organismer, uorganiske Helheder); det er jo allerede
Analogi, naar et personligt Væsen søger at sætte sig ind i, hvad der kan have
Værdi for andre personlige Væsener. En Tendens, som mærkes, kalder vi en Trang.
Der vil altsaa være et indre Forhold mellem Begreberne Trang og Værdi. Naar det,
Trangen gaar ud paa, bliver Genstand for Bevidsthed forud for
Tilfredsstillelsen, bliver Trangen til Drift\footnote{Om Forholdet mellem Trang
og Drift se min \textit{Psykologi}\textsuperscript{6}, VI B, 2 c; C, 1. --- VII
B, 1 a. --- Det er \textsc{Hobbes} og \textsc{Spinoza}, vi skylder den tydelige
Skelnen mellem Trang og Drift. Især har \textsc{Hobbes} sat Trangen i
Forbindelse med en begyndende Bevægelse i en vis Retning. Han ledes her af en
Analogi med \textsc{Galilei}'s Kraftbegreb (momento, impeto), der finder
Anvendelse, hvad enten et Legeme er i Hvile eller i Bevægelse. Jfr.
\textit{Elements of Law}, I, 7, 2 (endeavour ɔ: internal beginning of animal
motion). \textit{De Corpore}, XXV, 12 (Tendensen til uhæmmet motus vitalis
bliver, naar Erfaring og Erindring er mulig,
% NOTE RUNS ON (pp. 67--68): the note continues at the foot of p. 68 with "til Trang, appetitus)".
til Trang, appetitus). En rent fysisk Definition af conatus (som motus per
spatium et tempus minus qvam datur) gives De Corpore, XV, 2. (Se om dette sidste
Punkt \textit{Den nyere Filosofis Historie}\textsuperscript{2}, I, p. 174.
\textsc{Tønnies}: \textit{Thomas Hobbes}\textsuperscript{2}, p. 116 f). --- For
\textsc{Spinoza}'s Vedkommende henvises til \textit{Ethica}, III, 6---9, hvor
der skelnes mellem Tendensen (conatus) til Selvopholdelse, Trangen (appetitus)
og Driften (cupiditas), hvilken sidste forudsætter Bevidsthed om Maalet. I nyere
Literatur fremtræder den samme Forskel hos \textsc{Stout} (impulse --- desire),
\textsc{Fouillée} (appétit --- désir) o.~A. I tysk Literatur betyder Trieb ofte
begge Dele, ikke til Fordel for Klarheden.}.
% SIGNATURE MARK (p. 67): "41*" at the foot; printer's apparatus, not transcribed.

% --- p. 68 ---
\opage{68}En objektiv Værdilære bliver mulig ved, at Trang og Værdi viser sig at
hænge sammen med et personligt Væsens Bestaaen som Helhed. De viser sig at
variere med Betingelserne for denne Bestaaen. En Undersøgelse af disse
Betingelser belyser Værdiernes Opstaaen og Bestaaen, og den gør det muligt at
skelne mellem ægte og falske Værdier, begrundet eller ubegrundet Trang, idet den
vil kunne vise, om den Trang, der føles, og den Vurdering, den fører til, svarer
til de virkelig givne eller mulige Livsbetingelser.

Værdibegreberne kan inddeles efter forskellige Synspunkter. I formel Henseende
kan der skelnes mellem elementære (sansebestemte) og ideelle
(forestillingsbestemte) Værdier, mellem umiddelbare og middelbare Værdier
(Formaal og Midler), mellem potentielle og aktuelle Værdier. I real Henseende
kan der, naar Indholdet lægges til Grund, skelnes mellem biologiske,
intellektuelle, æstetiske, etiske og religiøse Værdier, og naar Omfanget lægges
til Grund mellem individuelle, sociale og
% PRINTED AS IS (p. 68): "og" doubled over the line end ("sociale og / og kosmiske"). A stray raised mark after "mellem" in the line above is not transcribed.
og kosmiske Værdier. Hvorvidt en rationel Værdilære kan gennemføres for alle
disse Værdiers Vedkommende, er en Sag for sig, og Spørgsmaalet derom hører under
Behandlingen af Vurderingsproblemet. I denne Sammenhæng kan vi holde os til de
reale Værdier, særlig efter Indholdssynspunktet, for at belyse Sammenhængen
mellem Værdibegrebet og Totalitetsbegrebet.

\medskip
25. At de forskellige reale Værdier, hvad Omfang angaar, forudsætter visse
Helheder med Tendens til at hævde sig, er øjensynligt. Hvert enkelt Individ,
hver social Gruppe og hvert Folk vurderer Forholdene udenfor sig og indenfor sig
efter, om Betingelserne, de virkelige eller de formodede, for Bestaaen ere
tilstede eller ikke. Men ogsaa fra Indholdssynspunktet kan det samme vises. Det
økonomiske Liv, det intellektuelle Liv, det æstetiske Liv, det etiske Liv og det
religiøse Liv --- det er ligesaa mange forskellige Helheder af Tanker og
Bestræbelser, hver med en Tendens til at afslutte sig i sig selv, og hver for
sig bestemte ved et herskende Formaal, med hvilket vedkommende Livsomraade staar
og falder.

\textit{a.} Kampen for det fysiske Livs Opholdelse kan gøre de økonomiske
(biologiske) Interesser til Et og Alt. Der er her en Side af Livet, som har sine
egne Love og i Kraft deraf afgrænser sig til en Helhed, selv om denne Helhed
% PRINTED AS IS (p. 68): „omfatttende“ (triple t).
igen kan gaa ind som Del i en mere omfatttende Livshelhed. Der foreligger her
ogsaa faktisk en egen Videnskab med sine egne Synspunkter og Metoder. (Se
nærmere Etik III, 14).

\textit{b.} At Videnskaben er et Exempel paa Helhedsdannelse, er forsøgt
godtgjort i foregaaende Afsnit af denne Afhandling. Det videnskabelige Arbejde
gaar ud paa
% --- p. 69 ---
\setcounter{footnote}{0}%
\opage{69}at opbygge en Tankeverden, indenfor hvilken ethvert Emne kan finde sin
Plads. Intellektuelle Værdier beror paa, hvad der fremmer dette Arbejde og fører
det i Retning af dets Maal.

\textit{c.} Æstetisk Værdi faar et Emne ved, at det fremtræder som en
karakteristisk Helhed og kan gøre sig gældende som saadan for umiddelbar
Anskuen, afset fra intellektuel og praktisk Interesse. Da Virkeligheden ikke
altid frembyder karakteristiske, fuldtud individualiserede Emner, kan Fantasien
træde til og give den Afrunding, uden hvilken et Hele ikke bliver til. Det er et
Forhold, der særlig kan studeres, hvor vi kender Emnerne (Motiverne) iforvejen
og kan sammenligne dem med, hvad Digteren eller Kunstneren føjer til gennem sit
Arbejde. Exempler herpaa har vi i de græske Tragikeres Forhold til Sagnhistorien
og i Shakespeare's Forhold til sine Forgængere (saaledes som dette er belyst i
\textsc{Schück}'s Bog om Shakespeare). Det er Billedet af en ejendommelig,
individuel Helhed der tilstræbes. Ved anvendt Kunst er Skønhedsindtrykket
betinget ved Forholdet mellem Genstandens Kontur og Ornamentets Kontur (mellem
„Konturlinie“ og „Ornamentlinie)“.
% PRINTED AS IS (p. 69): the closing parenthesis stands inside the closing quote, „Ornamentlinie)“.
For at der skal opnaas Skønhed, maa de nemlig enten falde sammen, eller ogsaa
maa deres Skæringspunkter ligge simpelt rytmisk, men dette er ogsaa Betingelsen
for, at et Helhedspræg kan opstaa. Ligesaa er det det simpelt rytmiske i Verset
i Modsætning til Prosa, der gør, at det fremtræder med større Helhedsvirkning,
og heri maa Grunden søges til den større æstetiske Værdi, Verset i sin
Begrænsning kan have\footnote{Denne Iagttagelse skyldes en ung Ven, der særlig
har beskæftiget sig med Ornamentik.}.

I filosofiske Systemer raader ofte en æstetisk Trang til harmonisk eller rytmisk
Sammenhæng, en Trang, der naturlig udspringer af den teoretiske Trang til at
paavise Sammenhæng i alt, selv i det mest modstridende. Nylig har \textsc{Mark
Baldwin} (Genetic Theory of Reality, 1915) hævdet, at en Verdensanskuelse, der
vil undgaa, eller føre ud over, al Relativitet, maa faa en æstetisk eller en
kunstnerisk Karakter; den maa blive „Pankalisme“. Ellers faar den en rent formal
Karakter, som i de Systemer, hvor Forholdet mellem Enhed og Mangfoldighed,
Identitet og Forskel spiller Hovedrollen. \textsc{Baldwin}'s Tankegang minder om
\textsc{Schelling}'s „System des transcendentalen Idealismus“ (1800), hvor
kunstnerisk Anskuen proklameres som eneste Mulighed til at naa ud over Livets
Modsætninger\footnote{\textit{Den nyere Filosofis Historie}\textsuperscript{2},
II, p. 156.}.

\textit{d.} Al etisk Vurdering gaar ud fra Interessen for en Helhed, hvis
Bestaaen og
% PRINTED AS IS (p. 69): "og" doubled over the line end ("Bestaaen og / og Udfoldelse").
og Udfoldelse skal hævdes og fremmes. Her kommer tydelig nok igen Omfangets
Synspunkt frem, idet Brydningen mellem etiske Systemer betinges ved, at det er
forskellige, mere eller mindre omfattende Helheder, som ligger til Grund. Det
kan være det enkelte Individs, eller Familiens, eller Standens, eller Nationens,
eller Menneskehedens Bestaaen, der betinger Vurderingen i det enkelte. I tredie
Kapitel af min Etik har jeg undersøgt disse Modsætninger og deres Betydning for
Etikens videnskabelige Karakter.

Undertiden mener man, at en Vurdering bliver objektiv blot derved, at den ikke
er individuel, men social. Deraf den hyppige Paastand, at en videnskabelig
% --- p. 70 ---
\setcounter{footnote}{0}%
\opage{70}Etik selvfølgelig maa være en Samfundsetik. \textsc{Heinrich
Maier}\footnote{\textit{Psychologie des emotionalen Denkens}, p. 652, jfr.
654---673.} vil begrænse Muligheden af objektiv Vurdering til saadanne Tilfælde,
hvor en „Relation til Mennesket som genus“ ligger til Grund. Men i sig selv kan
en Vurdering ud fra en mere begrænset Helhed være ligesaa konsekvent og ligesaa
objektivt begrundet som en Vurdering ud fra den mest omfattende Helhed, vi kan
forestille os. Egoisten kan f.~Ex. have fuld objektiv Ret i, at denne eller hin
Adfærd fra hans Side er en nødvendig Betingelse for Opnaaen af hans Formaal. Der
er altid i en etisk Dom en af Præmisserne, som bestemmes ved Forholdet til den
bestemte Helhed, hvis Bevarelse og Udvikling det gælder. Og det er denne Præmis,
der volder Problemet; ti der er ingen logisk Nødvendighed for, at en vis bestemt
Helhed skal bestemme Vurderingen.

En sammenlignende Etik vil have at undersøge, hvorvidt de Vurderinger, der
anstilles ud fra Forholdet til forskellige Helheder som Grundlag, kan føre til
Resultater, der stemmer overens (afset fra Motiveringen). Ved en saadan
Undersøgelse vil det vise sig, at der er visse etiske Kategorier, som anvendes
paa alle etiske Standpunkter, selv om de faar forskellig Plads. Saadanne
Begreber er Pligt, Dyd, Gode, Ret. De udtrykker alle ét Forhold mellem en enkelt
Handling og den Helhed, ud fra hvilken den vurderes. I Pligten udtrykkes, at
Helhedens Krav bestemmer Handlingen, i Dyden, at Handlingen udspringer af
umiddelbar Enhedsfølelse med Helheden; et Gode er alt, hvad der tjener Helhedens
Bestaaen; en Ret udtrykker, at Handlingen ikke er betinget ved Helhedens Krav,
men heller ikke strider mod dette\footnote{Smlgn. \textit{Den menneskelige
Tanke}, p. 350---361.}.

\textit{e.} Religionsfilosofien fører til en tilsvarende Opfattelse. Al Religion
betinges og bæres af en Trang til at fastholde den stadige Mulighed for, at den
Helhed, til hvilket Mennesket føler sig knyttet, eller som det selv mener at
være, kan bestaa og udvikle sig. Det er det, der menes, naar Religion siges at
være Tro paa Værdiens Bestaaen. Det er stadig et Rige, et større eller et
mindre, som Mennesket hævder i sin Religion\footnote{Smlgn., foruden min
\textit{Religionsfilosofi}, \textit{Den menneskelige Tanke}, p. 375 f.}.

\medskip
26. Overalt, hvor der i vor Erfaring fremtræder Emner med Helhedspræg, kan vi
have Interesse, i hvert Tilfælde en intellektuel Interesse, af at undersøge
deres Oprindelse og Betingelserne for deres Bestaaen og Udvikling. Det kan ofte
være vanskeligt her at afgrænse de rent intellektuelle Interesser fra andre
Interesser. Æstetisk kan vi saaledes „føle os ind“ eller „leve os ind“ i
Naturgenstande eller i historiske Personligheder og Begivenheder, som ikke har
nogen praktisk Interesse for os. Da al Værdi er knyttet til en Helhed som
Forudsætning, vil det ogsaa naturlig kunne ske, at en hvilkensomhelst Helhed
fremtræder for os som Værdicentrum, idet den ad Analogiens Vej faar et Skær af
Interessen for den Helhed, vi selv er eller tilhører.

% --- p. 71 ---
\setcounter{footnote}{0}%
\opage{71}I de moderne Udviklingsteorier, saaledes som de er fremtraadte i de to
sidste Aarhundreder, gør dette Synspunkt sig tydelig gældende. Forsaavidt
Menneskets Stilling i Universet er bestemt ved, hvad Skæbne det organiske Liv,
Jordkloden og Solsystemet har haft og har, er det ikke blot Analogi, der
betinger Interessen for biologisk, geologisk og astronomisk Udvikling. Men der
knytter sig ingen saadan Interesse til en enkelt Stjernes, et enkelt Bjergs
Historie eller til en Egns Geologi, naar vi ikke selv bor indenfor deres
Grænser. Og paa lignende Maade forholder det sig med Interessen for andre Tider
og Folkeslag, og for Personligheder, der har levet under helt andre Forhold end
vore. \textsc{Rickert} har her opstillet det interessante Begreb „historisk
Centrum“, et Begreb, der saa at sige staar paa Overgangen mellem det rent
teoretiske Totalitetsbegreb og Værdibegrebet. Historieforskningen vælger
saadanne Emner, der frembyder Karakteren af en Helhed, idet de staar baade som
Samlingssted og som Udgangspunkt for hele Rækker af Begivenheder. En saadan
Helhed --- være sig en enkelt Personlighed, et Folk, en Institution, en Kultur
--- staar som værdifuld i Lyset af, hvad der er samlet og afsluttet i den, og
hvad der udgaar fra den, selv om den hverken i sine Aarsager eller i sine
Virkninger har praktisk Indflydelse paa den Helhed, Historikeren selv er eller
lever i. Historikeren maa kunne sætte sig ind i, hvad der, givet en vis Helhed
--- en vis Tid, et vist Folk, visse Opgaver --- faar Værdi, og det er det, der
for ham bestemmer Forholdet mellem de enkelte Dele af den Fremstilling, han
giver. Det er ingen egentlig Vurdering, Historikeren giver (som Historiker), men
til Grund for hans Betragtning ligger et Forhold til Værdier (Beziehung auf
Werte), der ikke behøver at være hans egne, men som han kan sætte sig ind
i\footnote{\textit{Die Grenzen der naturwiss.
Begriffsbildung}\textsuperscript{2}. p. 325---329.}.

Mangler dette Forhold til Værdier ganske, hører den historiske Interesse op, og
vi kommer da til, hvad en anden skarpsindig Historiefilosof, \textsc{Georg
Simmel}, har kaldt „den historiske Tærskel“\footnote{\textit{Probleme der
Geschichtsphilosophie}\textsuperscript{2}, p. 131.}. Denne Tærskel kan
naturligvis flytte sig, sænkes eller hæves, ligesom \textsc{Herbarts} og
\textsc{Fechners} psykologiske Tærskel, alt efter som det historiske Materiale
gør det muligt at anlægge Helhedssynspunkter eller ikke.

Et tredie historisk Begreb hænger nøje sammen med de to nævnte Begreber, nemlig
Begrebet om historisk Storhed. Denne falder ikke sammen med moralsk Storhed, men
bestemmes ved Omfanget og Rækkeviden af de Kræfter, der virker til Frembringelse
af en ny Helhed i national eller kulturel Henseende. Saa voldsomt og selvisk en
historisk Personlighed kan vise sig at være, saa vil der dog --- ved virkelig
Storhed --- være hvad \textsc{Burckhardt} kalder „en hemmelighedsfuld Koincidens
mellem Individets Egoisme og Helhedens Nytte eller Storhed“. \textsc{Alexander}
gjorde ikke sine Erobringer og stiftede ikke sit Verdensrige for Humanitetens og
Kulturens Skyld, men hans Gerning førte dog ud over de Grænser,
Nationalitetsforskellighederne hidtil havde afstukket i disse Henseender. Store
Mænd er ifølge \textsc{Burckhardt} nød-%
% PRINTED AS IS (pp. 71--72): p. 71 ends "nød-" and p. 72 opens with the whole
% word "nødvendige"; the syllable is set twice. Hyphen kept, both halves as printed.
% --- p. 72 ---
\setcounter{footnote}{0}%
\opage{72}nødvendige, for at „den verdenhistoriske Bevægelse periodisk med et
Ryk kan frigøre sig fra uddøde Livsformer og reflekterende
Snak“\footnote{\textit{Weltgeschichtliche Betrachtungen}, p. 245, 252. --- Jfr.
\textsc{Hegel}'s aandfulde Udvikling om „Verdensfornuftens List“ i at benytte
menneskelige Lidenskaber, i Indledningen til hans Historiens Filosofi.}.

De forskellige historiske Skoler vilde ofte være uenige om, hvor de historiske
Centra ligger, hvor den historiske Tærskel er at sætte, og hvorvidt Begrebet
historisk Storhed skal finde Anvendelse i det enkelte Tilfælde. Der kan være en
Tendens til at holde med de Overvundne, eller til at holde med de Sejrende; der
kan lægges Vægt paa Kulturen eller paa Staten, paa de spontant virkende Kræfter
eller paa Organisation gennem Anvendelse af Magt. Under alle disse Modsætninger
vil det stedse vise sig at være Forholdet mellem Helheder og Betingelser for
disses Bestaaen, der ligger til Grund for de forskellige Synsmaader. Og ad
Analogiens Vej vil disse Modsætninger tilsidst, men ogsaa først tilsidst, kunne
føres tilbage til Modsætninger mellem etiske Standpunkter.

\medskip
27. Gennem Totalitetsbegrebet viser Værdibegrebet tilbage til Naturvidenskab og
Historieforsken. Ti dersom der ikke i Tilværelsen, saaledes som denne nu engang
er beskaffen, havde været Mulighed for Helheders Opstaaen, --- dersom altsaa
alle Aarsagsrækker havde været divergerende eller parallele, --- saa vilde
Værdibegrebet ikke kunne være opstaaet. Det er et Bidrag til Tilværelsens
Karakteristik, at der kan opstaa og bestaa Totaliteter indenfor den.

Hvorledes de Totaliteter, hvis Existens Værdien forudsætter, er opstaaede,
behøver ikke at være afgørende for selve Værdien. Værdi er en Kendsgerning, der
ikke omstødes ved den Maade, paa hvilken den er bleven til, ligesaa lidt som et
Menneskes Værdi forringes ved hans Forfædres slette Færd. Selv en saa uhildet
Mand som \textsc{Guyau} mente dog, at et Værdifuldt ikke kan opstaa af noget,
der ingen Værdi har eller endog har det modsatte af Værdi. Det er en teologisk
Tankegang, der her gør sig gældende, og som konsekvent maatte føre til, at al
Forsken efter Værdiers Tilblivelse blev forbudt i Værdiens eget Navn, ligesom
Lohengrin forbyder Elsa at spørge om hans Navn og Herkomst. Men det er ikke let
at hæmme Trangen til at spørge og forske, den Trang, som \textsc{Wagner} kalder
„Wissens Sorge tragen“\footnote{%
% VERSE IN NOTE (p. 72): four lines of Lohengrin, set in the print as a
% block beside the note numeral; line ends kept with \\ inside the footnote.
Nie sollst du mich befragen,\\
noch Wissens Sorge tragen,\\
woher ich kam der Fahrt,\\
noch wie mein Nam' und Art.}. Der er heller ingen Grund til nogen Betænkelighed.
Værdier opstaar ifølge de store Love, der gør Helhedsdannelser mulige under
visse Betingelser, og netop derved hænger, som allerede bemærket, Værdibegrebet
sammen med den videnskabelige Erkendelse og kan selv blive Indhold af Videnskab.

Men Virkelighedens Verden er mere omfattende end Værdiens Verden. Her er
Erfaringer at gøre, som kan begunstige en Følelse af Humor overfor Verden, som
den er, --- dels fordi Værdier ofte virkeliggøres ad Veje, vi ikke kan tillægge
% --- p. 73 ---
\setcounter{footnote}{0}%
\opage{73}Værdi i og for sig, dels fordi Værdiens Verden fremtræder som en Oase
indenfor en Ørken\footnote{\textit{Den store Humor}, p. 70---73.}. Selv om man
kunde tænke sig et værdifuldt Maal for Virkelighedens uoverskuelige
Aarsagsrækker, vilde vi dog ikke kunne kende dette Maal, og vi maatte i vor
Vurdering af Tilværelsen holde os til de Dele og de Strækninger, vi kan
overskue. Den Opfattelse, at alt i Naturen og Historien fører i Retning af en
Fuldendelse, en Harmonisering, en Virkeliggørelse af alle Idealer, skyldes
kiliastiske Forestillinger og optimistiske Postulater. En saadan naiv Teleologi
har ingen fast Fod under Fødderne. Støtter man sig til Erfaringen, maa man
snarere give \textsc{Burckhardt} Ret i, at i Historien som i Naturen anvendes
der uhyre Foranstaltninger og en uforholdsmæssig Larm med smaa Resultater til
Slut\footnote{\textit{Weltgeschichtliche Betrachtungen}, p. 183.}. Der er
Ødselhed i Historien som i Naturen. Men Værdibegrebet som Kategori har heller
ikke teleologiske Betragtninger som Konsekvens. Ligesaa lidt som nogen anden
Kategori, kan Værdibegrebet bruges som eneste Grundbegreb, ud fra hvilket alle
andre Tankeformer skulde kunne udledes i deres relative Berettigelse. Det er
endog den mest konkrete og komplexe af vore Kategorier, forudsætter de andre,
uden at forudsættes af dem.

Modsætningen mellem menneskelige Værdier og Virkelighedens Verden kom til klar
Bevidsthed ved det nye Verdensbillede, \textsc{Kopernicus} gav Stødet til.
Problemet blev stillet saa skarpt som ingensinde før: Livet er forsvindende i
Forhold til Jordmassen og Jorden i Forhold til Stjerneverdenen. Selv om Livet
sejrede over Materien, saa at det „élan de vie“, som en aandfuld Filosof taler
om, stedse kunde holde Stand mod Tendensen til Fordeling og Ligevægt, saa
begynder Disharmonierne igen indenfor Livets Verden, idet den ene Livshelhed kun
bestaar paa den andens Bekostning, som Ulven overfor Lammet. Hvis hver
Livshelhed skal betegne et Naturformaal, er der Strid i selve Værdiernes Verden.
Og hvis man istedetfor Teleologi blot taler om Orthogenesis (§~20), saa viser
Erfaring, at Udviklingen kan løbe sig fast i ubehjælpsomme, ikke levedygtige
Former.

Kunde det endda vises, at alle Sammenstød mellem Helheder indbyrdes, eller
mellem Helheder og Elementer var frugtbare, tjente til at føre Udviklingen
videre, vilde Værdibegrebet trods alt kunne hævde en afsluttende Stilling i vor
Tankeverden. Men har \textsc{Goethe} ikke Ret overfor \textsc{Hegel}, naar han
hævder, at der gives en absolut Tragik, ikke blot Konflikter, gennem hvilke
Sandhed og Ret fremtræder med ny Klarhed?

Vor Verdens- og Livsanskuelse maa, hvis den vil bevare Sammenhængen med rationel
Tænken, gaa andre Veje end Teleologien. Om der gives saadanne kan Kategorilæren
ikke for sig alene afgøre.
% END OF CH. V (p. 73): a short centred rule below the notes closes the
% chapter. SIGNATURE LINE at the foot: "D. K. D. Vidensk. Selsk. Skr., 7. Række,
% hist. og filos. Afd. III, 2" and the sheet numeral "42"; printer's apparatus,
% not transcribed.

% --- p. 74 ---
% CHAPTER HEAD (p. 74) as printed at the head of a new page: two centred lines
% ("VI." / title with final stop) over a short rule. The Indhold (p. 80) gives
% the same title. First division indented.
\section*{VI.\ Totalitet som Grænsebegreb.}
\addcontentsline{toc}{section}{VI.\ Totalitet som Grænsebegreb}

\opage{74}28. Udtrykket Grænsebegreb er indført af \textsc{Kant} til at betegne
et Begreb, vi nødvendigvis kommer til, men som vi ikke kan fuldbyrde, ikke give
positiv Betydning. Et saadant Begreb fandt \textsc{Kant} i Begrebet „Tingen i
sig selv“. Da der blandt vor Erkendelses Forudsætninger nødvendigvis findes en
subjektiv Faktor, der lægger sig for Dagen i Tænkningens Grundformer, kan vi
ikke erkende Tilværelsen, som den er i sig selv, og dog maa vi, mener
\textsc{Kant}, nødvendigvis antage, at Tilværelsen har en Beskaffenhed, ganske
afset fra den Maade, paa hvilken vi efter vor Tænknings Natur opfatter den.
Denne Paastand var, som bekendt, en stor Anstødssten for \textsc{Kant}'s
nærmeste Efterfølgere. Det vilde jo ogsaa være mærkeligt, om Erkendelsen selv
under sin logiske Fremskriden skulde føre til et Punkt, hvor den selv ikke mere
skulde gælde. Et saadant Punkt fandt \textsc{Kant} i det Spørgsmaal om Grunden
til de Emner (det „Stof“), under hvis Bearbejdelse den menneskelige Tanke faar
Brug for sine Former, sine Forudsætninger. Men --- netop ifølge \textsc{Kant}'s
egen Filosofi kan et saadant Spørgsmaal ikke opkastes! Ti ifølge „Kritik der
reinen Vernunft“ har Tanken ingen anden Opgave end den at bearbejde givne Emner
og forklare det ene Emne ved det andet. Hvor intet Emne er givet, kan der kun
gives formal Videnskab (Logik og Matematik), som netop ser bort fra alle
specielle Emner og undersøger Erkendelsens rene Former. I real Videnskab
(Naturvidenskab, Psykologi og Historie) bearbejdes Emner saaledes, at den
rationelle Sammenhæng, \textsc{Kant} kalder „Erfaring“, bliver mulig. Hverken
formal eller real Videnskab støder paa Tingen i sig selv, hin fordi den
overhovedet ingen „Ting“ forudsætter, denne, netop fordi den stedse forudsætter
„Ting“ som \emph{givne}.

Men \textsc{Kant} har et andet Grænsebegreb foruden „Tingen i sig selv“, og det
er Begrebet om absolut Totalitet. Dette Begreb betegner Noget, der ikke kan være
Emne for os, da ethvert Emne er begrænset, enhver given Helhed staar overfor en
Omverden, med hvilken den staar i Vekselvirkning. Begrebet absolut Totalitet har
efter \textsc{Kant} den Betydning at pege ud over enhver Grænse, man kunde være
tilbøjelig til at stanse ved, --- at henvise til nye Opgaver, hvergang Arbejdet
med de hidtil foreliggende Emner synes at være fuldendt. Ogsaa dette Begreb
kommer vi ifølge \textsc{Kant} med Nødvendighed til, da vor Tankes Grundform er
Syntesen, der stedse kræver det, der hidtil er sammenfattet, sammenfattet med
nye Elementer.

Ved nærmere Betragtning viser det sig nu, at den Grænse, der ligger i Begrebet
absolut Totalitet (i hvad \textsc{Kant} kalder „Ideer“), falder sammen med den
Grænse, Begrebet „Ting i sig selv“ skulde angive. Hvor \textsc{Kant} kritiserer
det dogmatiske Begreb om „Verden“ som given absolut Helhed, udtaler han: „I den
empiriske Regres kan der ikke træffes nogen Erfaring om en absolut
Grænse~.~.~.~. Grunden dertil er, at en saadan Erfaring maatte indeholde en
Begrænsning af et
% --- p. 75 ---
\setcounter{footnote}{0}%
% QUOTES (p. 75): the closing mark after umuligt answers the opening mark on
% p. 74; the fragment alone counts one closing mark over.
\opage{75}Fænomen ved Intet,... hvilket er umuligt“\footnote{\textit{Kritik der
reinen Vernunft}\textsuperscript{1}, p. 517.}. Denne Udtalelse rammer „Tingen i
sig selv“ (der aldrig kan blive Fænomen) saavel som den absolute Totalitet (der
aldrig kan foreligge som given). Og det er klart, at til „Tingen i sig selv“ kunde
vi kun komme, hvis vor Erkendelse var fuldendt, dannede en absolut Totalitet;
først da kunde vi komme til at spørge om Udspringet for de Emner, vi havde
arbejdet med; indtil da havde vi nok at gøre med at sammenarbejde de nye Emner,
der frembød sig. Men saa staar vi jo netop ved den anden Grænse, ved det Ideal,
Begrebet absolut Totalitet stiller, og stiller paa ethvert Trin af vor Erkendelses
Udvikling. Stedse er Opgaven den samme: ud over enhver given Helhed at søge
en større Sammenhæng, en større Helhed, og indenfor enhver given Helhed at søge
tilbage til simplere Elementer, som saa igen vil vise sig at være Helheder i
det smaa.

I \textsc{Kant}'s „Antiteser“, som han forgæves søger at paadutte en dogmatisk
Tankegang, er i Virkeligheden hans egen Filosofi indeholdt. Der opereres i
„Antiteserne“ med \textsc{Kant}'s eget centrale Begreb om „mulig Erfaring“, og
han siger udtrykkelig, at de udtrykker det rent videnskabelige
Standpunkt\footnote{\textit{Kritik der reinen Vernunft}\textsuperscript{1}, p.
468 f.}. De gør intet absolut afsluttet System muligt; men det er jo ogsaa den
egentlige Konsekvens af \textsc{Kant}'s Filosofi. Saalænge Tanken er vaagen,
finder der en stadig Brydning Sted mellem et Givet og de Opgaver, det stiller.
Erkendelsen arbejder stedse paa at komme til at udgøre en fuldkommen Helhed, men
saalænge Forholdet mellem Emner og Form gentager sig, stiller Opgaven sig paa
ny. ---

Som det i de foregaaende Afsnit af denne Afhandling er paavist, er
Totalitetsbegrebet ikke blot en speciel Kategori, der har sin bestemte Plads i
Kategoriernes Række men det udtrykker tillige en Ejendommelighed ved enhver af
vore Tankeformer. Derfor kan \textsc{Kant}'s skarpe Sondring mellem Kategori og
Ide heller ikke fastholdes\footnote{\textit{Kontinuiteten i \textsc{Kant}'s
filosofiske Udviklingsgang}, p. 29---35.}.
% p. 75, note 3: the print has „KANT's“ with a small raised apostrophe
% (second reading 2026-09-30, confirmed at 800 dpi; the work's own title
% has „Kants“). Earlier UNSURE reading „KANTs“ withdrawn.

\medskip
29. Ser vi nu nærmere paa Totalitetsbegrebet som Grænsebegreb, frembyder
% PRINTED AS IS (p. 75): „pet“ for „det“ (wrong sort; found in the second
% reading 2026-09-30, clear at 400--500 dpi).
pet, hvad vi ogsaa maatte vente hos et Grænsebegreb, dels en negativ, dels en
positiv Side.

\textit{a.} Den negative Side, som fremkommer ved, at enhver Helhed ---
Tankehelhed eller Emnehelhed --- bestemmes ved sit Forhold til noget Andet end
den selv, giver igen Anledning baade til en formal (logisk) og en real
(empirisk) Betragtning.

% ENUMERATOR (p. 75): the second „a.“ is read as Greek α, since β. follows
% on p. 76 and b. (Latin) on p. 77.
α. Intet Begreb kan bestemmes uden ved sit Forhold til andre Begreber,
ingen Dom begrundes uden ved sit Forhold til andre Domme. Ethvert Begreb og
enhver Dom er et Punkt, hvor Tankens ene Passerben sættes, for at der saa kan
søges et Punkt, hvor det andet kan sættes. Dette er det, som fører til at
opstille Relation (resp. Syntese) som det første Led i Kategoriernes Række.
Derfor vil en
% SIGNATURE MARK "42*" at the foot of p. 75 omitted as printer's apparatus.
% --- p. 76 ---
\setcounter{footnote}{0}%
\opage{76}absolut Afslutning stride mod Tankens egne Love. Der endes i hvert
Tilfælde med et Problem, naar det andet Passerben intet Punkt kan finde.

Naar spekulative Systemer lærte en Afslutning i Tanken om Noget, der „bestaar
i sig selv og forstaas ved sig selv“, saa var det en Overgang fra Tænkning til
Mystik. Intet bestaar ved sig selv, og Intet forstaas ved sig selv. Spekulative
Filosofer, lige fra \textsc{Platon} og \textsc{Aristoteles} af, giver os Valget
imellem at anerkende Noget, som forstaas ved sig selv og bestaar ved sig selv,
og at ende i Skepticisme. Saaledes siger \textsc{Spinoza} (Den korte Traktat I,
7, 9), at naar man ikke kan erkende det, der ligger til Grund for Alt, kan man
heller ikke erkende Enkelttingene, som har deres Grund i det. Selv i nyeste Tid
kommer en saadan ontologisk Tankegang frem. \textsc{Bertrand Russell} lægger den
til Grund for sin Kritik af den Opfattelse, ifølge hvilken Kriteriet for Sandhed
er den indbyrdes Sammenhæng mellem vore Tanker og Iagttagelser. Resultatet af
denne Opfattelse bliver Skepticisme, paastaar
% MARK (p. 76): a stray tick stands before „RUSSELL:“ and a stray dash in
% the left margin by „at Sandheden maa vindes“; not transcribed.
\textsc{Russell}: ti en fuldkommen og altomfattende Sammenhæng, kan ikke
paavises, og vi kunde derfor, ifølge hint Sandhedskriterium, ikke være sikre paa
Gyldigheden af en eneste Dom\footnote{\textit{On the nature of truth.}
(Proceedings of the Aristotelian Society. 1906---07). Smlgn. ovenfor § 18 a.}.
Men den Opfattelse, \textsc{Russell} bekæmper, siger i Grunden kun, at Sandheden
maa vindes ved stadigt Arbejde, og at de Sandheder, vi allerede mener at have
vundet, stedse paany maa prøves gennem deres Forhold til nye Tanker og
Iagttagelser. Her er et Arbejde, der ikke kan stanse. Og det fortsatte Arbejde
hviler paa den stadige Forudsætning, at Sandheden maa udgøre en Helhed. Selv de
sikreste Resultater, som ere opnaaede, viser sig at være begrænsede og at trænge
til Supplering. Der kan opnaas en stedse inderligere Sammenhæng mellem et stedse
større Indbegreb af Iagttagelser. Men der vil aldrig kunne opnaas begrundet
Forvisning om, at man nu er ved „Noget, der bestaar i sig selv og forstaas ved
sig selv“. Og --- selv om det var muligt at naa saa vidt, saa vilde Sandheden jo
ogsaa da bestaa i en Sammenhæng, nemlig Sammenhæng med det i sig selv
Bestaaende, og der kunde da være Mulighed for adskillig Tvivl, naar denne
Sammenhæng skulde staa sin Prøve i det Enkelte.

β. Saavidt den formale Betragtning. Men hvad der gælder for alle Tanker og
Iagttagelser, gælder ogsaa for alle Emner. Ethvert Emne, der kan være Genstand
for Iagttagelse og for tænkende Bearbejdelse, viser sig at være underlagt
bestemte Betingelser og at være bundet til bestemte Grænser, saa at man for at
forstaa det, maa tage Hensyn til hele den Sammenhæng, i hvilken det forekommer.
Ethvert Emne forholder sig til en Omverden, der enten frembyder Støtte eller
Hindring for dets Bestaaen. Jo mere Videnskaben gaar over til at anvende
Indsigten i de almene Love til at forstaa individuelle Helheder og finde
Betingelserne for deres Bestaaen, desto mere vil denne Synsmaade vise sig i sin
Berettigelse og i sin Betydning. Helt op i de højeste Former for det
menneskelige Aandsliv kan denne Undersøgelse finde Opgaver. ---

Erkendelsesproblemets Rod ligger i den Omstændighed, at Sandheden, baade som
real og som formal Sandhed, maa udgøre en Helhed. Istedetfor den gamle
% --- p. 77 ---
\setcounter{footnote}{0}%
% NOTE (p. 77, note 1): comma after (1912) at the image; the layer reads a stop.
\opage{77}Modsætning mellem „det Absolute“ og „det Relative“ maa man nu, som det
med Rette er sagt\footnote{\textsc{Léon Brunschvicg}: \textit{Les Étapes de la
Philosophie mathématique} (1912), p. 414.}, sætte Modsætningen mellem det Totale
og det Partielle.

\textit{b.} Den positive Side ved Undersøgelsen af en Helhed angaar for det
første Forholdet mellem Enhed og Mangfoldighed indenfor den. Enheden lægger sig
for Dagen i den gennemgaaende Lovmæssighed, ifølge hvilken Delene høre sammen,
--- Mangfoldigheden i de Elementer, der er forbundne i Kraft af denne Lov, men
hvert for sig kan besidde Egenskaber, der ikke blot er bestemte ved den. Som det
har vist sig i det foregaaende, frembyder de forskellige Helheder store
Uligheder, hvad Forholdet mellem Enhed og Mangfoldighed angaar. Og de filosofiske
Systemer udviser her en af deres mest karakteristiske Forskelligheder, alt
eftersom Vægten lægges paa det sammenholdende Baand eller paa de mangfoldige
Elementer, det forbinder. I Filosofiens Historie kommer denne Forskellighed
tydelig frem, naar man sammenligner \textsc{Parmenides} med \textsc{Demokrit},
\textsc{Platon} med \textsc{Aristoteles}, \textsc{Spinoza} med \textsc{Leibniz},
\textsc{Hegel} med \textsc{Herbart}, \textsc{Bradley} med \textsc{James}.
Systemernes Vanskelighed kommer, som \textsc{Wundt} træffende har gjort
opmærksom paa (se overfor § 10), deraf, at et enkelt Begreb forsøges gjort til
Et og Alt, uagtet det til sit Korrelat har et andet Begreb. Saaledes er netop
Forholdet mellem Enhed og Mangfoldighed. I den religiøse Forestillingskreds, som
tilsidst uvilkaarlig følger samme Love som alt andet menneskeligt Tankeliv, fører
det omtalte Modsætningsforhold til Dannelsen af to forskellige
Afslutningsbegreber, som hvert for sig vilde forudsætte Muligheden af absolut
Totalitet: Begrebet Gud og Begrebet Verden. Det egentlige Afslutningsbegreb
maatte her være et Begreb, der kunde forbinde begge disse Begreber --- ligesom
det hos \textsc{Platon} vilde være et Begreb, der kunde forbinde Ideer og
Fænomener, og hos \textsc{Spinoza} et Begreb, der kunde forbinde Substans og
Modi\footnote{Smlgn. min \textit{Religionsfilosofi}\textsuperscript{2}, p.
59---62.}. Hvis der skulde kunne dannes et saadant Begreb, maatte det være et
typisk Individualbegreb. Ti den Tilværelse, i hvilken vi lever, staar som et
Unicum, et stort Individ, hvis karakteristiske Ejendommelighed er bestemt ved de
faktisk gældende Loves Indhold og ved de faktisk givne Elementers Beskaffenhed.
\textsc{Spinoza} er maaske den, der klarest har set dette. Han hævder nemlig, at
Gud (deus s. substantia s. natura) ikke kan gaa ind under noget Almenbegreb (Ep.
50), og at hele Universet er et eneste stort Individ (Ethica II. Lemma 7.
Schol.). Det er derfor, at han i sin Erkendelseslære ikke vil blive staaende ved
de almene Love, men kræver en scientia intuitiva (en syntetisk Intuition), det
vil sige, en Erkendelse, som paa én Gang er universel og individuel. Ethvert
Forsøg paa en videnskabelig Verdensanskuelse gaar ud paa at danne et typisk
Individualbegreb, der indeholder en Karakteristik af Tilværelsen, denne
Tilværelse, som paa en Gang afgiver Emne for vor Tanke og selv er en Helhed, i
hvilken selve vor Tanke er en Del. Her ligger Roden til det kosmologiske
(metafysiske) Problem. ---

Foruden Modsætningen mellem Enhed og Mangfoldighed gør der sig for enhver
Helheds Vedkommende en Modsætning gældende mellem Værdi og Virkelighed,
% --- p. 78 ---
\opage{78}idet de indre Elementer og de ydre Omstændigheder baade kan støtte og
hæmme dens Bestaaen. Helhedens Bestaaen kan især bero paa de indre Elementers
Harmoni og Energi; men Spiren til Opløsning kan ofte netop ligge i disse
Elementers Beskaffenhed og indbyrdes Forhold. Og ligeledes forholder det sig med
de ydre Omstændigheder. Værdi beror paa de indre Elementers og de ydre
Omstændigheders Forhold til Helheden, naar dennes Bestaaen betragtes som
Grundværdi. Da ogsaa det Værdifulde er en Virkelighed, er det egentlig en
Modsætning mellem en værdifuld Virkelighed og al anden (indifferent eller
hæmmende) Virkelighed, vi her har for os. Spørgsmaalet bliver her, om vi kan
danne et Begreb, der kan forbinde begge til en Helhed. Her ligger Roden til det
religiøse Problem.
% END OF THE TREATISE'S TEXT: p. 78 ends with „religiøse Problem.“ over a
% short centred rule; the rest of the page is blank. (A stray dash stands in
% the left margin below the text; the layer's „c<“ at the foot answers to no
% visible character.) Not transcribed.
\begin{center}
\rule{5em}{0.4pt}
\end{center}

\clearpage
% --- p. 79 ---
% KATEGORITAVLE (p. 79), no folio printed. Heading in capitals, centred,
% over a centred parenthesis and a short rule. The table set as printed with
% a four-column tabular (roman numeral | arabic numeral | letter | text):
% group heads bold („(Værdibegreber).“ roman), the right brace spanning e.--h.
% with „Rationale Rækker.“ to its right set by a nested tabular inside
% $\left.\ \right\}$. A double rule closes the page. Show-through from the
% facing pages is visible on the scan; not transcribed.
\section*{KATEGORITAVLE}
\addcontentsline{toc}{section}{Kategoritavle}

\begin{center}
\opage{79}(Smlgn. Den menneskelige Tanke 1910)\\[0.5ex]
\rule{4em}{0.4pt}

\bigskip
\begin{tabular}{@{}r@{\ }r@{\ }r@{\ }l@{}}
I. & \multicolumn{3}{@{}l}{\textbf{Fundamentale Kategorier.}}\\
 & 1. & \multicolumn{2}{@{}l}{Syntese --- Relation.}\\
 & 2. & \multicolumn{2}{@{}l}{Kontinuitet --- Diskontinuitet.}\\
 & 3. & \multicolumn{2}{@{}l}{Lighed --- Forskel.}\\
 & & a. & Kaotisk Forskelsrække.\\
 & & b. & Ubestemt varierende Forskelsrække.\\
 & & c. & Regelmæssig varierende Forskelsrække.\\
 & & d. & Identisk varierende Forskelsrække.\\
 & & \multicolumn{2}{@{}l}{$\left.\begin{tabular}{@{}r@{\ }l@{}}
   e. & Fremskridende Forskelsrække.\\
   f. & Partiel Identitetsrække.\\
   g. & Gensidighedsrække.\\
   h. & Absolut Identitetsrække.
   \end{tabular}\ \right\}$\ \ Rationale Rækker.}\\[2ex]
II. & \multicolumn{3}{@{}l}{\textbf{Formale Kategorier.}}\\
 & 1. & \multicolumn{2}{@{}l}{Identitet.}\\
 & 2. & \multicolumn{2}{@{}l}{Analogi (Reduktion af Kvalitetsforhold).}\\
 & & a. & Tid.\quad b.\ Tal.\quad c.\ Grad.\quad d.\ Sted.\\
 & 3. & \multicolumn{2}{@{}l}{Negation.}\\
 & 4. & \multicolumn{2}{@{}l}{Rationalitet.}\\[2ex]
III. & \multicolumn{3}{@{}l}{\textbf{Reale Kategorier.}}\\
 & 1. & \multicolumn{2}{@{}l}{Kausalitet.}\\
 & 2. & \multicolumn{2}{@{}l}{Totalitet.}\\
 & 3. & \multicolumn{2}{@{}l}{Udvikling.}\\[2ex]
IV. & \multicolumn{3}{@{}l}{\textbf{Ideale Kategorier} (Værdibegreber).}\\
 & 1. & \multicolumn{2}{@{}l}{Formale Værdiforhold.}\\
 & 2. & \multicolumn{2}{@{}l}{Reale Værdiforhold.}
\end{tabular}

\bigskip
\rule{5em}{0.4pt}\\[-0.9ex]
\rule{5em}{0.4pt}
\end{center}

\clearpage
% --- p. 80 ---
% INDHOLD (p. 80), no folio printed. Heading in spaced capitals, centred,
% over a short rule; „Side“ heads the page-number column; dot leaders. The
% chapter titles here differ from the heads in the text (II--IV drop
% „-begrebet“; V reads „Værdikategorierne“ for „Værdi“): transcribed as
% printed on this page.
\section*{INDHOLD}
\addcontentsline{toc}{section}{Indhold}

\noindent\opage{80}%
\begin{tabular}{@{}r@{\ \ }p{0.72\textwidth}@{\ }r@{}}
 & & Side\\
I. & Kategorilærens Stilling i Filosofien\dotfill & 3---23\\
II. & Totalitet og de fundamentale Kategorier\dotfill & 23---31\\
III. & Totalitet og de formale Kategorier\dotfill & 31---42\\
IV. & Totalitet og de reale Kategorier\dotfill & 42---66\\
V. & Totalitet og Værdikategorierne\dotfill & 66---73\\
VI. & Totalitet som Grænsebegreb\dotfill & 74---78
\end{tabular}

\begin{center}
\rule{5em}{0.4pt}
\end{center}

% TRYKFEJL (p. 80): heading in spaced capitals, centred, over a short rule;
% eight lines, „p.“ on the first and a ditto dash „-“ on the rest; a
% short rule closes the page. Transcribed line by line as printed.
\section*{TRYKFEJL}
\addcontentsline{toc}{section}{Trykfejl}

\begin{center}
\begin{tabular}{@{}l@{\ }l@{}}
p. & 31, lin.\ 6 f.\ n.\ deres --- læs: den.\\
- & 33, lin.\ 21 f.\ o.\ Individualforestillinger --- læs: Individualbegreber.\\
- & 36, lin.\ 1 f.\ o.\ fra --- læs: for.\\
- & 36, lin.\ 20 f.\ o.\ Erfaringer --- læs: Erfaringen.\\
- & 36, anden Note, lin.\ 3 f.\ o.\ samtlige --- læs: sanselige.\\
- & 44, lin.\ 19 f.\ o.\ Farveaktens --- læs: Farvesansens.\\
- & 45, første Note.\ 340 --- læs: 388.\\
- & 64, lin.\ 3 f.\ n.\ første --- læs: faste.
\end{tabular}

\bigskip
\rule{5em}{0.4pt}
\end{center}

\end{document}
